% 用 pdflatex 编译（ctexart 会自动调用 CJK 支持）
\documentclass[UTF8,a4paper,11pt]{ctexart}
\usepackage{anyfontsize}
\usepackage{geometry}
\usepackage{fancyhdr}
\usepackage{amsmath,amssymb,amsthm,mathtools}
\usepackage[T2A]{fontenc}
\usepackage[russian,english]{babel}
\usepackage{mathrsfs}
\usepackage{array}
\usepackage{booktabs}
\usepackage{tabularx}
\usepackage{url}
\usepackage{tikz}
\usetikzlibrary{arrows.meta,positioning,calc,shapes.geometric}
\geometry{a4paper,scale=0.8}
\setlength{\headheight}{18pt}
\allowdisplaybreaks[2]
\emergencystretch=3em
\newcommand{\ii}{\mathrm{i}}
\newcommand{\QQ}{\mathbb Q}
\newcommand{\ZZ}{\mathbb Z}
\newcommand{\NN}{\mathbb N}
\title{\textbf{\huge{分圆方程与根式去嵌套}}\\[6pt]
\large{分圆方程的一般化根式解法 \;|\; 根式可开方性的判定与算法}}
\author{\large{\textit{Rinchi}}}
\date{\kaishu 2026 年 10 月}
\setcounter{tocdepth}{3}
\pagestyle{fancy}
\fancyhead[L]{分圆方程与根式去嵌套}
\begin{document}
\theoremstyle{definition}
\newtheorem*{solve}{解}
\newtheorem*{syoumei}{证}
\newtheoremstyle{myplain}
{\topsep}{\topsep}{\itshape}{0pt}{\bfseries}{}{1em}
{\thmname{#1}\thmnumber{ #2.}\thmnote{ #3}}
\newtheoremstyle{mydefinition}
{\topsep}{\topsep}{\normalfont}{0pt}{\bfseries}{}{1em}
{\thmname{#1}\thmnumber{ #2.}\thmnote{ #3}}
\theoremstyle{myplain}
\newtheorem{theorem}{定理}[section]
\newtheorem{proposition}[theorem]{命题}
\newtheorem{lemma}[theorem]{引理}
\newtheorem{corollary}[theorem]{推论}
\theoremstyle{mydefinition}
\newtheorem{definition}[theorem]{定义}
\newtheorem{example}[theorem]{例}
\newtheorem{remark}[theorem]{注}
\newtheorem{notation}[theorem]{记号}
\newtheorem{convention}[theorem]{约定}
\rmfamily
\maketitle
\begin{abstract}
本文由两条线索合成；第一条求分圆方程 $x^{n}=1$ 的显式根式解，
给出全阶法、半阶法与周期下降法三条路线：
前两条各用一次高次开方把全部单位根一次写出，形式上最短，
却没有把那个高次根号真正开出来；第三条沿子域逐层下降，
每层只开一个小的素数次方，手算步骤最长，但每一步的被开方数都落在已求出的实数域里，
因而可以把根号开到底；
第二条判定嵌套根式 $\sqrt[n]{\alpha}$ 能否去嵌套并求出结果，
给出范数必要条件、两层封闭公式、Landau 单项式引理与完备判定，
配以六步手算流程。
\end{abstract}
\tableofcontents
\pagebreak

\section*{前言}
分圆方程，即为形如$x^{n}=1$的方程，它的$n$个根在复平面上把单位圆$n$等分，因而得名；
求这类方程的根式解，是古典代数里最有分量的问题之一：
1796 年，十九岁的 C.F.\ Gauß 发现正十七边形可以只用圆规和直尺作出来，
五年之后他把完整的答案写进《算术研究》，
其中既给出了分圆方程根式可解性的证明，也给出了具体的算法；
本文不去重复那份证明，而是把其中几条算法的骨架完整地摆出来，
写清各条路线分别在算什么、能算到什么程度；\par

全文沿两条线索展开：其一求$x^{n}=1$的显式根式解，
给出全阶法、半阶法与周期下降法三条路线；
其二判定一个嵌套根式$\sqrt[n]{\alpha}$能否去嵌套，并给出求值的方法；
这两条线索并非彼此独立——前两条路线最后都停在一个高次根号上，
而那个根号的被开方数本身就是若干层根式的嵌套，
能不能把它真正开出来，正是第二条线索要回答的问题；\par

三条路线可以一句话概括：全阶法把全部$P-1$个非平凡根交给一个$Q$次根号，
半阶法借助共轭配对把那个根号降到$H$次，
周期下降法则把一个大根号换成一串小根号；
前两条路线的公式短、思路净，却没有把那个高次根号真正开出来；
第三条路线的流程长，但每一步都落在已经求出的实数上，
直到得出每个单位根的实部与虚部；\par

本文的结构安排如下：第一章交代问题的由来、要用到的四件工具，
以及三条路线共用的套路；第二章给出全阶法与半阶法两条一次开方的路线，
并用五次、十七次、七次三个算例说明各自的计算量；
第三章给出周期下降法，并沿十三次与十七次两个算例逐层走到底；
第四章判定嵌套根式能否去嵌套，给出两种「可去」的含义、完备判据与求值方法；
附录收拢全书记号、术语与参考文献；\par


\section{从正十七边形到三条路线}
\subsection{正十七边形与分圆方程的由来}
分圆方程就是形如
\[
 x^{n}=1\qquad(n\in\NN_{+})
\]
的方程；在复平面上，它的$n$个根把单位圆$n$等分，因而得名「分圆」；
我们在初中就认识$x^{2}=1$的两个根$\pm1$，
学过复数之后又认识了$x^{3}=1$与$x^{4}=1$的根，
熟悉黄金分割率与$\sin36^{\circ}$的人大致还知道$x^{5}=1$的根；
那么一般的$n$次单位根，也就是$x^{n}=1$的根式解，又该怎么得到呢；

作为合数的情形并不带来新的困难；按算术基本定理，
正整数$n$总可以写成若干个素数幂的乘积：
\[
 n=\prod_{k=1}^{s}P_k^{\alpha_k}
\]
只要方程左端能提出一个$P_s^{\alpha_s-1}\prod_{k<s}P_k^{\alpha_k}$次幂，就得到：
\[
 \Bigl(x^{P_s^{\alpha_s-1}\prod_{k<s}P_k^{\alpha_k}}\Bigr)^{P_s}=1
\]
于是先解一个$P_s$次分圆方程，再逐层回代，$n$次的情形就化归为素数次的情形；
因此本文只讨论素数次的单位根；\par

1796 年 3 月 30 日，十九岁的 C.F.\ Gauß 发现：正十七边形可以只用圆规和直尺作出来；
用今天的语言说，这件事等价于$\omega_{17}$的实部与虚部都能写成嵌套的二次根式；
五年后他把完整的答案写进《算术研究》，
其中既给出了分圆方程根式可解性的证明，也给出了具体的算法；\par


最后说明本文的记号分工：序列指标与函数变量用小写拉丁字母，
过程中的具体量用大写拉丁字母，过程中的函数（各种表）用小写希腊字母，
预制映射用大写希腊字母，单位根只由$\omega$承担，虚数单位一律写作$\ii$；
三条路线用的预解式统一记作$Z$，只用上标的阶区分；
后文的字母分工表列出本文记号与三套原始记号之间的对照；
\subsection{需要用到的四件工具}

\subsubsection{模$P$的乘法群与数论阶}
对正整数$m$与整数$a$，记$(a,m)$为$a$与$m$的最大公因数；
同余关系把一个整数集分成若干个剩余类，这一节只看与$m$互素的那些类；

对正整数$m$，把模$m$的全部互素剩余类收集起来：
\[
 \mathscr{G}_m:=\{[a]_m:\ 1\le a\le m,\ (a,m)=1\}
\]
并在其中定义乘法为「先相乘再取模」；可以逐条验证：
乘法满足结合律；$[1]_m$是单位元；
每个类$[a]_m$都有逆元，因为$(a,m)=1$时$ax\equiv1\pmod m$总有解；
于是$(\mathscr{G}_m,\cdot)$构成一个群，称为模$m$的乘法群；
特别地，$m=P$为素数时，$1$到$P-1$的每个整数都与$P$互素，
所以：
\[
 \mathscr{G}_P=\{1,2,\dots,P-1\},\qquad |\mathscr{G}_P|=Q=P-1
\]
这里$[a]_P$就简写成$a$本身，下文对模$P$的运算一律如此；\par

群中元素的个数称为群的阶；对$a\in\mathscr{G}_m$，
使得$a^{k}\equiv1\pmod m$成立的最小正整数$k$称为$a$对模$m$的数论阶，
记作$\operatorname{ord}_m(a)$；有限群里每个元素的阶都有限，
并且$a^k\equiv1\pmod m$成立的充要条件是$\operatorname{ord}_m(a)$整除$k$；
举例来说，$\operatorname{ord}_7(2)=3$，因为：
\[
 2^1\equiv2,\qquad 2^2\equiv4,\qquad 2^3\equiv8\equiv1\pmod7
\]
而$2$的更低次幂都不等于$1$；\par

\subsubsection{原根与编号表}
若$a$对模$m$的数论阶恰好等于群的阶$\varphi(m)$，
就称$a$为模$m$的一个原根；原根把整个群「排成一行」：
若$G$是模$P$的原根，则$G^0,G^1,\dots,G^{Q-1}$模$P$两两不同余，
因而恰好遍历$1,2,\dots,P-1$的全部非零剩余类；

模$P$的原根总是存在的（这一条也是 Gauß 证明的），找起来也并不费力：
只需要从$2,3,4,\dots$逐个试，看哪一个的阶等于$Q$即可；
例如模$5$的原根可以取$G=2$：
\[
 2^0,2^1,2^2,2^3\equiv1,2,4,3\pmod5
\]
四个数恰好把非零剩余类跑了一遍；模$13$的原根可以取$G=6$，
模$17$的原根可以取$G=3$，这两张表在后文的算例里会从头用到尾；\par

给定原根$G$之后，可以写两张互查的表；一张是原根表：
\[
 G^a=[G^a]_P,\qquad a=0,1,\dots,Q-1
\]
另一张是它的反查，称为编号表；对$1\le x\le P-1$，
方程$G^{\delta_G(x)}\equiv x\pmod P$在$0\le\delta_G(x)<Q$内有唯一解，
由此定义的$\delta_G(x)$就是$x$在原根表里的行号：
\[
 \delta_G:\{1,2,\dots,P-1\}\longrightarrow\{0,1,\dots,Q-1\}
\]
我们把$G^a$表的两行上下对调就得到它；\par

编号表是整篇文章的计算入口；它把「模$P$的乘法」
翻译成「模$Q$的加法」：$xy$的编号是$\delta_G(x)+\delta_G(y)$模$Q$；
后文所有公式里的指数都来自这张表，而不是来自任何角的计算；\par

\subsubsection{本原单位根}
欧拉公式给出$e^{2\pi t\ii}=1$对一切整数$t$成立，因此分圆方程可以写成：
\[
 x^{n}=e^{2\pi t\ii}
\]
它的通解是$x=e^{2\pi t\ii/n}$；为讨论方便，记：
\[
 \omega_n^t:=e^{2\pi t\ii/n}
\]
于是$\omega_n^0,\omega_n^1,\dots,\omega_n^{n-1}$恰好是$x^n=1$的全部根，
并且它们在乘法下构成一个$n$阶循环群；
若$(t,n)=1$，则称$\omega_n^t$为$n$次本原单位根；
通俗地说，本原单位根就是在$n$次单位根里「新出现」的那些，
它们不是任何一个更低次分圆方程的根；
最小的几个例子是$\omega_2^1=-1$、$\omega_3^1=\dfrac{-1+\sqrt3\,\ii}{2}$、
$\omega_4^1=\ii$；

$n\le5$时的单位根都能直接写成根式，后文的低次单位根表把常用的几组列出，
算到中间需要时就近查用；次数更高时，这些辅助单位根本身也要按同一套算法递归构造，
所以它们不是「先假设已经知道」的量；\par

\subsubsection{Lagrange 预解式}
前三件工具都是铺垫，真正把单位根与根式联系起来的是这一件；

对素数$P$，设$G$是原根，$\omega_Q$是任一个本原$Q$次单位根；
我们把$P$个$P$次单位根按原根表的次序排成一列：
\[
 \omega_P^{G^0},\ \omega_P^{G^1},\ \dots,\ \omega_P^{G^{Q-1}}
\]
再用$\omega_Q$的各次幂给它们加权重，做一次长度为$Q$的离散 Fourier 变换：
\[
 Z^{(Q)}_j:=\sum_{k=0}^{Q-1}\omega_Q^{jk}\,\omega_P^{G^k},\qquad j=0,1,\dots,Q-1
\]
这一族量就称为 Lagrange 预解式；它的好处有两层；\par

第一层是「反演得回去」；由于$\omega_Q$是$Q$次本原单位根，
几何级数恒等式$\sum_{j=0}^{Q-1}\omega_Q^{jk}=0$（$Q\nmid k$时）给出：
\[
 \omega_P^{G^k}=\frac1Q\sum_{j=0}^{Q-1}\omega_Q^{-jk}\,Z^{(Q)}_j
\]
只要能算出全部$Z^{(Q)}_j$，把它按上式线性组合就得到全部单位根；
从预解式到单位根的这一转折，是三条算法共用的骨架；\par

第二层是「$Z^{(Q)}_j$之间有乘法关系」；若把$Z^{(Q)}_r$与$Z^{(Q)}_1$相乘，
乘积中$k$与$k+1$的配对会自然出现，整理之后总能写成：
\[
 Z^{(Q)}_r Z^{(Q)}_1=J_r\,Z^{(Q)}_{r+1}
\]
其中$J_r$是一个只依赖模$P$计数、不含任何未知单位根的数；
这条关系的含义是：只要知道$Z^{(Q)}_1$，其余$Z^{(Q)}_j$都能由它乘上已知系数得到；
于是反演式里那些未知量，实际上只是同一个数的不同幂次；
把$j=1$那条链走到底，我们还会得到$Z^{(Q)}_1$的$Q$次幂等于一个已知的多项式值，
于是「求全部单位根」这件事被压缩成「开一个根号」；
\subsection{三条路线共用的套路：把根写成预解式的线性组合}
在把三条算法并排摆开之前，先把它们共用的三步固定下来；
下面这套做法在后面每一节都会重新出现一次，只是「对谁做 Fourier」不同；\par

\textbf{第一步：选一组中间量，并对它做 Fourier 变换；}
中间量是一族关于$P$次单位根的和，Fourier 变换的阶数决定了后面要开多少次方；
全阶法对全部$Q$个非零剩余类做，阶数$Q$；
半阶法对$H$个实迹做，阶数$H$；
周期下降法把$H$拆成一条链$H=L_1L_2\cdots L_g$，每一层只对$L_r$个块做，阶数$L_r$；\par

\textbf{第二步：把相邻两项的乘积写成已知系数乘下一项；}
这一步是每一种方法里唯一需要算的地方，而它总是化归为有限计数：
模$P$的两个剩余类相乘、相加落在哪一类里，数一遍就行；
写出来的系数都是整数、有理数或者低次单位根的有理式，不含未知量；\par

\textbf{第三步：把中间量反演回去；}
由第一步的 Fourier 变换是正交的，中间量总能用几何级数反演成原来的和，
$\sum_j\omega_Q^{jk}=0$（$jk$不是$Q$的倍数时）就是全部依据；\par

以最简单的情形为例，$P=3$，取原根$G=2$、辅助单位根$\omega_Q=\omega_2^1=-1$；
预解式是$Z^{(2)}_1=\omega_3^1-\omega_3^2$，
而$\omega_3^1+\omega_3^2=-1$，所以：
\[
 \omega_3^1=\frac{-1+Z^{(2)}_1}{2},\qquad \omega_3^2=\frac{-1-Z^{(2)}_1}{2}
\]
再算$(Z^{(2)}_1)^2=(\omega_3^1-\omega_3^2)^2=\omega_3^2-2+\omega_3^1=-3$，
而$-3$的平方根写成实部加虚部就是$\sqrt3\,\ii$，于是：
\[
 \omega_3^1=\frac{-1+\sqrt3\,\ii}{2},\qquad \omega_3^2=\frac{-1-\sqrt3\,\ii}{2}
\]
三步走完，全部三次单位根就齐了；后文三种方法给出的都是同一件事：
一组中间量、一组系数、一个反演式，区别只在中间量怎么选、系数怎么算，
以及最后那个根号能不能开得动。
\pagebreak

\section{两条一次开方的路线：全阶法与半阶法}
\subsection{全阶法：一次开$P-1$次方写出全部根}
\subsubsection{算法}
全阶法的中间量就是上一节的$Z^{(Q)}_j$，一次把全部$Q$个非零剩余类都用上；

我们取本原$Q$次单位根$\omega_Q$（它本身的根式表达由更低次的分圆方程给出，
次数$Q<P$，所以这是一个规模更小的同类问题），定义：
\[
 Z^{(Q)}_j=\sum_{k=0}^{Q-1}\omega_Q^{jk}\,\omega_P^{G^k},\qquad j=0,1,\dots,Q-1
\]
两个端点值可以立刻算出来：$Z^{(Q)}_0=\sum_k\omega_P^{G^k}=-1$（全部非平凡根之和），
以及$Z^{(Q)}_Q=Z^{(Q)}_0=-1$；\par

我们把$Z^{(Q)}_r$与$Z^{(Q)}_1$相乘并按$k$与$k+1$配对，得到乘法关系：
\[
 Z^{(Q)}_r Z^{(Q)}_1=J_r\,Z^{(Q)}_{r+1},\qquad 1\le r\le Q-2
\]
其中$J_r$是只依赖原根表的一个整数组合：
\[
 J_r=\sum_{\substack{0\le a,b\le Q-1\\ G^a+G^b\equiv1\pmod P}}\omega_Q^{r a+b}
\]
只要手上已经有$G^a$表与$\omega_Q$的取值，这个系数就能逐项数出来，
不需要任何未知的单位根；另有一对互补的关系：
\[
 Z^{(Q)}_j Z^{(Q)}_{Q-j}=(-1)^jP,\qquad
 \overline{Z^{(Q)}_j}=(-1)^jZ^{(Q)}_{Q-j}
\]
它说明预解式两两互为共轭（差一个符号），模长平方全都已知；\par

我们把第一条关系从$r=1$用到$r=j-1$，连乘得到$Z^{(Q)}_j$用$Z^{(Q)}_1$表示的式子：
\[
 Z^{(Q)}_j=\frac{\bigl(Z^{(Q)}_1\bigr)^{j}}{B_j},\qquad
 B_j=\prod_{r=1}^{j-1}J_r,\qquad 1\le j\le Q-1
\]
再用第二条关系取$j=1$：$Z^{(Q)}_1Z^{(Q)}_{Q-1}=-P$，
把它与$j=Q-1$的上式联立，就得到关键的一步：
\[
 \bigl(Z^{(Q)}_1\bigr)^{Q}=E^{(Q)}_{P,G}(\omega_Q),\qquad
 E^{(Q)}_{P,G}(w)=\Bigl(-P\prod_{r=1}^{Q-2}J_r(w)\Bigr)\bmod\mu_Q(w)
\]
右端的含义是：把每个$J_r$里的$\omega_Q$换成一个未定元$w$，
乘起来得到关于$w$的整系数多项式，再对分圆多项式$\mu_Q$取余；
取余不改变在$w=\omega_Q$处的值，却把多项式的次数降到可控范围；
于是$E^{(Q)}_{P,G}$是一个完全可以写下来的整系数多项式，
$Z^{(Q)}_1$就是它的一个$Q$次方根：
\[
 U^{(Q)}:=\sqrt[Q]{E^{(Q)}_{P,G}(\omega_Q)}
\]
把回代与反演合在一处，全部非平凡单位根由一个公式给出：
\[
 \omega_P^{G^k}=\frac1Q\Bigl(-1+\sum_{j=1}^{Q-1}\frac{\omega_Q^{-jk}}{B_j}\,\bigl(U^{(Q)}\bigr)^{j}\Bigr),
 \qquad k=0,1,\dots,Q-1
\]
公式里只有一个根号，就是那个$Q$次根；其余全是四则运算与已知的$\omega_Q$的幂。\par

\subsubsection{五次的计算：一次开四次方}
我们取$P=5$，原根$G=3$，辅助单位根取本原四次单位根$\omega_Q$，
即$\omega_Q^2=-1$、$\omega_Q^4=1$；四个预解式是：
\[
 \begin{aligned}
 Z^{(4)}_1&=\omega_5+\omega_Q\,\omega_5^{3}+\omega_Q^2\omega_5^{4}+\omega_Q^3\omega_5^{2}\\[2pt]
 Z^{(4)}_2&=\omega_5-\omega_5^{3}+\omega_5^{4}-\omega_5^{2}\\[2pt]
 Z^{(4)}_3&=\omega_5-\omega_Q\,\omega_5^{3}+\omega_5^{4}-\omega_Q^3\omega_5^{2}\\[2pt]
 Z^{(4)}_4&=-1
 \end{aligned}
\]
把它们按$\omega_5$的幂次整理，逐项乘开可以数出三条关系：
\[
 \bigl(Z^{(4)}_1\bigr)^2=(-1+2\omega_Q)Z^{(4)}_2,\qquad
 Z^{(4)}_1Z^{(4)}_3=-5,\qquad
 \bigl(Z^{(4)}_1\bigr)^4=-15-20\,\omega_Q
\]
第一条给出$Z^{(4)}_2=\dfrac{\bigl(Z^{(4)}_1\bigr)^2}{-1+2\omega_Q}$，
第二条给出$Z^{(4)}_3=-\dfrac{5}{Z^{(4)}_1}$，
第三条则说明$Z^{(4)}_1$是一个四次方根；把$\omega_Q$按$\omega_Q^2=-1$化简，
$-15-20\,\omega_Q$的模长是$\sqrt{15^2+20^2}=25$，
于是它可以写成$25$乘上一个$\omega_Q$的幂与一个正实数的四次根：
\[
 Z^{(4)}_1=\omega_Q^{m}\sqrt{5}\Bigl(\sqrt{\frac{5-\sqrt5}{10}}+\sqrt{\frac{5+\sqrt5}{10}}\,\ii\Bigr)
 =\frac{\omega_Q^{m}}{2}\Bigl(\sqrt{10-2\sqrt5}+\sqrt{10+2\sqrt5}\,\ii\Bigr),
 \qquad m=0,1,2,3
\]
四个分支$m$逐个代入反演式，四次单位根就全部到手；
我们把结果按$\omega_5$的幂次顺序排好，并把实部与虚部分开写：
\[
 \begin{aligned}
 \omega_5^{1}&=\dfrac{-1+\sqrt5}{4}+\dfrac{\sqrt{10+2\sqrt5}}{4}\,\ii\\[6pt]
 \omega_5^{2}&=\dfrac{-1-\sqrt5}{4}+\dfrac{\sqrt{10-2\sqrt5}}{4}\,\ii\\[6pt]
 \omega_5^{3}&=\dfrac{-1-\sqrt5}{4}-\dfrac{\sqrt{10-2\sqrt5}}{4}\,\ii\\[6pt]
 \omega_5^{4}&=\dfrac{-1+\sqrt5}{4}-\dfrac{\sqrt{10+2\sqrt5}}{4}\,\ii
 \end{aligned}
\]
这四个式子与大家熟悉的黄金分割表达式一致：
例如$\omega_5^1+\omega_5^4=2\cos\frac{2\pi}{5}=\frac{-1+\sqrt5}{2}$；

从这道题里可以看清全阶法的形状：计算量集中在一处，
也就是把$\bigl(Z^{(4)}_1\bigr)^4$那个四次根开出来；
$P$大一些时，这一步的内容不变，只是次数从$4$变成$Q$。\par

\subsubsection{计算量}
全阶法需要数出$J_1,J_2,\dots,J_{Q-2}$共$Q-2$个系数，
每个系数是一个长度约$Q$的求和，
即在$P-1$个剩余类里逐对找出满足$a+b\equiv1$的配对；
按朴素方式数一遍，代价是$(Q-2)(Q-1)$量级的整数运算；
这对纸笔而言不算轻松，但每一步都是纯计数，没有需要「猜」的地方。
\subsection{半阶法：借助共轭配对把开方次数减半}
\subsubsection{算法}
上一节最后那对互补关系$Z^{(Q)}_jZ^{(Q)}_{Q-j}=(-1)^jP$透露了一件事：
预解式天然成对出现，而互为共轭的一对给出的是同一个实数信息；
半阶法正是从这里出发，把「有效分量」的个数减少一半；

因为$G^{H}\equiv-1\pmod P$（这是原根的基本性质），
原根表的后半段是前半段的相反数，于是$\omega_P^{G^{a+H}}=\omega_P^{-G^{a}}$，
两者互为共轭，它们的和是实数：
\[
 S_{H,a}:=\omega_P^{G^{a}}+\omega_P^{-G^{a}}=2\cos\frac{2\pi G^{a}}{P},
 \qquad 0\le a<H
\]
这$H$个实数称为实迹，它们是全部$P$个单位根按共轭配对之后剩下的一半信息，
外加一个明显的已知量$\sum_aS_{H,a}=-1$；\par

我们取本原$H$次单位根$\omega_H$（它与$\omega_Q$的关系是$\omega_H=\omega_Q^2$），
对$H$个实迹做一次长度为$H$的 Fourier 变换，得到半阶预解式：
\[
 Z^{(H)}_k=\sum_{a=0}^{H-1}\omega_H^{ka}S_{H,a}
 =\sum_{a=0}^{H-1}\omega_H^{ka}\bigl(\omega_P^{G^{a}}+\omega_P^{-G^{a}}\bigr),
 \qquad k=0,1,\dots,H-1
\]
半阶预解式不是另起炉灶的新对象：
把$\omega_P^{G^{a+H}}=\omega_P^{-G^a}$代回全阶预解式的定义，
偶数位的分量恰好就是它，也就是：
\[
 Z^{(H)}_k=Z^{(Q)}_{2k},\qquad 0\le k<H
\]
全阶法用$Q$个分量描述全部$P$个单位根，半阶法只取其中下标为偶数的$H$个，
因为奇数位那一半与它们互为共轭，不带来新信息；
分量个数由$Q-1$降到$H-1$，这就是次数减半的来源；\par

乘法关系的形式与全阶法完全一样：
\[
 Z^{(H)}_r Z^{(H)}_1=K_r\,Z^{(H)}_{r+1},\qquad
 Z^{(H)}_k Z^{(H)}_{H-k}=P,\qquad
 \overline{Z^{(H)}_k}=Z^{(H)}_{H-k}
\]
\[
 K_r=\sum_{\substack{0\le a,b\le Q-1\\G^a+G^b\equiv1\pmod P}}\omega_H^{r a+b}
\]
终点关系从$(-1)^kP$变成$P$，是因为半阶的特征标在$-1$处取值为$1$，
配对时不再出现符号；记累乘$C_j=\prod_{r=1}^{j-1}K_r$，则：
\[
 Z^{(H)}_j=\frac{\bigl(Z^{(H)}_1\bigr)^{j}}{C_j},\qquad
 \bigl(Z^{(H)}_1\bigr)^{H}=E^{(H)}_{P,G}(\omega_H),\qquad
 E^{(H)}_{P,G}(w)=\Bigl(P\prod_{r=1}^{H-2}K_r(w)\Bigr)\bmod\mu_H(w)
\]
同样地，右端是一个可以写出来的整系数多项式在$w=\omega_H$处的值；
把$U^{(H)}$取作它的一个$H$次方根，规范化之后的分量是：
\[
 \bigl(U^{(H)}\bigr)^{H}=E^{(H)}_{P,G}(\omega_H),\qquad
 \frac{Z^{(H)}_j}{C_j}=\frac{\bigl(U^{(H)}\bigr)^{j}}{C_j},\qquad 1\le j\le H-1
\]
反演给出全部实迹，再由实迹取出每个单位根：
\[
 -1+\sum_{j=1}^{H-1}\omega_H^{aj}\,\frac{\bigl(U^{(H)}\bigr)^{j}}{C_j}
 =H\Bigl(\omega_P^{G^{-a}}+\omega_P^{-G^{-a}}\Bigr),
 \qquad 0\le a<H
\]
最后一步是解二次方程：已知$z+z^{-1}=S$，则$z$是$y^2-Sy+1=0$的根，
用求根公式取出实部与虚部即可；整条路线上只有一个根号，次数是$H$；
半阶法与全阶法之间还有一层换算关系，见后文「与半阶法的换算」一节。\par

\subsubsection{十七次的计算：一次开八次方}
$P=17$时$Q=16$、$H=8$，全阶法要开$16$次方，半阶法则只需要一个$8$次根；
上一节的公式给出$\bigl(U^{(H)}\bigr)^{8}=E^{(H)}_{17,3}(\omega_H)$，
其中$\omega_H$是八次本原单位根；
把$\omega_H$的取值代进去，右端化为具体的数，
$U^{(H)}$就是一个八次方根，进而反演出八个实迹；
再对每个实迹解一次二次方程，得到全部十六个非平凡单位根；
这一路走下来，八次根号是唯一的「高次」障碍。

\subsubsection{七次的计算：辅助根恰好是三次单位根}
十七次那个例子里，辅助单位根$\omega_H$是八次单位根，它自己也要靠递归解决；
换一个更小的例子：$P=7$；这时$Q=6$、$H=3$，
辅助单位根是三次单位根，而三次单位根可以直接查表；

我们取模$7$的原根$G=3$；原根表是：
\[
 G^0,G^1,\dots,G^5=1,3,2,6,4,5
\]
六个非零剩余类恰好各出现一次；
本原三次单位根按低次单位根表展开，实部与虚部分开写：
\[
 \omega_H=\omega_3^1=\frac{-1+\sqrt3\,\ii}{2},\qquad
 \omega_H^2=\frac{-1-\sqrt3\,\ii}{2}
\]

第一步，写三个实迹；第$0$层的原根表把六个非零剩余类排成一行，
每个都要与它的相反数配成一对；三对的和就是三个实迹：
\[
 S_{3,0}=\omega_7^{1}+\omega_7^{6}=2\cos\frac{2\pi}{7},\qquad
 S_{3,1}=\omega_7^{3}+\omega_7^{4}=2\cos\frac{6\pi}{7},\qquad
 S_{3,2}=\omega_7^{2}+\omega_7^{5}=2\cos\frac{4\pi}{7}
\]
三者之和为$-1$，这是第一道验算；它们满足三次方程：
\[
 y^3+y^2-2y-1=0
\]
这个三次方程的判别式是：
\[
 1^2(-2)^2-4(-2)^3-4\cdot1^3(-1)-27(-1)^2+18\cdot1\cdot(-2)(-1)
 =4+32-4-27+36=49=7^2
\]
又是一个完全平方，所以三个根都是实数——
七次与十三次一样，都落在实三次域里；\par

第二步，数系数；$H=3$时只有$K_1$一个非平凡系数，
按定义在$1\le a,b\le6$里数满足$3^a+3^b\equiv1\pmod7$的配对：
\begin{center}
\begin{tabular}{c|cccccc}
\toprule
$s$ & 1&2&3&4&5&6\\
\midrule
$3^a$ & 3&2&6&4&5&1\\
\midrule
使$3^b\equiv1-3^a$的$b$ & 4&2&6&3&5&---\\
$\omega_H^{a+b}$ & $\omega_H^5=\omega_H^2$ & $\omega_H^4=\omega_H$ & $\omega_H^9=\omega_H^0$ & $\omega_H^7=\omega_H$ & $\omega_H^{10}=\omega_H$ & ---\\
\bottomrule
\end{tabular}
\end{center}
最后一行是把指数按$\omega_H^3=1$约化之后的结果；
六个根里$x=1$自身要排除（$3^a\equiv1$须$1-1=0$，不是非零剩余类），
所以一共五项，相加得：
\[
 K_1=\omega_H^2+\omega_H+1+\omega_H+\omega_H=1+3\omega_H=\frac12-\frac{3\sqrt3}{2}\ii
\]
与$K_1\overline{K_1}=7$相符：设$c+d\omega_H$的共轭是$c+d\omega_H^2$，
两者相乘用$\omega_H+\omega_H^2=-1$、$\omega_H\omega_H^2=1$化简，
得到$(c+d\omega_H)(c+d\omega_H^2)=c^2-cd+d^2$；
把$c=1$、$d=3$代进去，正好是$1-3+9=7$；\par

另一半阶关系给出$v_2=v_1^{2}/K_1$，所以非平凡分量只有一个独立值；\par

第三步，写半阶预解式；按定义$v_j=\sum_a\omega_H^{ja}S_{3,a}$，
只要算出$v_1$，反演式里就只剩它和它的共轭：
\[
 v_1=S_{3,0}+\omega_H S_{3,1}+\omega_H^2 S_{3,2}
\]
它的模长平方可以直接验证：
\[
 v_1\overline{v_1}=P=7
\]
而它的三次方落在$\QQ(\omega_H)$里，化简之后是一个很短的式子：
\[
 v_1^{3}=-7-21\omega_H
\]
把它写成「实部加虚部乘$\ii$」的样子，只要把
$\omega_H=\dfrac{-1+\sqrt3\,\ii}{2}$代进去：
\[
 -7-21\omega_H=-7+21\cdot\frac12-21\cdot\frac{\sqrt3}{2}\ii
 =\frac72-\frac{21\sqrt3}{2}\ii
\]
于是公共根式是：
\[
 v_1^{\,3}=E^{(H)}_{7,3}(\omega_H)=-7-21\omega_H
 =\frac72-\frac{21\sqrt3}{2}\ii
\]
它只是一个「有理数加有理数乘$\sqrt3\,\ii$」的复数，
而不是十七次那种嵌套许多层的开方量；
这里不必再引入新字母：公共根式就是$v_1$本身，三次方直接对上；
式子里只剩一个立方根与$\sqrt3\,\ii$两样东西；\par

第四步，反演；把$v_1$与它的共轭$\overline{v_1}$代入：
\[
 S_{3,a}=\frac{-1+\omega_H^{-a}v_1+\omega_H^{a}\overline{v_1}}{3},\qquad a=0,1,2
\]
三个数相加等于$-1$，与第一道验算一致；
我们把它们按数值排一下，最大的那个是：
\[
 S_{3,0}=\frac{-1+v_1+\overline{v_1}}{3}=2\cos\frac{2\pi}{7}
\]
这正是$\omega_7+\omega_7^{-1}$，也是三个实迹里最大的一个；
把$v_1$的实部记作一个符号，反演式就化成很省事的样子：
\[
 3S_{3,0}+1=v_1+\overline{v_1}
\]
于是另外两个实迹也能用它表示；具体地，用上一节的三个数
写出它们的公共三次方程再化简，得到两条恒等式：
\[
 S_{3,1}=-S_{3,0}^2-S_{3,0}+1,\qquad S_{3,2}=S_{3,0}^2-2
\]

第五步，收尾；收尾公式要查的两个下标是
$a_1=[\delta_3(1)]_3=[0]_3=0$与$b_1=[\delta_3(2)]_3=[2]_3=2$，
于是$W_0=S_{3,0}$、$W_2=S_{3,2}$，代入：
\[
 \omega_7^1=\frac{W_0}{2}+\frac12\sqrt{2-W_2}\ \ii
 =\frac{S_{3,0}}{2}+\frac12\sqrt{2-S_{3,2}}\ \ii
\]
被开方数在这里还能再化简一步：由$S_{3,2}=S_{3,0}^2-2$，
\[
 2-S_{3,2}=4-S_{3,0}^2
\]
而$S_{3,0}=2\cos\frac{2\pi}{7}$落在$(-2,2)$内，所以$4-S_{3,0}^2$是正数；\par

把七个$t$逐个查表代进去，就得到全部七个单位根；
记$S=S_{3,0}$，三个实迹分别是
\[
 W_0=S,\qquad W_1=-S^2-S+1,\qquad W_2=S^2-2
\]
再记住$2-W_2=4-S^2$、$2-W_1=S^2+S+1$、$2-W_0=2-S$，
七个根就是：
\[
 \begin{aligned}
 \omega_7^{1}&=\frac{S}{2}+\frac{\sqrt{4-S^2}}{2}\,\ii,
 &\omega_7^{2}&=\frac{S^2-2}{2}+\frac{\sqrt{S^2+S+1}}{2}\,\ii,\\[4pt]
 \omega_7^{3}&=\frac{-S^2-S+1}{2}+\frac{\sqrt{2-S}}{2}\,\ii,
 &\omega_7^{4}&=\frac{-S^2-S+1}{2}-\frac{\sqrt{2-S}}{2}\,\ii,\\[4pt]
 \omega_7^{5}&=\frac{S^2-2}{2}-\frac{\sqrt{S^2+S+1}}{2}\,\ii,
 &\omega_7^{6}&=\frac{S}{2}-\frac{\sqrt{4-S^2}}{2}\,\ii
 \end{aligned}
\]
六个非平凡根两两共轭，与$\omega_7^{7-t}=\overline{\omega_7^{t}}$相符；
再补上$\omega_7^0=1$，七个根齐了；\par

最后一件事是把$S$写成根式；$S$是$y^3+y^2-2y-1=0$的最大根，
解这个三次方程（同样要经过复数开立方）；前面已经算出
$v_1^{\,3}=-7-21\omega_H=\dfrac72-\dfrac{21\sqrt3}{2}\ii$，
它的模长是：
\[
 \Bigl|\frac72-\frac{21\sqrt3}{2}\ii\Bigr|
 =\sqrt{\frac{49}{4}+\frac{1323}{4}}=\sqrt{343}=7\sqrt7
\]
所以它的立方根模长恰好是$\sqrt7$；取模长$\sqrt7$、
辐角为$\dfrac13\arg\bigl(\frac72-\frac{21\sqrt3}{2}\ii\bigr)$的那一支，
记作：
\[
 U=\sqrt[3]{\frac72-\frac{21\sqrt3}{2}\ii},
 \qquad |U|=\sqrt7
\]
它就是三次方程求根公式里的那一项；把$3S=2\operatorname{Re}U-1$代回，
三个实迹就都成了$\operatorname{Re}U$的整系数多项式：
\[
 W_0=\frac{2\operatorname{Re}U-1}{3},\qquad
 W_1=\frac{-4(\operatorname{Re}U)^2-2\operatorname{Re}U+11}{9},\qquad
 W_2=\frac{4(\operatorname{Re}U)^2-4\operatorname{Re}U-17}{9}
\]
（这三条可以逐条代回去验算：把
$(\operatorname{Re}U)^2=\dfrac{2\sqrt7\operatorname{Re}U+7}{4}$
代进后两条，就化简成$W_0$的二次式；）\par

两个立方根的被开方数互为共轭：
$28-84\sqrt3\,\ii=8\cdot\bigl(\frac72-\frac{21\sqrt3}{2}\ii\bigr)$，
而$8=2^3$是立方数，开立方时正好把$2$提到根号外面：
\[
 \sqrt[3]{28-84\sqrt3\,\ii}
 =\sqrt[3]{8}\cdot\sqrt[3]{\frac72-\frac{21\sqrt3}{2}\ii}=2U
\]
两者的模长相同：$|28-84\sqrt3\,\ii|^2=28^2+84^2\cdot3=784+21168=21952=28^3$，
被开方数的模长是$28^{3/2}$，开立方之后是$28^{1/2}=2\sqrt7$；\par

于是七个根完整写开；下式只用根式与数字，不引入任何字母记号：
{\small
\[
 \begin{aligned}
 \omega_7^{1}&=\frac{-2+\sqrt[3]{28-84\sqrt3\,\ii}+\sqrt[3]{28+84\sqrt3\,\ii}}{12}\\[2pt]
 &\qquad+\frac{\sqrt{84+(3+3\sqrt3\,\ii)\sqrt[3]{28+84\sqrt3\,\ii}
   +(3-3\sqrt3\,\ii)\sqrt[3]{28-84\sqrt3\,\ii}}}{12}\ \ii,\\[8pt]
 \omega_7^{2}&=\frac{-4+(-1+\sqrt3\,\ii)\sqrt[3]{28-84\sqrt3\,\ii}
   +(-1-\sqrt3\,\ii)\sqrt[3]{28+84\sqrt3\,\ii}}{24}\\[2pt]
 &\qquad+\frac{\sqrt{84+(3+3\sqrt3\,\ii)\sqrt[3]{28-84\sqrt3\,\ii}
   +(3-3\sqrt3\,\ii)\sqrt[3]{28+84\sqrt3\,\ii}}}{12}\ \ii,\\[8pt]
 \omega_7^{3}&=\frac{-4+(-1+\sqrt3\,\ii)\sqrt[3]{28+84\sqrt3\,\ii}
   +(-1-\sqrt3\,\ii)\sqrt[3]{28-84\sqrt3\,\ii}}{24}\\[2pt]
 &\qquad+\frac{\sqrt{84-6\sqrt[3]{28-84\sqrt3\,\ii}-6\sqrt[3]{28+84\sqrt3\,\ii}}}{12}\ \ii,\\[8pt]
 \omega_7^{4}&=\frac{-4+(-1+\sqrt3\,\ii)\sqrt[3]{28+84\sqrt3\,\ii}
   +(-1-\sqrt3\,\ii)\sqrt[3]{28-84\sqrt3\,\ii}}{24}\\[2pt]
 &\qquad-\frac{\sqrt{84-6\sqrt[3]{28-84\sqrt3\,\ii}-6\sqrt[3]{28+84\sqrt3\,\ii}}}{12}\ \ii,\\[8pt]
 \omega_7^{5}&=\frac{-4+(-1+\sqrt3\,\ii)\sqrt[3]{28-84\sqrt3\,\ii}
   +(-1-\sqrt3\,\ii)\sqrt[3]{28+84\sqrt3\,\ii}}{24}\\[2pt]
 &\qquad-\frac{\sqrt{84+(3+3\sqrt3\,\ii)\sqrt[3]{28-84\sqrt3\,\ii}
   +(3-3\sqrt3\,\ii)\sqrt[3]{28+84\sqrt3\,\ii}}}{12}\ \ii,\\[8pt]
 \omega_7^{6}&=\frac{-2+\sqrt[3]{28-84\sqrt3\,\ii}+\sqrt[3]{28+84\sqrt3\,\ii}}{12}\\[2pt]
 &\qquad-\frac{\sqrt{84+(3+3\sqrt3\,\ii)\sqrt[3]{28+84\sqrt3\,\ii}
   +(3-3\sqrt3\,\ii)\sqrt[3]{28-84\sqrt3\,\ii}}}{12}\ \ii
 \end{aligned}
\]
}
三个根号内的量（即$2-W_0$、$2-W_2$、$2-W_1$）都是正的实数：
它们的形状是$2-2\cos\frac{2\pi k}{7}$，而每个实迹都严格落在$(-2,2)$内，
所以开方时一律取正实根；\par

这样写开后，式中除虚数单位$\ii$本身以外不含任何字母：
只有两个被开方数互为共轭的立方根、若干实数的平方根，
以及$(3\pm3\sqrt3\,\ii)$这样的实系数；
没有比$\sqrt[3]{\ }$更高的嵌套，也没有$\sqrt{-1}$这类写法；\par

这一例的$\omega_H$是三次单位根，直接查表即可，不必递归；
而$H=3$是半阶法能取到的最小非平凡值（$H=1$时$P=3$，没有非平凡分量），
所以七次是半阶法最简的完整算例。\par

\subsubsection{计算量}
半阶法要数$H-2$个系数，每个仍是长度约$Q$的求和，
总代价约$(H-2)(P-2)$次整数运算，比全阶法少了一半多；
$P=31$时全阶法是$29\cdot30=870$，半阶法降到$13\cdot29=377$。
\subsection{转折：两条路线都没有把根式开完}
前两节的两套算法都很短：一个高次根号，一个线性组合，全部单位根就出来了；
但它们交付的式子并不是可以照着写下来的数；
我们把两边的核心式子并排放在一起看：
\[
 U^{(Q)}=\sqrt[Q]{E^{(Q)}_{P,G}(\omega_Q)},\qquad
 U^{(H)}=\sqrt[H]{E^{(H)}_{P,G}(\omega_H)}
\]
这里有三个问题，一个比一个实际；

\textbf{第一，被开方数本身含有辅助单位根；}
$E^{(Q)}_{P,G}(\omega_Q)$里含$\omega_Q$，而$\omega_Q$是$Q$次单位根；
$E^{(H)}_{P,G}(\omega_H)$里含$\omega_H$，而$\omega_H$是$H$次单位根；
要写出$\omega_Q$或$\omega_H$的表达式，就要解次数更低的同类分圆方程；
这件事可以递归下去（次数严格变小，必然终止），
但刚才的公式本身没有做这件事：
它把「求$P$次单位根」的困难，换成了「求$Q$次单位根」的困难加一个$Q$次开方，
并没有把辅助单位根的实部与虚部真正写出来；\par

\textbf{第二，开方没有开到底；}
式子里出现的是$\sqrt[Q]{\ }$或$\sqrt[H]{\ }$本身，而不是它的展开；
所谓「开到底」，指的是两件具体的事：
把被开方数里平方因子、完全$d$次幂因子提出来；
把$Q$次（或$H$次）根写成若干个小次根式的嵌套；
全阶法与半阶法都没有做这一步；
$E^{(Q)}_{P,G}(\omega_Q)$只是合格地被写成了一个整系数多项式的取值，
它是否包含完全$Q$次幂因子、能否提因子，公式完全没有交代；\par

\textbf{第三，分支没有被规定；}
$U^{(Q)}$的$Q$个取值给出的是同一组单位根的重新排列，
而公式的右端$\omega_P^{G^k}$是按照固定编号写出来的；
到底是哪一个分支对应哪一列单位根，需要额外说明；
半阶法的情况相同：$U^{(H)}$的$H$个取值给出$H$个实迹的置换，
而实迹与编号之间的对应仍需另行确定；
仅凭约定「取主值」是不够的，因为复根式的主值在计算过程中并不保持相容；\par

一句话总结：这两套算法给出的是\textbf{形式解}，
它们证明了根式解存在，并且把「求根」化归为「开一个根」，
但没有完成最后那一步开方，也没有把结果整理成$a+b\ii$的样子；
要得到真正能写下来的数，我们就需要把那个大根号拆成开得动的小根号，
这就是下一节的周期下降法要做的事；\par

拆根号这件事本身并不只是把$\sqrt[ab]{Z}$改写成$\sqrt[a]{\sqrt[b]{Z}}$；
纯粹改写根号指标不会告诉我们任何新信息：
哪个中间量落在哪个数域里、系数是多少、各分支之间怎么保持相容，
都需要重新组织；周期下降法的下降，下降的是周期长度，
增加的是被确定下来的已知子域，这才使计算真的向前走。
\pagebreak

\section{周期下降法：逐层下降的根式解}
\subsection{方法总述}
周期下降法有一套自己的用字规则，与前面几节的$Z$、$U$、$E$互不冲突：
指标与变量用小写拉丁字母（块乘法的计数系数用西里尔大写$\textcyrillic{Б}$），
系数、常数、比值用大写拉丁字母，
函数与表用小写希腊字母，算子、泛函与矩阵用大写希腊字母，
集合与数域用大写花体拉丁字母；
背景记号$\omega_n^t$、$G$、$P$、$Q$、$H$、$\delta_G$与前面几节相同，不作改动；

手算之前只需准备三张表：原根表$G^0,G^1,\dots,G^{Q-1}$，
由它反查的编号表$\delta_G$，以及每算完一层就抄下来的块和表$S_{r,0},\dots,S_{r,D_r-1}$。\par

\subsubsection{周期：把根分块}
前两节都利用了同一件事：单位根可以配对，因此可以分组；
周期下降法把这个想法推到底——既然按共轭配对能把$P-1$个指数分成$H$组，
那就继续往下分，直到每组只剩两个指数；

具体做法是先取定一条分解链；把$H$按任意次序分解成素因子的乘积：
\[
 H=\prod_{a=1}^{g}L_a,\qquad D_0=1,\qquad D_r=\prod_{a=1}^{r}L_a\quad(1\le r\le g)
\]
这里$g$是链长，$L_1,L_2,\dots,L_g$是各个素因子；连乘的哑下标用$a$，
与层号$r$区别开；重复的素因子逐个保留，例如$H=8$写成$2\cdot2\cdot2$就是三层；
$D_r$的意义是「第$r$层一共有多少个块」，而每一块里有$N_r=Q/D_r$个指数；
第$r$层的第$s$个块定义为：
\[
 \mathscr{C}_{r,s}:=\bigl\{[G^{\,s+D_rz}]_P:\ 0\le z<N_r\bigr\},\qquad 0\le s<D_r
\]
也就是说，$\mathscr{C}_{r,s}$是原根表中那些行号与$s$相差$D_r$的整数倍的位置；
用编号表写：$x$属于哪一块，只看它的行号模$D_r$的余数：
\[
 x\in\mathscr{C}_{r,s}
 \iff \delta_G(x)\equiv s\pmod{D_r}
\]
这样一来，各块互不相交、合起来正好是全部$Q$个非零剩余类，
每块里有$N_r=Q/D_r$个指数；
（这些块就是由$G^{D_r}$生成的子群的陪集，
所以它们对乘法是「等距」的：$x\in\mathscr{C}_{r,s}$时，
$x$乘以$G^{D_r}$得到的还在同一块里；）
$D_0=1$时全部$Q$个指数都在一块里，这就是起点：
\[
 S_{0,0}=\omega_P^{G^0}+\omega_P^{G^1}+\cdots+\omega_P^{G^{Q-1}}=-1
\]
$D_r=H$时每块只剩两个指数，且这两个指数互为相反数，块和恰好是实迹$S_{H,s}$；\par

每个块的和记作：
\[
 S_{r,s}:=\sum_{x\in\mathscr{C}_{r,s}}\omega_P^{\,x}
\]
$r$从$0$增到$g$的过程，就是把$-1$这一块逐层细分、直到分成$H$个实迹的过程；
这条链上每一层对应一个中间数域：
\[
 \QQ=\mathscr{K}_1\subset\mathscr{K}_{D_1}\subset\cdots\subset\mathscr{K}_H\subset\QQ(\omega_P)
\]
每一步的相对次数正好是$L_r$；相对次数为$L_r$的意思是，
第$r$层的全部$D_r$个块和，可以用第$r-1$层的$D_{r-1}$个块和加上一个$L_r$次根式写出；
这就是「每层只开一个$L_r$次方」的来历；\par

块和之间的分裂关系是本算法的出发点：
\[
 \mathscr{C}_{r-1,s}=\coprod_{b=0}^{L_r-1}\mathscr{C}_{r,s+D_{r-1}b},
 \qquad
 S_{r-1,s}=\sum_{b=0}^{L_r-1}S_{r,s+D_{r-1}b}
\]
也就是说，一个粗块恰好分成$L_r$个细块，细块序号$b$只走$0,\dots,L_r-1$，
对应的块号是$s+D_{r-1}b$而不是$b$；
第$r$层要解决的问题就是：已知若干个粗块，恢复其中更细的分组；\par

\begin{proposition}
对$1\le r\le g$，记$\mathscr{K}_{D_r}=\QQ(S_{r,0})$，则$[\mathscr{K}_{D_r}:\QQ]=D_r$，
且$S_{r,0},S_{r,1},\dots,S_{r,D_r-1}$构成$\mathscr{K}_{D_r}$在$\QQ$上的一组基；
又$\mathscr{K}_{D_{r-1}}\subset\mathscr{K}_{D_r}$，$\mathscr{K}_{D_r}/\mathscr{K}_{D_{r-1}}$是次数为$L_r$的循环扩张；
当$D_r\mid H$时每块中的根成共轭对出现，所以所有的$S_{r,s}$都是实数；
\end{proposition}
\begin{proof}
自同构$\sigma:\omega_P\mapsto\omega_P^{G}$把$S_{r,s}$送到$S_{r,s+1}$，而$\sigma^{D_r}$固定各块和；
若$\sum_sc_sS_{r,s}=0$，则多项式$\sum_sc_s\sum_{x\in\mathscr{C}_{r,s}}y^x$的次数不超过$P-1$、
常数项为零且在$y=\omega_P$处取零，由$\mu_P=1+y+\cdots+y^{P-1}$的不可约性它只能是零，
故各$c_s=0$；于是所有$S_{r,s}$互不相同又线性无关，
$S_{r,0}$有恰好$D_r$个共轭，$[\mathscr{K}_{D_r}:\QQ]=D_r$；
包含关系与相对次数由循环群的子群关系得到；实性来自$-1=G^{H}$落在每块所对的子群中；
\end{proof}

最上面一层$D_g=H$的块特别简单，因为$\mathscr{C}_{H,s}=\{G^s,\,P-G^s\}$：
\[
 S_{H,s}=\omega_P^{G^s}+\omega_P^{-G^s}=2\cos\frac{2\pi G^s}{P}
\]
这$H$个数就是最终要恢复的实迹。\par

\subsubsection{块乘法的计数系数}
要让「用上一层的块算出这一层的块」成为可能，得先知道块与块相乘会发生什么；
我们把两个块和相乘，共$N_r^2$项$\omega_P^{\,x+y}$，再按指数合并同类项；
记两个纯计数：
\[
 \textcyrillic{Б}_{r,s_1,s_2,s_3}
 =\#\{(x,y):\ x\in\mathscr{C}_{r,s_1},\ y\in\mathscr{C}_{r,s_2},\ [x+y]_P\in\mathscr{C}_{r,s_3}\},
\]
\[
 \beta_{r,s_1,s_2}
 =\#\{(x,y):\ x\in\mathscr{C}_{r,s_1},\ y\in\mathscr{C}_{r,s_2},\ x+y\equiv0\pmod P\}.
\]
前者数「$x+y$落在第$s_3$块里」的有序对，后者数$x+y\equiv0$的有序对；
后一类贡献的是$\omega_P^0=1$，不落在任何块和里，必须单独记账；
各块是$G^{D_r}$的陪集，计数对乘$G^{D_r}$不变，
所以落在$\mathscr{C}_{r,s_3}$里的对数平摊到该块的$N_r$个指数上，每个指数被数到同样多次，故
\[
 S_{r,s_1}S_{r,s_2}
 =\sum_{s_3=0}^{D_r-1}\frac{\textcyrillic{Б}_{r,s_1,s_2,s_3}}{N_r}\,S_{r,s_3}
 +\beta_{r,s_1,s_2}.
\]
两个系数都是整数（$\textcyrillic{Б}$被$N_r$整除）；再由$\sum_{s_3}S_{r,s_3}=-1$，
即$1=-\sum_{s_3}S_{r,s_3}$，把常数项并入块和：
\begin{equation}
\label{eq:blockmul}
 S_{r,s_1}S_{r,s_2}
 =\sum_{s_3=0}^{D_r-1}
 \Bigl(\frac{\textcyrillic{Б}_{r,s_1,s_2,s_3}}{N_r}-\beta_{r,s_1,s_2}\Bigr)S_{r,s_3}.
\end{equation}
这就是第$r$层块和的整数系数乘法表；$\textcyrillic{Б}$与$\beta$都能逐对数目出来，
不需要知道任何一个单位根的数值。例如十七次第一层，把$8\times8$对指数逐一相加：
\[
 S_{1,0}^2+S_{1,1}^2=17,\qquad S_{1,0}+S_{1,1}=-1 .
\]

\subsubsection{周期方程：由计数写出块和的极小多项式}
固定$s_1$，把式~\eqref{eq:blockmul}的系数排成$D_r$阶整数矩阵
\[
 (\Lambda_{r,s_1})_{s_3,s_2}
 =\frac{\textcyrillic{Б}_{r,s_1,s_2,s_3}}{N_r}-\beta_{r,s_1,s_2},
\]
它就是「乘以$S_{r,s_1}$」在$\QQ$-基$S_{r,0},\dots,S_{r,D_r-1}$下的矩阵
（这组基的线性无关性见前面对$[\mathscr{K}_{D_r}:\QQ]=D_r$的证明）；
于是$S_{r,s_1}$在$\QQ$上的极小多项式就是它的特征多项式：
\begin{equation}
\label{eq:period}
 \Phi_r(y)=\det\bigl(yI-\Lambda_{r,s_1}\bigr)
 =y^{D_r}+c_1y^{D_r-1}+\cdots+c_{D_r}.
\end{equation}
系数全是有理整数，只由计数算出，不含任何根式；
$\Phi_r$的$D_r$个根正是第$r$层的全部块和$S_{r,s}$；
这就是第$r$层的\textbf{Gauss 周期方程}：每个块和都满足它，
而它本身可以逐对数出来，不必先解任何方程；三个小例：
\[
 \begin{aligned}
  P=7,\ r=1:\quad & y^3+y^2-2y-1,\\
  P=13,\ r=1:\quad & y^3+y^2-4y+1,\\
  P=17,\ r=3:\quad & y^8+y^7-7y^6-6y^5+15y^4+10y^3-10y^2-4y+1 .
 \end{aligned}
\]
最后一条正是正十七边形所依赖的周期方程。

\subsubsection{两个预制映射}
余下的计算都在「$P$行的表」上进行：一张表$\zeta$的行号是模$P$的指数$0\le x<P$，
第$x$行的值是$\zeta(x)$；
我们只需要画一张$P$行的表格，不必真的写满$P$行，
因为下文每张表都只在少数几行上取非零值；
这两个映射都定义在这个表的集合上，
定义式里的$\zeta$与$\eta$是任意表的占位符（行号都取$0\le x<P$），
$f$是求和用的第二行号（同样是占位符，$0\le f<P$），它们只在定义式中出现，用完即弃；

\textbf{（1）折叠算子（一元，带参数）；}对任意表$\eta$与$\zeta$规定
\[
 \Psi_\eta(\zeta)(x)=\sum_{f=0}^{P-1}\zeta(f)\,\eta\bigl([x-f]_P\bigr),\qquad 0\le x<P
\]
并把它的迭代记为
\[
 \Psi_\eta^{0}(\zeta)=\zeta,\qquad
 \Psi_\eta^{\,n+1}(\zeta)=\Psi_\eta\bigl(\Psi_\eta^{\,n}(\zeta)\bigr)
\]
手算的做法是：只看表$\zeta$的非零行，若第$f$行取到非零值，
就把表$\eta$向右平移$f$格（越界的行号减$P$）、乘上该行的值，
最后把所有非零行这样得到的结果叠加；
$\Psi_\eta$是「沿表$\eta$折叠一次」，$\Psi_\eta^{\,n}$是折叠$n$次；
下文只用它算三张表，代入的具体表见后；\par

\textbf{（2）求值泛函；}对任意表$\zeta$与$0\le r\le g$规定
\[
 \Xi_r(\zeta)=\zeta(0)+\sum_{s=0}^{D_r-1}\zeta(G^s)S_{r,s}
 =\sum_{s=0}^{D_r-1}\bigl(\zeta(G^s)-\zeta(0)\bigr)S_{r,s}
\]
两个式子的相等只用到$\sum_{s=0}^{D_r-1}S_{r,s}=-1$；
当$\zeta$在每个$\mathscr{C}_{r,s}$上取常值时，$\Xi_r(\zeta)$就是把表$\zeta$与第$r$层的块和配起来的结果；
我们只需要查$\zeta$在$0$行与$G^0,G^1,\dots,G^{D_r-1}$这$D_r$行的值；
特别地$r=0$时
\[
 \Xi_0(\zeta)=\bigl(\zeta(G^0)-\zeta(0)\bigr)S_{0,0}=\zeta(0)-\zeta(1)
\]
这个常数项容易被漏掉，而它往往决定整层的正负号，十七次第一层就是这种情形。\par

\subsubsection{三个表与符号的闭式写法}
上文对权重表只给了逐块赋值的描述，这里补上闭式，以便后面写公式时直接引用；

\textbf{指示函数；}整数条件$\mathscr{P}$成立时取$1$、不成立时取$0$：
\[
 \iota_{\mathscr{P}}=
 \begin{cases}
 1, & \mathscr{P}\ \text{成立},\\
 0, & \text{否则},
 \end{cases}
\]
特别地，「$x$属于集合$\mathscr{S}$」就写成$\iota_{x\in\mathscr{S}}$；\par

\textbf{权重表；}第$r$层第$j$个权重表$\chi_{r,j}$在第$b$个细块上取$\omega_{L_r}^{\,jb}$、
在其余位置取$0$；它有两种等价的闭式写法；
一是写成细块指示表$\rho_{r,b}$的线性组合，
\[
 \chi_{r,j}(x)=\sum_{b=0}^{L_r-1}\omega_{L_r}^{\,jb}\,\rho_{r,b}(x),\qquad
 \rho_{r,b}(x)=\prod_{\substack{0\le u<D_r\\ u\ne D_{r-1}b}}
 \frac{[\delta_G(x)]_{D_r}-u}{D_{r-1}b-u}
\]
第二个式子是拉格朗日插值：它的分子是「$x$的块号与第$u$块之差」、
分母是「第$D_{r-1}b$块与第$u$块之差」；
当$x$落在$\mathscr{C}_{r,D_{r-1}b}$里时分子分母逐项相等、连乘为$1$，
而当$x$落在别的块里时，连乘中总有一个因子为零、整个括号为$0$；
实际计算时不必真的去做这个连乘，按前一句话逐块赋值即可；
闭式的用处是让公式里不必再写「若$x$属于某块则取某值」这样的分句；
二是直接对指数本身插值：
\[
 \chi_{r,j}(x)=\sum_{t=0}^{D_r-1}
 \Bigl(\prod_{\substack{a=0\\a\ne t}}^{D_r-1}
 \frac{x-G^{\,t+D_{r-1}a}}{G^{\,t+D_{r-1}t}-G^{\,t+D_{r-1}a}}\Bigr)\omega_{L_r}^{\,jt}
\]
第$t$个括号里，分子是「$x$与第$t$块内第$a$个指数的差」，
分母是「第$t$块自己的两个指数之差」；当$x$落在$\mathscr{C}_{r,t}$里时，
分子分母的连乘逐个相等、比值为$1$，而当$x$落在别的块里时，
连乘中总有一个因子为零、整个括号为$0$；
于是整个和恰好实现「$x$落在第$t$块就乘$\omega_{L_r}^{jt}$」；\par

权重表还有一处要点：它只在$\mathscr{C}_{r-1,0}$上有非零值，
即只在第$r-1$层第一个粗块被拆开的那$L_r$个细块上取值，
其余$P-1-|\mathscr{C}_{r-1,0}|$行全是$0$；\par

\textbf{虚部符号；}收尾公式里要用一个只取$\pm1$的量：
\[
 \varepsilon_t=\operatorname{sgn}\sin\frac{2\pi t_0}{P},\qquad t_0=[t]_P\ne0
\]
其中$\operatorname{sgn}$是符号函数；
这里用三角函数只是为了让写法紧凑，它只负责给出符号，不参与任何数值计算；
本算法本身只用四则运算与开方。\par

\subsubsection{每层要做的六步}
有了块乘法与两个预制映射，第$r$层的计算可以固定成六步；

\textbf{第~0~步：两张表；}按$P$求出原根$G$，
写出原根表$G^a$与编号表$\delta_G$；
再把$H$分解成$L_1L_2\cdots L_g$，写出$D_0,D_1,\dots,D_g$；\par

\textbf{第~1~步：起点；}写$S_{0,0}=-1$；\par

\textbf{第~2~步：写陪集与权重表；}按定义列出第$r$层的$D_r$个块$\mathscr{C}_{r,s}$，
然后给那些「落在同一个细块$\mathscr{C}_{r,D_{r-1}b}$里的指数」赋值$\omega_{L_r}^{jb}$（$0\le b<L_r$），
其余位置赋$0$；逐块赋值的写法与取值表见「第 2 步之二：填权重表」；
其中$\omega_{L_r}$是本层的辅助单位根，$L_r\le5$时直接查低次单位根表展开成根式，
$L_r\ge7$时对它自身再走一遍本文的算法（次数更小，必然终止）；
填表时把$\omega_{L_r}^{jb}$换成实部加虚部乘$\ii$的根式，
后面所有加法都在实部与虚部里各自进行；
权重表只在本层要拆开的那些块上取值，
它把第$r-1$层的第$0$块按$L_r$个新块重新标注了一遍；\par

\textbf{第~3~步：折叠、求值、写开方量；}把权重表按上一层的块结构折叠若干次，
得到两张只依赖$x$的表$\varphi_r$与$\xi_{r,j}$，再用上一层的块和把它们求值：
\[
 F_r=\Xi_{r-1}(\varphi_r),\qquad
 T_{r,j}=\Xi_{r-1}(\xi_{r,j}),\qquad
 Y_{r,j}=\frac{T_{r,j}}{F_r}
\]
求和只跑遍第$r-1$层的$D_{r-1}$个块，
而$\varphi_r$、$\xi_{r,j}$的每一个值都是有限计数
（把$x$写成若干个落在指定块里的指数之和的方法数），
所以$F_r,T_{r,j}$都是上一层块和的有理式，不含任何未知量；
当$L_r\le3$时，由共轭配对可以证明只有$j=1$这一个非平凡分量，
于是$T_{r,j}$不必计算，$Y_{r,j}$也不用算；\par

\textbf{第~4~步：开方与反演；}取一个$R_r$使$R_r^{L_r}=F_r$，
再由有限 Fourier 反演得到第$0$块里的$L_r$个新块和$X_{r,0},\dots,X_{r,L_r-1}$，
反演式见后文「一层下降」定理；
这一步把前几步的量一次用尽：$R_r$提供被开方数，$Y_{r,j}$提供各分量之间的比例，
$\omega_{L_r}^{-js}$把分量重新合成块和；
这里有一个必须遵守的约定：$R_r$可以任取一个$L_r$次根，
但$Y_{r,j}R_r^{\,j}$不能各自另取分支；
$L_r$个分量来自同一个根$R_r$的幂次，
若把每个分量都当成独立的根号来选，得到的$X_{r,s}$一般不是实数，
整层的块和就乱了；$L_r=3$时两个非平凡分量恰好互为共轭，
这一步只需要一个共轭关系$Y_{r,2}R_r^{\,2}=\overline{R_r}$；\par

\textbf{第~5~步：恢复整层；}第~4~步只给出第$0$块里的$L_r$个块和$X_{r,0},\dots,X_{r,L_r-1}$，
第$r$层一共有$D_r$个块，剩下的是把它们按块结构平移出来；
这件事由上一层的块和加上一组有理多项式完成：
每个块和都能写成入口值$X_{r,0}$的一个多项式$\theta_{s}$，
系数由上一层块和的有理组合给出；终层（$D_r=H$时）改用递推多项式：
\[
 V_0=2,\qquad V_1=y,\qquad V_{n+1}=yV_n-V_{n-1},\qquad
 S_{H,s}=V_{u_s}(X_{g,0}),\qquad u_s=\min\bigl([G^s]_P,\,P-[G^s]_P\bigr)
\]
这一族多项式的来源很朴素：若$y=z+z^{-1}$，则$z^n+z^{-n}=V_n(y)$；
于是知道了全部实迹，$\omega_P$的每一个幂次与它的逆元的和都到手了；\par

\textbf{第~6~步：收尾；}把$H$个实迹整理成全部单位根，见后文「收尾」一节。\par

\subsubsection{三条手算约定}
上面几步里有三处容易出错的地方，这里单独提出来；

\textbf{其一，共轭分支不能独立取；}开方时只取定一个$R_r$，
同一个预解式的共轭分量由
\[
 Y_{r,L_r-1}R_r^{\,L_r-1}=\overline{R_r},
 \qquad
 \overline{R_r}=\frac{\sqrt[L_r]{F_r\overline{F_r}}}{R_r^{\,L_r-1}}
\]
给出，其中分子取正实根；分量是同一个数的幂，不是一批独立的根号：
第$j$个非平凡分量写成$Y_{r,j}R_r^{\,j}$，
把$R_r$与$\overline{R_r}$分别当作两个独立的$L_r$次根，
会得到不是实数的「块和」；\par

\textbf{其二，$L_r\le3$时不必算$\xi_{r,j}$；}因为此时共轭配对已经把非平凡分量配完，
$j=1$是唯一需要处理的分量，而$T_{r,1}=F_r$，即$Y_{r,1}=1$，
所以$L_r=2$时只有$Y_{r,1}R_r=R_r$，
$L_r=3$时只有$Y_{r,1}R_r=R_r$与$Y_{r,2}R_r^{\,2}=\overline{R_r}$，
中间没有别的分量；$L_r\ge5$时才需要$T_{r,2},\dots,T_{r,L_r-2}$；\par

\textbf{其三，$L_r=2$时块内只有两个成员；}此时两个细块的和恒等于$S_{r-1,s}$，
所以恢复步退化为「取与$X_{r,0}$相容的那一个成员」，
不必真的去解矩阵；十七次的三层全是这种情形，
这也是它全由二次根式构成的原因。\par

\subsubsection{与另外两条路线的关系}
周期下降法与半阶法用的是同一族预解式，只是分组尺度不同；
半阶法把$H$个实迹一次变换，于是出现一个$H$次根；
周期下降法先把实迹按$D$个一组并成粗块，粗块自身已经在次数更小的域里，
再把每一层的剩余循环交给本层的局部预解式；

把$H=DL$里的$L$次循环单独拿出来，令$\omega_L=\omega_H^{\,D}$，
就得到第$s$组的局部预解式$K_{s,j}$；它与后文「与半阶法的换算」一节里的$K_{r,s,j}$
在最上面一层（$r=g$）是同一个量，那里给出完整的下标写法与换算关系；
半阶法那个$H$次根，可以沿分解链逐段拆开：
例如$H=8=2\cdot2\cdot2$时，一个$8$次根可以换成一串二次根；
$H=6=3\cdot2$时，一个$6$次根可以换成三次根再套二次根；
这正是上一节所说「把开不动的大根号换成开得动的小根号」的具体含义，
完整的换算关系见后文「与半阶法的换算」一节。
\subsection{第$r$层上逐步计算的展开}
这一节把六步里的第~2~步到第~4~步展开；此时手边已有上一层（第$r-1$层）的全部
$D_{r-1}$个块和$S_{r-1,0}$到$S_{r-1,D_{r-1}-1}$，
要算出的是第$r$层块$0$里的$L_r$个块和$X_{r,0},\dots,X_{r,L_r-1}$。\par

\subsubsection{第 2 步之一：写块表}
第$r$层的第$s$块从$G^s$出发，每次乘$G^{D_r}$，共取$N_r=Q/D_r$个：
\[
 \mathscr{C}_{r,s}=\bigl\{G^s,\;G^{s+D_r},\;G^{s+2D_r},\;\dots,\;G^{s+(N_r-1)D_r}\bigr\}
 \quad(\text{行号按模 }P\text{ 取值}).
\]
我们不必把$D_r$个块全部写出：本层真正要用到的只有$\mathscr{C}_{r,0}$与
被拆开的那$L_r$个细块$\mathscr{C}_{r,D_{r-1}b}$（$0\le b<L_r$），
其余块只是最后核对$\sum_sS_{r,s}=-1$时才露面。

\subsubsection{第 2 步之二：填权重表}
权重表$\chi_{r,j}$只在一处有值：
\[
 x\in\mathscr{C}_{r,D_{r-1}b}\ \Longrightarrow\ \chi_{r,j}(x)=\omega_{L_r}^{\,jb},
 \qquad 0\le b<L_r
\]
其余行取$0$；三层的例子写出来就是：
\begin{center}
\begin{tabular}{c|c|c|c|c|c}
\toprule
细块 $b$ & 块 & $\chi_{r,1}$ & $\chi_{r,2}$ & $\chi_{r,3}$ & $\cdots$ \\
\midrule
$0$ & $\mathscr{C}_{r,0}$ & $1$ & $1$ & $1$ & $\cdots$ \\
$1$ & $\mathscr{C}_{r,D_{r-1}}$ & $\omega_{L_r}$ & $\omega_{L_r}^{2}$ & $\omega_{L_r}^{3}$ & $\cdots$ \\
$2$ & $\mathscr{C}_{r,2D_{r-1}}$ & $\omega_{L_r}^{2}$ & $\omega_{L_r}^{4}$ & $\omega_{L_r}^{6}$ & $\cdots$ \\
其余行 & — & $0$ & $0$ & $0$ & $0$ \\
\bottomrule
\end{tabular}
\end{center}
（$\chi_{r,j}$的行号仍然是模$P$的指数，上表只是把每个细块上的取值抄出来；）
填表时把$\omega_{L_r}^{jb}$按低次单位根表换成实部加虚部乘$\ii$的根式，
后面所有加法都在实部与虚部里各自进行。

\subsubsection{第 3 步：折叠、求值与系数}
先算折叠表；手算顺序是沿表$\chi_{r,1}$一次一次折叠：
\[
 \Psi_{\chi_{r,1}}(\iota_{x=0}),\quad
 \Psi_{\chi_{r,1}}^{2}(\iota_{x=0}),\quad\dots,\quad
 \varphi_r=\Psi_{\chi_{r,1}}^{\,L_r}(\iota_{x=0})
\]
每一步都按折叠算子的做法平移叠加；得到的表有一个便于验算的性质：
它在每一个$\mathscr{C}_{r-1,s}$上是常数，即
\[
 \varphi_r(G^{D_{r-1}}x)=\varphi_r(x),\qquad 0\le x<P
\]
若某张表不满足这条，说明折叠里有一行算错了；
同理，交叉表$\xi_{r,j}=\Psi_{\chi_{r,1}}^{\,L_r-j}(\chi_{r,j})$也在每个$\mathscr{C}_{r-1,s}$上是常数；
按第二条手算约定，$L_r\le3$时这一步可以整段跳过；

然后只查$0$行与$G^s$（$0\le s<D_{r-1}$）这几行，做求值；
把求值泛函$\Xi_{r-1}$写成展开式就是：
\begin{align}
 F_r&=\sum_{s=0}^{D_{r-1}-1}\Bigl(\varphi_r(G^s)-\varphi_r(0)\Bigr)S_{r-1,s},\\
 T_{r,j}&=\sum_{s=0}^{D_{r-1}-1}\Bigl(\xi_{r,j}(G^s)-\xi_{r,j}(0)\Bigr)S_{r-1,s},
 \qquad Y_{r,j}=\frac{T_{r,j}}{F_r}.
\end{align}
$F_r$与$T_{r,j}$都是上一层块和的有理系数组合，没有任何新根号；
$Y_{r,1}=1$恒成立，可以用它做一次自检；\par

\begin{theorem}[一层下降]
按上面两条式子算出的$F_r$非零；
取一个取定的$L_r$次根$R_r$使$R_r^{\,L_r}=F_r$，则
\[
 X_{r,s}=\frac1{L_r}\Bigl(S_{r-1,0}
 +\sum_{j=1}^{L_r-1}\omega_{L_r}^{-js}\,Y_{r,j}R_r^{\,j}\Bigr),
 \qquad 0\le s<L_r
\]
这$L_r$个数恰好是$S_{r,0},S_{r,D_{r-1}},\dots,S_{r,(L_r-1)D_{r-1}}$
的一个循环排列，特别地每个都是实数；
\end{theorem}
\begin{proof}
令$\tau=\sigma^{D_{r-1}}$；它把$S_{r,D_{r-1}b}$送到$S_{r,D_{r-1}(b+1)}$，
所以本层的预解式$\sum_b\omega_{L_r}^{jb}S_{r,D_{r-1}b}$是$\tau$的特征向量，特征值为$\omega_{L_r}^{-j}$；
于是它的$L_r$次幂、以及它与自己一次幂之积，都被$\tau$固定，
落在上一层块和域与本层辅助根域的合成域中，而按求值泛函求值恰好得到
$F_r$与$T_{r,j}$，故$Y_{r,j}=T_{r,j}/F_r$，其中$R_r$取$F_r$的某一支$L_r$次根；
把这一支代入反演式并作有限 Fourier 反演，
得$X_{r,s}$是那$L_r$个块和的一个循环排列；
不同的分支只造成一次循环平移。
\end{proof}

\subsubsection{第 4 步：开方与反演}
到这里只剩一次开方，分支任意：
\[
 R_r^{\,L_r}=F_r .
\]
共轭分量不另取根，按第一条手算约定由$R_r$定出；
$L_r\le3$时分量表就齐了，$L_r\ge5$时中间分量用$Y_{r,j}R_r^{\,j}$；
最后按反演式写出$L_r$个$X_{r,s}$，
取$X_{r,0}$交给恢复步骤；

\begin{proposition}
不管取哪个分支，都有
\[
 \sum_{s=0}^{L_r-1}X_{r,s}=S_{r-1,0},
 \qquad
 X_{r,s}\ \text{是实数}.
\]
\end{proposition}
\begin{proof}
第一式由$\sum_{s=0}^{L_r-1}\omega_{L_r}^{-js}=0$（$j\ne0$）得到；
实性由$\overline{Y_{r,j}R_r^{\,j}}=Y_{r,L_r-j}R_r^{\,L_r-j}$与$\omega_{L_r}^{-(L_r-j)s}$的配对得到。
\end{proof}
\subsection{层与层之间的迭代：整层的恢复}
上一节的输出只是块$0$里的$L_r$个块和，
而做下一层时求值式要用到第$r$层的全部$D_r$个块和，
所以必须把剩下的$D_{r-1}-1$个粗块补齐；
补齐的过程只用有理数四则运算，不增加新根号；
下面按「上一层交出什么、这一层还回什么」的顺序写清楚。

\subsubsection{上一层交出来的清单}
第$r-1$层做完后，手边有这些量，下面每一条都带上一层或本层的下标：
\begin{enumerate}
 \item 上一层的块数$D_{r-1}$与本层的$D_r=D_{r-1}L_r$；
 \item 本层的辅助单位根$\omega_{L_r}$；
 \item 上一层的完整块和表$S_{r-1,0},S_{r-1,1},\dots,S_{r-1,D_{r-1}-1}$；
 \item 上一层已经验算过的块内和关系$S_{r-2,s}=\sum_bS_{r-1,s+D_{r-2}b}$；
\end{enumerate}
$r=1$时手边只有$S_{0,0}=-1$一个数，$D_0=1$。

\subsubsection{这一层要还回去的清单}
\begin{enumerate}
 \item 本层的完整块和表$S_{r,0},S_{r,1},\dots,S_{r,D_r-1}$。
 \item 其中必须成立两条恒等式，作为交出去的验算：
       \[
        \sum_{s=0}^{D_r-1}S_{r,s}=-1,
        \qquad
        \sum_{b=0}^{L_r-1}S_{r,s+D_{r-1}b}=S_{r-1,s}\quad(0\le s<D_{r-1}).
       \]
\end{enumerate}

\subsubsection{块的结构}
第$r$层的$D_r$个块和按上一层分成$D_{r-1}$个粗块，每块$L_r$个：
\[
 \text{第 }s\text{ 个粗块}=\{S_{r,s},\,S_{r,s+D_{r-1}},\,S_{r,s+2D_{r-1}},\dots,
 S_{r,s+(L_r-1)D_{r-1}}\},\qquad 0\le s<D_{r-1}
\]
而$X_{r,0}$是第$0$个粗块中的某一个，其余$X_{r,s}$是第$0$个粗块的其余成员
（顺序由一个整体平移决定）；
所以恢复要做的是：从第$0$个粗块的一个成员出发，还原出全部$D_r$个块和。

\subsubsection{三条恢复途径}
按$D_{r-1}$与$D_r$的大小，我们走三条途径之一；

\textbf{（甲）$D_{r-1}=1$：不需要恢复；}
此时粗块只有一个，反演式的$L_1$个数就是第$1$层的全部块和：
\[
 S_{1,s}=X_{1,s}
 =\frac1{L_1}\Bigl(-1+\sum_{j=1}^{L_1-1}\omega_{L_1}^{-js}\,Y_{1,j}R_1^{\,j}\Bigr),
 \qquad 0\le s<L_1
\]
十三次与十七次的第一层都走这一条；\par

\textbf{（乙）$D_{r-1}>1$且$D_r<H$：用恢复矩阵；}
先按$\rho_{r,0}(x)=\iota_{x\in\mathscr{C}_{r,0}}$列出第$0$个细块的指示表，
再递推幂表
\[
 \lambda_{r,0}=\iota_{x=0},\qquad
 \lambda_{r,h}=\Psi_{\rho_{r,0}}(\lambda_{r,h-1})
 \quad(1\le h<D_r)
\]
它是纯计数：$\lambda_{r,h}(x)$就是把$x$写成$h$个
$\mathscr{C}_{r,0}$中元素之和（可重复、计次序）的方法数；
然后填恢复矩阵
\[
 \Theta_{r,s,h}=\lambda_{r,h}(G^s)-\lambda_{r,h}(0),
 \qquad 0\le s,h<D_r
\]
它的第$0$列全是$-1$；填好以后列出方程组
\[
 X_{r,0}^{\ h}=\sum_{s=0}^{D_r-1}\Theta_{r,s,h}S_{r,s},
 \qquad h=0,1,\dots,D_r-1
\]
其中$X_{r,0}$已知、$S_{r,s}$待求；把$D_r$个方程按$h$从小到大排好，
用消元法解出全部$S_{r,s}$即可；等价地，把$\Theta_r$求逆后直接读出多项式
\[
 \theta_s(y)=\sum_{h=0}^{D_r-1}(\Theta_r^{-1})_{h,s}\,y^{h}
 =\sum_{h=0}^{D_r-1}
 \frac{(-1)^{s+h}\det \Theta_r[\widehat s\mid\widehat h]}{\det \Theta_r}\,y^{h}
\]
其中$\Theta_r[\widehat s\mid\widehat h]$表示从$\Theta_r$中删去第$s$行与第$h$列所得的方阵的行列式；
再令
\[
 S_{r,s}=\theta_s(X_{r,0}),\qquad 0\le s<D_r.
\]
注意$\theta_0(y)=y$恒成立，因为$S_{r,0}$就是$X_{r,0}$在某个平移下的对应物；
这一条可以当作矩阵求逆的自检；$\Theta_r$的元素全是整数，行列式与逆的元素都是有理数，
所以这一步只用四则运算；十七次的第二层走这一条，那里$D_2=4$，矩阵只有四行四列；\par

\textbf{（丙）$D_r=H$：用终层递推；}
最后一层不填矩阵：最上面一层的块和$S_{H,s}$就是$\omega_P^{G^s}+\omega_P^{-G^s}$，
而$X_{r,0}$是某个$\omega_P^{n}+\omega_P^{-n}$，于是对它作多项式运算即可逐个得到：
\[
 S_{r,s}=V_{u_s}(X_{r,0}),\qquad u_s=\min\bigl([G^s]_P,\,P-[G^s]_P\bigr)
\]
其中$V_n$由$V_0=2$、$V_1=y$、$V_{n+1}=yV_n-V_{n-1}$递推，
它满足$V_n(z+z^{-1})=z^n+z^{-n}$；我们先把$V_1,\dots,V_H$列成一张系数表，
再把$X_{r,0}$代进去；$u_s$也先列成一张$H$行的表；
十三次与十七次的最后一层都走这一条。\par

\subsubsection{逐块核对块内和}
按上面任一条路径得到第$r$层全部块和之后，先核对两条恒等式
$\sum_sS_{r,s}=-1$与$\sum_bS_{r,s+D_{r-1}b}=S_{r-1,s}$，
再进入下一层；若某一个粗块的和与$S_{r-1,s}$不符，
说明该粗块的成员取成了别的粗块的成员，此时应回头检查开方分支与恢复矩阵的填入顺序。
\subsection{十三次的计算：先开三次根，再开二次根}
取$P=13$、原根$G=6$，分解链取$L_1=3$、$L_2=2$，于是
\[
 Q=12,\quad H=6,\quad D_0=1,\quad D_1=3,\quad D_2=6=H,\quad
 N_1=4,\quad N_2=2.
\]
两次开方的次数分别是$3$与$2$，都开得动；
十三次之所以适合做样板，是因为它两层各走一条不同的恢复途径。

\subsubsection{第 0 步：两张表}
原根表与编号表如下，$s$是行号：
\begin{center}
\begin{tabular}{c|cccccccccccc}
\toprule
$s$ & 0&1&2&3&4&5&6&7&8&9&10&11\\
\midrule
$G^s=[6^s]_{13}$ & 1&6&10&8&9&2&12&7&3&5&4&11\\
\bottomrule
\end{tabular}
\end{center}
把上面一行与下一行对调，得到编号表：
\begin{center}
\begin{tabular}{c|cccccccccccc}
\toprule
$x$ & 1&2&3&4&5&6&7&8&9&10&11&12\\
\midrule
$\delta_6(x)$ & 0&5&8&10&9&1&7&3&4&2&11&6\\
\bottomrule
\end{tabular}
\end{center}
两张表就是全部计算的入口，后面每一步都只查表与数数；
验算：$G^0=1$、$G^{12}=1$，且十二个数两两不同，确实是一个原根。

\subsubsection{第 1 层：$L_1=3$，$D_1=3$}
第 2~步：分块；第$1$层要把全部$12$个指数分成$3$块，每块$N_1=4$个；
按定义$\mathscr{C}_{1,s}=\{G^{s+3z},\ z=0,1,2,3\}$，即行号相差$3$的整数倍的那些位置：
\begin{center}
\begin{tabular}{c|l|l}
\toprule
$s$ & $\mathscr{C}_{1,s}$（行号） & 对应的指数 \\
\midrule
$0$ & $\{0,3,6,9\}$ & $\{1,8,12,5\}$ \\
$1$ & $\{1,4,7,10\}$ & $\{6,9,7,4\}$ \\
$2$ & $\{2,5,8,11\}$ & $\{10,2,3,11\}$ \\
\bottomrule
\end{tabular}
\end{center}
三块合起来恰好是$\{1,2,\dots,12\}$，每块四个指数，分块无误；

权重表$\chi_{1,1}$在这三块上分别取值$1$、$\omega_3$、$\omega_3^2$，其余行取$0$；
按块的顺序分组写出来：
\begin{center}
\begin{tabular}{c|cccccccccccc}
\toprule
行号 $x$ & 1&4&5&6&7&8&9&12 & 2&3&10&11\\
\midrule
$\chi_{1,1}(x)$ & $1$&$\omega_3$&$1$&$\omega_3$&$\omega_3$&$1$&$\omega_3$&$1$
 & $\omega_3^2$&$\omega_3^2$&$\omega_3^2$&$\omega_3^2$\\
\bottomrule
\end{tabular}
\end{center}

第 3~步：折叠与求值；折叠的含义是：
把$x$写成三个非零指数之和，把三处权重相乘，再把全部写法累加起来；
这个累加给出三个计数$(c_0,c_1,c_2)$，
分别表示三个指数落在第$0$、$1$、$2$块的次数，于是折叠值就是：
\[
 \varphi_1(x)=c_0+c_1\omega_3+c_2\omega_3^2
\]
我们按「三个指数各落在哪一块」分类计数即可；结果是：
\begin{center}
{\small
\begin{tabular}{c|c|c|c}
\toprule
行号 $x$ & $(c_0,c_1,c_2)$ & $c_0+c_1\omega_3+c_2\omega_3^2$ 化简 & 说明 \\
\midrule
$0$ & $(24,72,36)$ & $24+72\omega_3+36\omega_3^2=24+72\omega_3-36-36\omega_3=-12+36\omega_3$ & 三个指数和为零 \\
其余 $x$ & $(46,42,45)$ & $46+42\omega_3+45\omega_3^2=1-3\omega_3$ & 每个非零行都一样 \\
\bottomrule
\end{tabular}} 
\end{center}
两行的计数总数分别是$132$与$133$，恰好是「三个非零指数之和为零」
与「之和为给定非零值」的有序三重组个数，可以用来检查数数有没有漏；
化简时只用$\omega_3^2=-1-\omega_3$这一个关系，上表右端两栏是同一个数；
另有$\sum_x\varphi_1(x)=\bigl(\sum_x\chi_{1,1}(x)\bigr)^3=1^3=1$，
而$1\cdot(-12+36\omega_3)+12\cdot(1-3\omega_3)=1$，与此相符；\par

求值时$D_0=1$，上一层只有一个块、块和是$S_{0,0}=-1$，
所以只需要查$G^0=1$这一行：
\[
 F_1=\bigl(\varphi_1(1)-\varphi_1(0)\bigr)S_{0,0}
 =\bigl((1-3\omega_3)-(-12+36\omega_3)\bigr)\cdot(-1)
 =13\bigl(-1+3\omega_3\bigr)
\]
这里的$\varphi_1(1)-\varphi_1(0)$是关键：折叠值在$x=0$处与其余行不同，
差值正是开方量的来源；\par

第 4~步：开方与反演；把$\omega_3=\dfrac{-1+\sqrt3\,\ii}{2}$代进去，
实部与虚部分开写：
\[
 F_1=13\Bigl(-1+3\cdot\frac{-1+\sqrt3\,\ii}{2}\Bigr)
 =13\cdot\frac{-2-3+3\sqrt3\,\ii}{2}
 =\frac{-65+39\sqrt3\,\ii}{2}
\]
共轭量只需把$\sqrt3\,\ii$换成$-\sqrt3\,\ii$：
\[
 F_1\overline{F_1}
 =\frac{-65+39\sqrt3\,\ii}{2}\cdot\frac{-65-39\sqrt3\,\ii}{2}
 =\frac{4225+4563}{4}=\frac{8788}{4}=2197=13^3
\]
结果是一个正实数，即本层的开方量开方之后得到的两个分量互为共轭；
按第一条手算约定，取一个$R_1$使：
\[
 R_1^{\,3}=13(-1+3\omega_3),\qquad
 \overline{R_1}=\frac{\sqrt[3]{F_1\overline{F_1}}}{R_1^{\,2}}
 =\frac{13}{R_1}
\]
（分子是$\sqrt[3]{13^3}=13$，这一步完全由范数决定，不需要另取一个根号；）
$L_1=3\le3$，按第二条手算约定不必算$\xi_{1,j}$；
若要核对，直接数（把$x$写成两个指数之和、按块配对）可得$T_{1,2}=13$，于是
\[
 Y_{1,2}R_1^{\,2}=\frac{13}{13(-1+3\omega_3)}R_1^{\,2}
 =\frac{R_1^{\,3}}{(-1+3\omega_3)R_1}
 =\frac{13}{R_1}=\overline{R_1}
\]
与共轭分量一致；把$Y_{1,1}R_1=R_1$、$Y_{1,2}R_1^{\,2}=\overline{R_1}=\dfrac{13}{R_1}$代入反演式，
得第一层的三个块和：
\[
 S_{1,s}=\frac{-1+\omega_3^{-s}R_1+\omega_3^{-2s}\dfrac{13}{R_1}}{3},
 \qquad s=0,1,2
\]
$R_1$可以任取一个三次根，但$\dfrac{13}{R_1}$不能另取一个独立的三次根；
三块的和加起来等于$-1$，这是第一道验算；\par

第 5~步：恢复；本层$D_0=1$，上一层的块只有一个，走恢复途径（甲）：
三个块和在算出来的时候就已经代表整层了，不需要再做恢复；
不过它们之间仍有两个恒等关系，可以留作验算，下一层也要用到：
\[
 S_{1,1}=S_{1,0}^2+S_{1,0}-3,\qquad S_{1,2}=-S_{1,0}^2-2S_{1,0}+2
\]
这两条正是恢复多项式$\theta_s$在$D_0=1$时的样子：本层不需要恢复，
它们可以用来验算$S_{1,1},S_{1,2}$；\par

我们再把$S_{1,0}$写成精确根式，好让后面的验算全部在根式之间进行；
$S_{1,0}$本身是从开方量与反演算出来的，但它落在数域里之后，
可以用一条更短的路写它：由反演式与上面的两个恒等式可以验证它满足三次方程：
\[
 y^3+y^2-4y+1=0
\]
它的判别式是：
\[
 4^2-4(-4)^3+(-4)^2-4\cdot1^3-27=16+256+16-4-27=169=13^2
\]
一个完全平方，所以三个根都是实数；而这三个根正是$\omega_{13}$的三个共轭配对和；
解这个三次方程，用 Cardano 公式把它写成根式——
注意三个根都是实根，所以必须经过复数开立方，得到的式子里必然带$\ii$：
\[
 S_{1,0}=\frac{-2+\sqrt[3]{-260+156\sqrt3\,\ii}+\sqrt[3]{-260-156\sqrt3\,\ii}}{6}
\]
这里有一处必须说清楚：两个立方根各有三个取值，
必须取成互为共轭的那一对分支，右端才会等于$S_{1,0}$；
换一对分支就得到另外两个块和；
这也正是「实根式写不出来」的具体含义：
它的表达式里必然出现$\sqrt3\,\ii$，
但两项的虚部两两抵消，最终结果仍是实数；\par

另一方面也说得通：正因为$S_{1,1},S_{1,2}$能由$S_{1,0}$的二次式表示，
三次方程的三个根才会落在同一个三次域里，
这也解释了为什么本层只需要开一个三次根。\par

\subsubsection{第 2 层：$L_2=2$，$D_2=6=H$}
第 2~步：分块；第$2$层把每块再对半切开，$N_2=2$，每一块里两个指数互为相反数：
\begin{center}
\begin{tabular}{c|cccccc}
\toprule
$s$ & 0&1&2&3&4&5\\
\midrule
$\mathscr{C}_{2,s}$ & $\{1,12\}$&$\{6,7\}$&$\{10,3\}$&$\{8,5\}$&$\{9,4\}$&$\{2,11\}$\\
\bottomrule
\end{tabular}
\end{center}
权重表$\chi_{2,1}$：第一个细块$\mathscr{C}_{2,0}=\{1,12\}$上取$1$，
第二个细块$\mathscr{C}_{2,3}=\{8,5\}$上取$-1$，
其余行取$0$（$D_1=3$，所以细块是$\mathscr{C}_{2,3b}$）；

第 3~步：折叠与求值；$\varphi_2=\Psi_{\chi_{2,1}}^{2}(\iota_{x=0})$只有$4\times4=16$对要数，
可以整张列出来：
\begin{center}
\begin{tabular}{c|cccc}
\toprule
$y+z$ & $1$ & $12$ & $5$ & $8$ \\
\midrule
$1$ & $2:+1$ & $0:+1$ & $6:-1$ & $9:-1$\\
$12$ & $0:+1$ & $11:+1$ & $4:-1$ & $7:-1$\\
$5$ & $6:-1$ & $4:-1$ & $10:+1$ & $0:+1$\\
$8$ & $9:-1$ & $7:-1$ & $0:+1$ & $3:+1$\\
\bottomrule
\end{tabular}
\end{center}
表中每一项的记号是「$y+z$落在哪个指数上：两点权重之积」；
把同行的和加起来，得
\begin{center}
\begin{tabular}{c|ccccccccc}
\toprule
$x$ & 0&2&3&4&6&7&9&10&11\\
\midrule
$\varphi_2(x)$ & $4$&$1$&$1$&$-2$&$-2$&$-2$&$-2$&$1$&$1$\\
\bottomrule
\end{tabular}
\end{center}
其余行取$0$；全部数之和为零，与$\sum_x\chi_{2,1}(x)=1+1-1-1=0$相符；\par

求值时$D_1=3$，要查$G^0=1$、$G^1=6$、$G^2=10$三行：
\[
 F_2=\bigl(0-4\bigr)S_{1,0}+\bigl((-2)-4\bigr)S_{1,1}+\bigl(1-4\bigr)S_{1,2}
 =-4S_{1,0}-6S_{1,1}-3S_{1,2}
\]
再用上面两个恒等关系把$S_{1,1},S_{1,2}$换成$S_{1,0}$：
\begin{align*}
 F_2&=-4S_{1,0}-6\bigl(S_{1,0}^2+S_{1,0}-3\bigr)-3\bigl(-S_{1,0}^2-2S_{1,0}+2\bigr)\\
 &=-4S_{1,0}-6S_{1,0}^2-6S_{1,0}+18+3S_{1,0}^2+6S_{1,0}-6\\
 &=-3S_{1,0}^2-4S_{1,0}+12
\end{align*}
这里三个块和原来都是三次根式，化简之后整层只剩一个含$S_{1,0}$的二次式，
这就是「下降」的样子：每走一层，未知量少一个，表达式短一截；\par

第 4~步：开方与反演；$L_2=2$，开方取正实根$R_2=\sqrt{F_2}$，
本层两块的和是$S_{1,0}$，于是两个成员是：
\[
 X_{2,0}=\frac{S_{1,0}+R_2}{2},\qquad
 X_{2,1}=\frac{S_{1,0}-R_2}{2}
\]
它们就是$\mathscr{C}_{2,0}$与$\mathscr{C}_{2,3}$上的两个实迹；
把$F_2$按块结构写出来就是$F_2=S_{1,0}^2-4S_{1,1}$，
于是$X_{2,0}$、$X_{2,1}$正是二次方程
\[
 y^2-S_{1,0}\,y+S_{1,1}=0
\]
的两个根，这一点与$L_r=2$时「块内只有两个成员」的约定完全一致；\par

第 5~步：恢复（终层）；本层$D_2=6=H$，走恢复途径（丙）；先把$V_n$与$u_s$列成表；
$V_n$由$V_{n+1}=yV_n-V_{n-1}$递推，$V_0=2$、$V_1=y$，前几个是：
\begin{align*}
 V_2&=yV_1-V_0=y^2-2\\
 V_3&=yV_2-V_1=y(y^2-2)-y=y^3-3y\\
 V_4&=yV_3-V_2=y(y^3-3y)-(y^2-2)=y^4-4y^2+2\\
 V_5&=yV_4-V_3=y^5-5y^3+5y\\
 V_6&=yV_5-V_4=y^6-6y^4+9y^2-2
\end{align*}
这六个多项式就是切比雪夫多项式：
若$y=z+z^{-1}$，则$V_n(y)=z^n+z^{-n}$，所以$V_{u_s}(X_{2,0})$恰好是
$\omega_{13}^{G^s}+\omega_{13}^{-G^s}$，与实迹的定义一致；
下标表与对应的多项式是（表中把取值点记作$y=X_{2,0}$，以免式子写得太长）：
\begin{center}
\begin{tabular}{c|cccccc}
\toprule
$s$ & 0&1&2&3&4&5\\
\midrule
$G^s=[6^s]_{13}$ & 1&6&10&8&9&2\\
$u_s$ & 1&6&3&5&4&2\\[2pt]
$V_{u_s}$ & $y$ & $y^6-6y^4+9y^2-2$
 & $y^3-3y$ & $y^5-5y^3+5y$ & $y^4-4y^2+2$ & $y^2-2$\\
\bottomrule
\end{tabular}
\end{center}
取$X_{2,0}=\dfrac{S_{1,0}+R_2}{2}$，按$S_{2,s}=V_{u_s}(X_{2,0})$求出六个实迹；
把$F_2=-3S_{1,0}^2-4S_{1,0}+12$与两个恒等关系代进去，六个实迹的近似值为
\begin{center}
\begin{tabular}{c|cccccc}
\toprule
$s$ & 0&1&2&3&4&5\\
\midrule
$S_{2,s}$ & $1.77091$&$-1.94188$&$0.24107$&$-1.49702$&$-0.70921$&$1.13613$\\
\bottomrule
\end{tabular}
\end{center}
六个数之和为$-1$；
这六个数与$2\cos\dfrac{2\pi u_s}{13}$逐个相等，
其中$u_s$依次取$1,6,3,5,4,2$；最大的是$s=0$那一项，$S_{2,0}=X_{2,0}$。\par

\subsubsection{第 6 步：收尾}
最大的实迹是$S_{2,0}=1.77091=X_{2,0}$（它等于$2\cos\frac{2\pi}{13}$），
按最大值判定式取$v=0$，于是$W_q=S_{2,q}$；
以$t=1$为例：$t_0=1$，$\delta_6(1)=0$给出$\alpha_1=0$；
$\delta_6(2)=5$给出$\beta_1=5$；
$1\le t\le H$所以$\varepsilon_1=1$；查上一张表，$\beta_1=5$那一项是：
\[
 W_5=S_{2,5}=V_2(X_{2,0})=X_{2,0}^2-2,\qquad X_{2,0}=\frac{S_{1,0}+R_2}{2}
\]
代入收尾公式：
\[
 \omega_{13}=\frac{W_0}{2}+\frac12\sqrt{2-W_5}\,\ii
 =\frac{X_{2,0}}{2}+\frac12\sqrt{2-X_{2,0}^2+2}\,\ii
 =\frac{S_{1,0}+R_2}{4}+\frac12\sqrt{4-\Bigl(\frac{S_{1,0}+R_2}{2}\Bigr)^2}\ \ii
\]
式中的两个量都是上一层的精确根式：
\[
 S_{1,0}=\frac{-2+\sqrt[3]{-260+156\sqrt3\,\ii}+\sqrt[3]{-260-156\sqrt3\,\ii}}{6},
 \qquad R_2=\sqrt{S_{1,0}^2-4S_{1,1}}
\]
其中$S_{1,1}=S_{1,0}^2+S_{1,0}-3$；式中的两个立方根取互为共轭的那一对分支，
所以$S_{1,0}$是实数，$R_2$取正实根；把两个恒等式代进去，
平方根号内的量是$S_{1,0}$的二次式：
\[
 S_{1,0}^2-4S_{1,1}=-3S_{1,0}^2-4S_{1,0}+12
\]
前者是正数（它正是上一层二次方程的判别式），
而$4-\Bigl(\frac{S_{1,0}+R_2}{2}\Bigr)^2$同样为正，
所以实部与虚部都是实数；
最后一道验算是$\omega_{13}$落在单位圆上；把$X_{2,0}=\dfrac{S_{1,0}+R_2}{2}$
与$R_2^2=S_{1,0}^2-4S_{1,1}$代进去，两个平方项分别是：
\[
 \Bigl(\frac{X_{2,0}}{2}\Bigr)^2=\frac{S_{1,0}^2+2S_{1,0}R_2+R_2^2}{16},
 \qquad
 \frac{2-X_{2,0}^2+2}{4}=\frac{16-S_{1,0}^2-2S_{1,0}R_2-R_2^2}{16}
\]
两者相加时，含$R_2$的项与$S_{1,0}^2$的项两两抵消，只剩$\dfrac{16}{16}=1$；
因此无论$S_{1,0}$取哪一个分支，右端都自动落在单位圆上，
这正是收尾公式的性质：实部取半个实迹，虚部取补量的一半，
两者的平方和恒等于$1$；
代入收尾公式：
\[
 \omega_{13}=\frac{W_0}{2}+\frac12\sqrt{2-W_5}\ \ii
\]
其余$t$用同一式逐个算，
十三个单位根就全部到手。
\subsection{十七次的计算：三层二次根式}
十七次是最有名的一个算例，因为它恰好对应正十七边形的尺规作图；
取$P=17$、原根$G=3$，$H=8$分解成$L_1=L_2=L_3=2$，
于是：
\[
 Q=16,\quad H=8,\quad D_0=1,\quad D_1=2,\quad D_2=4,\quad D_3=8=H.
\]
三层全部是二次扩张，所以全部结果都能用嵌套的二次根式写出；
这个算例也把三条恢复途径都走了一遍：第一层走（甲），第二层走（乙），第三层走（丙）。

\subsubsection{第 0 步：两张表}
原根表如下：
\begin{center}
\begin{tabular}{c|cccccccc}
\toprule
$s$ & 0&1&2&3&4&5&6&7\\
\midrule
$G^s=[3^s]_{17}$ & 1&3&9&10&13&5&15&11\\
\midrule
$s$ & 8&9&10&11&12&13&14&15\\
\midrule
$G^s=[3^s]_{17}$ & 16&14&8&7&4&12&2&6\\
\bottomrule
\end{tabular}
\end{center}
编号表（反查）为$\delta_3(1)=0$，$\delta_3(2)=14$，$\delta_3(3)=1$，$\delta_3(4)=12$，$\delta_3(5)=5$，$\delta_3(6)=15$，
$\delta_3(7)=11$，$\delta_3(8)=10$，$\delta_3(9)=2$，$\delta_3(10)=3$，$\delta_3(11)=7$，$\delta_3(12)=13$，
$\delta_3(13)=4$，$\delta_3(14)=9$，$\delta_3(15)=6$，$\delta_3(16)=8$。

\subsubsection{第 1 层：$L_1=2$}
$\mathscr{C}_{1,0}$是二次剩余集，$\mathscr{C}_{1,1}$是非剩余集：
\[
 \mathscr{C}_{1,0}=\{1,2,4,8,9,13,15,16\},\qquad
 \mathscr{C}_{1,1}=\{3,5,6,7,10,11,12,14\}.
\]
权重表$\chi_{1,1}$在$\mathscr{C}_{1,0}$上取$1$、在$\mathscr{C}_{1,1}$上取$-1$，它就是 Legendre 符号；
$\varphi_1=\Psi_{\chi_{1,1}}(\iota_{x=0})$的两个值都能手算：
$x=0$时每一对$(f,-f)$都贡献$\chi_{1,1}(f)\chi_{1,1}(-f)$，而$-1=G^8$落在$\mathscr{C}_{1,0}$内，
所以$f$与$-f$落在同一块、两个因子同号，$16$对全是$+1$，得$\varphi_1(0)=16$；
$x=1$时求和化为$\sum_{f\ne0,1}\chi_{1,1}(f)\chi_{1,1}(1-f)$，按$f=2,\dots,16$逐项数，
$15$项中$7$项取$+1$、$8$项取$-1$，得$\varphi_1(1)=-1$；
求值时$D_0=1$，只需查$G^0=1$一行：
\[
 F_1=\bigl(\varphi_1(1)-\varphi_1(0)\bigr)S_{0,0}=\bigl(-1-16\bigr)\cdot(-1)=17
\]
这个$-\varphi_1(0)$的修正项决定了整个符号；
开方取正实根$R_1=\sqrt{17}$，按恢复途径（甲）得第一层的两个块和：
\[
 S_{1,0}=\frac{-1+\sqrt{17}}{2}=1.56155,\qquad
 S_{1,1}=\frac{-1-\sqrt{17}}{2}=-2.56155.
\]
两者之和为$-1$，乘积为$\dfrac{1-17}{4}=-4$，都是整数运算，可以作为验算。

\subsubsection{第 2 层：$L_2=2$，$D_2=4$}
细块为
\[
 \mathscr{C}_{2,0}=\{1,13,16,4\},\quad \mathscr{C}_{2,1}=\{3,5,12,14\},\quad
 \mathscr{C}_{2,2}=\{9,15,8,2\},\quad \mathscr{C}_{2,3}=\{10,11,7,6\}.
\]
权重表$\chi_{2,1}$：第一个细块$\mathscr{C}_{2,0}$上取$1$，第二个细块$\mathscr{C}_{2,2}$上取$-1$
（$D_1=2$，所以细块是$\mathscr{C}_{2,2b}$），其余行取$0$；
$\varphi_2=\Psi_{\chi_{2,1}}^{2}(\iota_{x=0})$有$8\times8=64$对，按行数后得到
\[
 \varphi_2(0)=8,\qquad
 \varphi_2(x)=-1\quad(x\in\mathscr{C}_{2,0}\cup\mathscr{C}_{2,2}),\qquad
 \varphi_2(x)=0\quad(\text{其余行}).
\]
求值时要查$G^0=1$与$G^1=3$两行：
\[
 F_2=\bigl((-1)-8\bigr)S_{1,0}+\bigl(0-8\bigr)S_{1,1}=-9S_{1,0}-8S_{1,1}
 =\frac{17-\sqrt{17}}{2}=6.43845.
\]
（用$S_{1,0}+S_{1,1}=-1$把$S_{1,1}$换成$-1-S_{1,0}$即得化简结果，
它同样是一个二次根式；）
开方得$R_2=\sqrt{F_2}=2.53741$，反演给出块$0$的两个成员：
\[
 X_{2,0}=\frac{S_{1,0}+R_2}{2}=2.04948,\qquad
 X_{2,1}=\frac{S_{1,0}-R_2}{2}=-0.48793.
\]
它们就是$\mathscr{C}_{2,0}$与$\mathscr{C}_{2,2}$上的两个实迹，记为$S_{2,0}=X_{2,0}$、$S_{2,2}=X_{2,1}$，
并满足$S_{2,0}+S_{2,2}=S_{1,0}$、$S_{2,0}S_{2,2}=-1$，与二次方程$y^2-S_{1,0}y-1=0$相符；

本层$D_1=2>1$而$D_2=4<8$，走恢复途径（乙）；先列第$0$个细块的指示表
$\rho_{2,0}=\iota_{x\in\mathscr{C}_{2,0}}$，再递推幂表，只记$0$行与$G^s$行
（$s=0,1,2,3$，即$1,3,9,10$）：
\begin{center}
\begin{tabular}{c|ccccc}
\toprule
$x$ & $0$ & $1$ & $3$ & $9$ & $10$\\
\midrule
$\lambda_{2,0}(x)$ & $1$&$0$&$0$&$0$&$0$\\
$\lambda_{2,1}(x)$ & $0$&$1$&$0$&$0$&$0$\\
$\lambda_{2,2}(x)$ & $4$&$0$&$2$&$1$&$0$\\
$\lambda_{2,3}(x)$ & $0$&$9$&$1$&$3$&$3$\\
\bottomrule
\end{tabular}
\end{center}
每行减去该行的$0$行值，得到恢复矩阵
\[
 \Theta_2=\begin{pmatrix}
 -1&1&-4&9\\
 -1&0&-2&1\\
 -1&0&-3&3\\
 -1&0&-4&3
 \end{pmatrix},
 \qquad \det \Theta_2=2 .
\]
$\lambda_{2,2}(0)=4$是把$0$写成两个$\mathscr{C}_{2,0}$中元素之和的方法数，
即四对$(1,16)$、$(16,1)$、$(4,13)$、$(13,4)$；$\lambda_{2,3}$的四个值同样按此逐条数出；\par

解方程组（或求逆）得$\theta_s$的系数：
\begin{align*}
 \theta_0(y)&=y,\\
 \theta_1(y)&=-\frac32+3y-\frac12y^3,\\
 \theta_2(y)&=-1-6y+y^2+y^3,\\
 \theta_3(y)&=\frac32+2y-y^2-\frac12y^3.
\end{align*}
代入$X_{2,0}=S_{2,0}=2.04948$：
\[
 S_{2,1}=\theta_1(X_{2,0})=-\frac32+3(2.04948)-\frac12(8.60857)=0.34415,
\]
同样得$S_{2,3}=\theta_3(X_{2,0})=-2.90570$，于是第二层四个块和齐全：
\[
 S_{2,0}=2.04948,\quad S_{2,1}=0.34415,\quad S_{2,2}=-0.48793,\quad S_{2,3}=-2.90570,
\]
四者之和为$-1$，且$S_{2,0}+S_{2,2}=S_{1,0}$、$S_{2,1}+S_{2,3}=S_{1,1}$，
与块内和的验算相符；把它们写成根式就是常见的
\[
 S_{2,0}=\frac{-1+s_1+s_2}{4},\quad S_{2,1}=\frac{-1-s_1+s_3}{4},\quad
 S_{2,2}=\frac{-1+s_1-s_2}{4},\quad S_{2,3}=\frac{-1-s_1-s_3}{4},
\]
其中$s_1=\sqrt{17}$，$s_2=\sqrt{34-2\sqrt{17}}$，$s_3=\sqrt{34+2\sqrt{17}}$，
三个都是正实数的平方根；
验算两条，全都只用根式；其一是$S_{2,0}+S_{2,2}=S_{1,0}$：
把$S_{2,0}$与$S_{2,2}$两式相加，$s_2$一正一负抵消，
剩下的恰好是$\dfrac{-1+s_1}{2}=S_{1,0}$；
其二是$S_{2,0}S_{2,2}=-1$：
\[
 S_{2,0}S_{2,2}
 =\frac{(-1+s_1)^2-s_2^2}{16}
 =\frac{(18-2s_1)-(34-2s_1)}{16}
 =\frac{-16}{16}=-1
\]
第一个等号是平方差公式，第二个等号用了$s_1^2=17$与$s_2^2=34-2s_1$，
$-2s_1$正是这样消掉的；
两条合起来，说明$S_{2,0},S_{2,2}$确实是$y^2-S_{1,0}y-1=0$的两个根，与块结构相符。\par

\subsubsection{第 3 层：$L_3=2$，$D_3=8=H$}
细块是八个共轭对：
\begin{center}
\begin{tabular}{c|cccccccc}
\toprule
$s$ & 0&1&2&3&4&5&6&7\\
\midrule
$\mathscr{C}_{3,s}$ & $\{1,16\}$&$\{3,14\}$&$\{9,8\}$&$\{10,7\}$&$\{13,4\}$&$\{5,12\}$&$\{15,2\}$&$\{11,6\}$\\
\bottomrule
\end{tabular}
\end{center}
权重表$\chi_{3,1}$：$\mathscr{C}_{3,0}=\{1,16\}$上取$1$，$\mathscr{C}_{3,4}=\{4,13\}$上取$-1$，其余行取$0$；
$\varphi_3=\Psi_{\chi_{3,1}}^{2}(\iota_{x=0})$同样只有$16$对，逐对数出
\begin{center}
\begin{tabular}{c|ccccccccc}
\toprule
$x$ & 0&2&15&5&14&3&12&8&9\\
\midrule
$\varphi_3(x)$ & $4$&$1$&$1$&$-2$&$-2$&$-2$&$-2$&$1$&$1$\\
\bottomrule
\end{tabular}
\end{center}
其余行取$0$；求值时查$G^0=1$、$G^1=3$、$G^2=9$、$G^3=10$四行：
\[
 F_3=(0-4)S_{2,0}+(-2-4)S_{2,1}+(1-4)S_{2,2}+(0-4)S_{2,3}
 =-4S_{2,0}-6S_{2,1}-3S_{2,2}-4S_{2,3}=2.82377.
\]
化简只需把上一层的四个根式代进去、再按$s_1^2=17$、$s_2^2=34-2s_1$、$s_3^2=34+2s_1$
整理，整个过程不出现任何小数，结果仍是一个嵌套二次根式：
\[
 F_3=S_{2,0}^2-4S_{2,1}=\frac{17+3s_1-s_2-2s_3}{4}.
\]
开方取正实根$R_3=\sqrt{F_3}=1.68041$，反演给出块$0$的两个成员：
\[
 X_{3,0}=\frac{S_{2,0}+R_3}{2}=1.86494,\qquad
 X_{3,1}=\frac{S_{2,0}-R_3}{2}=0.18454.
\]
本层是终层，走恢复途径（丙）；先列$u_s$表
\begin{center}
\begin{tabular}{c|cccccccc}
\toprule
$s$ & 0&1&2&3&4&5&6&7\\
\midrule
$G^s=[3^s]_{17}$ & 1&3&9&10&13&5&15&11\\
$[17-G^s]$ & 16&14&8&7&4&12&2&6\\
$u_s$ & 1&3&8&7&4&5&2&6\\
\bottomrule
\end{tabular}
\end{center}
取$X_{3,0}=1.86494$，按$S_{3,s}=V_{u_s}(X_{3,0})$得八个实迹：
\begin{center}
\begin{tabular}{c|cccccccc}
\toprule
$s$ & 0&1&2&3&4&5&6&7\\
\midrule
$S_{3,s}$ & $1.86494$&$0.89148$&$-1.96595$&$-1.70043$&$0.18454$&$-0.54733$&$1.47802$&$-1.20527$\\
\bottomrule
\end{tabular}
\end{center}
八个实迹之和为$-1$；
由$V_n$的基本性质（$y=z+z^{-1}$时$V_n(y)=z^n+z^{-n}$），
取$X_{3,0}=2\cos\dfrac{2\pi}{17}$就会有$S_{3,s}=2\cos\dfrac{2\pi u_s}{17}$，
于是八个实迹的精确表达一次到手（同样把取值点记作$y=X_{3,0}$）：
\begin{center}
\begin{tabular}{c|cccccccc}
\toprule
$s$ & 0&1&2&3&4&5&6&7\\
\midrule
$u_s$ & 1&3&8&7&4&5&2&6\\
$S_{3,s}$ & $y$ & $y^3-3y$ & $V_8(y)$ & $V_7(y)$
 & $y^4-4y^2+2$ & $y^5-5y^3+5y$ & $y^2-2$ & $V_6(y)$\\
\bottomrule
\end{tabular}
\end{center}
其中$V_7$与$V_8$由同一递推$V_{n+1}=yV_n-V_{n-1}$给出，
与十三次递推到$V_6$的做法完全一样，这里不再逐个写开；
把$y=X_{3,0}$代入这八个多项式，就是八个实迹的精确根式；
其中$V_1$那一格给出$S_{3,0}=X_{3,0}$，
把它按上一层的结果代开：
\[
 S_{3,0}=\frac{S_{2,0}+R_3}{2}
 =\frac{-1+\sqrt{17}+\sqrt{34-2\sqrt{17}}
 +2\sqrt{17+3\sqrt{17}-\sqrt{34-2\sqrt{17}}-2\sqrt{34+2\sqrt{17}}}}{8}
\]
这个数正是$2\cos\dfrac{2\pi}{17}$，而$s=4$那一格给出的
$S_{3,4}=X_{3,0}^4-4X_{3,0}^2+2$正是$2\cos\dfrac{8\pi}{17}$，
两者都是纯嵌套二次根式。

\subsubsection{第 6 步：收尾}
最大的实迹是$S_{3,0}$，于是$v=0$、$W_q=S_{3,q}$；以$t=1$为例；
收尾公式要查的两个下标是$\alpha_1=[\delta_3(1)]_8=[0]_8=0$与
$\beta_1=[\delta_3(2)]_8=[14]_8=6$，$\varepsilon_1=1$，于是：
\[
 W_0=S_{3,0}=X_{3,0},\qquad W_6=S_{3,6}=V_2(X_{3,0})=X_{3,0}^2-2
\]
代入收尾公式：
\[
 \omega_{17}=\frac{W_0}{2}+\frac12\sqrt{2-W_6}\,\ii
 =\frac{X_{3,0}}{2}+\frac12\sqrt{4-X_{3,0}^2}\,\ii
\]
实部与虚部都是$X_{3,0}$的二次式，$X_{3,0}$本身是上一式那个嵌套根式；
把$X_{3,0}$按上一层的四个块和代开，实部就是常见的那个式子：
\[
 \operatorname{Re}\omega_{17}
 =\frac{-1+s_1+s_2+2\sqrt{17+3s_1-s_2-2s_3}}{16}
\]
其中$s_1=\sqrt{17}$、$s_2=\sqrt{34-2\sqrt{17}}$、$s_3=\sqrt{34+2\sqrt{17}}$；
这个式子的来源就是前面$S_{3,0}$的那条式子：把它除以$2$，
立刻得到这里的实部，两个式子只差一个因子$2$；
虚部是$\dfrac12\sqrt{4-X_{3,0}^2}$，展开后同样是嵌套二次根式，
化简时也只用$s_1^2=17$、$s_2^2=34-2s_1$、$s_3^2=34+2s_1$三个关系；
单位圆上的验算与上一节相同：两个平方项相加时含$X_{3,0}$的项刚好抵消，
结果恒为$1$；
十七次全部由二次根式构成，这正是正十七边形可以尺规作图的代数原因。
\subsection{收尾：从实迹到每个单位根的实部与虚部}
走完最后一层之后，手边有$H$个实数$S_{H,s}$（$0\le s<H$）；
把这一族实数整理成全部单位根，只需三步，全用查表与四则运算完成；
每个$t$都是独立算出来的，不需要逐次相乘；

\textbf{第一步：定出最大实迹；}目的是把编号对齐，让「哪一个实迹对应$\omega_P$」这件事
不依赖任何关于角的大小的判断；做法可以完全代数地写出来：
\[
 \prod_{\substack{0\le c<H\\c\ne s}}
 \frac12\Bigl(1+\frac{S_{H,s}-S_{H,c}}{\sqrt{(S_{H,s}-S_{H,c})^2}}\Bigr)
 =\begin{cases}1,&S_{H,s}\ \text{是最大的},\\0,&\text{其余情形},\end{cases}
\]
分母的平方根取非负实根，所以$\dfrac{S_{H,s}-S_{H,c}}{\sqrt{(S_{H,s}-S_{H,c})^2}}$
在$S_{H,s}>S_{H,c}$时等于$1$、在$S_{H,s}<S_{H,c}$时等于$-1$；
只要$S_{H,s}$不是最大的，括号里就有一个因子为$0$；
这$H$个数互不相同（因为$z+z^{-1}$相同只可能是$z$与$z^{-1}$），所以最大值唯一；
把使上式左端等于$1$的那个下标记为$v$；
两个实数的大小也可以用
\[
 \max(y_1,y_2)=\frac{y_1+y_2+\sqrt{(y_1-y_2)^2}}{2}
\]
完全代数地写出来，不需要比较角的大小；\par

\textbf{第二步：把整层平移到标准位置；}定义
\[
 W_q=S_{H,[v+q]_H},
 \qquad 0\le q<H
\]
它是把最大实迹按$\omega_P$的指数整体平移$q$格的结果；按最上面一层的块和公式，$W_q$正是
\[
 W_q=\omega_P^{G^q}+\omega_P^{-G^q}
\]
其中$G$仍是开始选定的那个原根；即两个下标之间只差一次整体平移，
而这次平移恰好把编号对准标准的幂指数；\par

\textbf{第三步：逐个写出单位根；}对$0\le t<P$，记$t_0=[t]_P$（于是$0\le t_0<P$）；
$t_0\ne0$时查编号表得
\[
 \alpha_t=[\delta_G(t_0)]_H,\qquad \beta_t=[\delta_G([2t_0]_P)]_H,
 \qquad
 \varepsilon_t=\begin{cases}1,&1\le t_0\le H,\\-1,&H<t_0<P,\end{cases}
\]
其中$\delta_G$的定义域是$1\le x\le Q$，所以$t_0=0$时不查表；
两个下标都按模$H$取余，故$\alpha_t,\beta_t\in[0,H)$，与符号函数给出的$\varepsilon_t$一致；
于是
\[
 \boxed{\ \omega_P^{t}=\frac{W_{\alpha_t}}{2}
 +\frac{\varepsilon_t}{2}\sqrt{2-W_{\beta_t}}\ \ii\ }
 \qquad(0\le t<P),
\]
$t_0=0$时直接写$1+0\,\ii$；
这一式把实部与虚部分开了：实部是$\dfrac{W_{\alpha_t}}{2}$，
虚部是$\dfrac{\varepsilon_t}{2}\sqrt{2-W_{\beta_t}}$，
根号取非负实根，符号由$t_0$落在上半圈还是下半圈决定；
三个输入全部来自查表：编号表给出两张表，实迹表给出两个数值。
\subsection{计算流程的图示}
\begin{center}
% 本图用到 tikz 的 arrows.meta、positioning、calc、shapes.geometric 四个库；
% 在导言区加 \usetikzlibrary{arrows.meta,positioning,calc,shapes.geometric} 更好，这里就地载入。
\usetikzlibrary{arrows.meta,positioning,calc,shapes.geometric}
\scalebox{0.92}{%
\begin{tikzpicture}[
  node distance=3.5mm,
  font=\small,
  box/.style={draw, rounded corners=3pt, align=left, text width=92mm, inner sep=3.5pt},
  dec/.style={draw, diamond, aspect=2.4, align=center, inner sep=1pt, font=\small},
  side/.style={draw, rounded corners=3pt, align=left, text width=48mm, inner sep=3.5pt},
  term/.style={draw, rounded corners=12pt, align=center, inner sep=4pt, text width=50mm},
  note/.style={draw, dashed, rounded corners=3pt, align=left, text width=92mm,
               inner sep=3.5pt, font=\footnotesize},
  ar/.style={-{Stealth[length=2.2mm]}, semithick},
  lbl/.style={font=\footnotesize, inner sep=2pt, fill=white}
]

% ============ 主链（整体左移 22mm，给右侧分支留位置）============
\node[term, xshift=-22mm] (start) {开始：给定奇素数 $P$};

\node[box, below=of start] (s1)
 {\textbf{预备}（①）：求原根 $G$，写原根表 $G^a$ 与编号表 $\delta_G$；
  把 $H$ 分解成 $L_1\cdots L_g$，写出 $D_0=1$、$D_r=\prod_{a=1}^{r}L_a$；
  置 $r=0$，$S_{0,0}=-1$};

\node[box, below=of s1] (s2)
 {\textbf{进层}（②）：$r\leftarrow r+1$；写块与细块
  $\mathscr{C}_{r,s}=\{[G^{s+D_rz}]_P\}$（$0\le s<D_r$）、$\mathscr{C}_{r,D_{r-1}b}$（$0\le b<L_r$）；
  记住块内和 $S_{r-1,s}=\sum_b S_{r,s+D_{r-1}b}$};

\node[box, below=of s2] (s3)
 {\textbf{权重表与折叠}（②）：按 $\omega_{L_r}$ 填
  $\chi_{r,j}(x)=\omega_{L_r}^{\,jb}$（$x\in\mathscr{C}_{r,D_{r-1}b}$），其余行取 $0$；
  折叠得 $\varphi_r=\Psi_{\chi_{r,1}}^{\,L_r}(\iota_{x=0})$ 与
  $\xi_{r,j}=\Psi_{\chi_{r,1}}^{\,L_r-j}(\chi_{r,j})$};

\node[box, below=of s3] (s4)
 {\textbf{求值}（②，只用上一层块和）：$F_r=\Xi_{r-1}(\varphi_r)$，
  $T_{r,j}=\Xi_{r-1}(\xi_{r,j})$，$Y_{r,j}=T_{r,j}/F_r$};

\node[box, below=of s4] (s5)
 {\textbf{开方与反演}（②）：取 $R_r$ 使 $R_r^{\,L_r}=F_r$（$\overline{R_r}$ 由 $R_r$ 定出，不另取根）；
  $X_{r,s}=\frac1{L_r}\bigl(S_{r-1,0}
   +\sum_{j=1}^{L_r-1}\omega_{L_r}^{-js}Y_{r,j}R_r^{\,j}\bigr)$，入口值 $X_{r,0}$};

% ============ 分支一 ============
\node[dec, below=of s5] (d1) {$D_{r-1}=1$？};
\node[side, right=6mm of d1] (A)
 {\textbf{甲}\ $D_{r-1}=1$：粗块只有一个，\\
  反演已给全层 $S_{1,s}=X_{1,s}$};

% ============ 分支二 ============
\node[dec, below=9mm of d1] (d2) {$D_r=H$？};
\node[side, right=6mm of d2] (C)
 {\textbf{丙}\ $D_r=H$（终层）：\\
  $S_{r,s}=V_{u_s}(X_{r,0})$；\\
  $V_0=2$，$V_1=y$，$V_{n+1}=yV_n-V_{n-1}$};

\node[box, below=11mm of d2] (s6)
 {\textbf{乙}\ 其余情形（②）：由 $\lambda_{r,h}=\Psi_{\rho_{r,0}}(\lambda_{r,h-1})$ 填
  $\Theta_{r,s,h}=\lambda_{r,h}(G^s)-\lambda_{r,h}(0)$，
  解 $\sum_s\Theta_{r,s,h}S_{r,s}=X_{r,0}^{\ h}$ 得 $S_{r,s}=\theta_s(X_{r,0})$；
  $L_r=2$ 时取与 $X_{r,0}$ 相容的成员};

% ============ 主链连接 ============
\draw[ar] (start) -- (s1);
\draw[ar] (s1) -- (s2);
\draw[ar] (s2) -- (s3);
\draw[ar] (s3) -- (s4);
\draw[ar] (s4) -- (s5);
\draw[ar] (s5) -- (d1);
\draw[ar] (d1) -- node[lbl, right=0.5mm] {否} (d2);
\draw[ar] (d1.east) -- node[lbl, above=0.5mm] {是} (A.west);
\draw[ar] (d2) -- node[lbl, right=0.5mm] {否} (s6);
\draw[ar] (d2.east) -- node[lbl, above=0.5mm] {是} (C.west);

% ============ 汇合 ============
\coordinate (m1) at ($(s6.south)+(0,-4mm)$);
\draw[ar] (s6.south) -- (m1);
\draw[ar] (C.east) -- ++(7mm,0) |- (m1);
\draw[ar] (A.east) -- ++(7mm,0) |- (m1);

\node[dec, below=of s6, yshift=-8mm] (d3) {$D_r=H$？};
\draw[ar] (m1) -- (d3.north);

% ============ 回环：回到 ② ============
\draw[ar] (d3.west) -- ++(-40mm,0) |- (s2.west);
\node[lbl, left=2mm] at ($(d3.west)+(-8mm,0)$) {否};

% ============ 收尾 ============
\node[box, below=of d3] (s7)
 {\textbf{得出根}（③）：由 $S_{H,s}$ 定出最大实迹的下标 $v$ 与平移量
  $W_q=S_{H,[v+q]_H}$；查表得 $\alpha_t,\beta_t,\varepsilon_t$；写出每个根
  $\omega_P^{t}=\frac{W_{\alpha_t}}{2}+\frac{\varepsilon_t}{2}\sqrt{2-W_{\beta_t}}\ \ii$};
\draw[ar] (d3.south) -- node[lbl, right=1mm] {是} (s7.north);

\node[term, below=of s7] (fin) {输出全部 $\omega_P^{t}$（$0\le t<P$）};
\draw[ar] (s7.south) -- (fin.north);

% ============ 旁注：验算 ============
\node[note, below=4mm of fin] (chk)
 {\textbf{每层算完先验算}：$\sum_sS_{r,s}=-1$；$\sum_bS_{r,s+D_{r-1}b}=S_{r-1,s}$；
  $T_{r,1}=F_r$；$R_r^{\,L_r}=F_r$ 且 $R_r\overline{R_r}>0$；$\sum_sX_{r,s}=S_{r-1,0}$。\\
  \textbf{收尾后验算}：$\sum_sS_{H,s}=-1$，且每个 $\omega_P^{t}$ 的模为 $1$};

\end{tikzpicture}
}
\end{center}
\subsection{每层计算后的验算清单}
下面十条按顺序做，每层只花几分钟，可以挡住绝大多数手算错误；
\begin{enumerate}
 \item \textbf{原根表}：$G^0=1$，$G^{Q}=1$，且$G^0,\dots,G^{Q-1}$两两不同；
 \item \textbf{块划分}：第$r$层的$D_r$个块合起来恰好是$\{1,2,\dots,P-1\}$，
       每块里正好$N_r=Q/D_r$个指数；
 \item \textbf{权重表}：$\chi_{r,j}$只在$\mathscr{C}_{r-1,0}$上非零，且在$\mathscr{C}_{r,D_{r-1}b}$上恒等于$\omega_{L_r}^{jb}$；
 \item \textbf{不变性}：$\varphi_r$与$\xi_{r,j}$在每个$\mathscr{C}_{r-1,s}$上的取值都一样；
       若某两行不等，说明折叠里少叠了一次；
 \item \textbf{折叠总量}：$\sum_{x=0}^{P-1}\varphi_r(x)=\Bigl(\sum_{x=0}^{P-1}\chi_{r,1}(x)\Bigr)^{L_r}$，
       例如十三次第二层左右都是$0$，十七次第一层也是$0$；
 \item \textbf{系数}：$F_r\ne0$且$T_{r,1}=F_r$，也就是$Y_{r,1}=1$；
 \item \textbf{开方}：$R_r^{\,L_r}=F_r$；再核对$R_r\overline{R_r}$，它必须是一个正实数；
 \item \textbf{反演和}：$X_{r,0}+\cdots+X_{r,L_r-1}=S_{r-1,0}$；
 \item \textbf{块内和}：$\sum_sS_{r,s}=-1$与$\sum_bS_{r,s+D_{r-1}b}=S_{r-1,s}$两条；前者对每层都成立；
 \item \textbf{终层与收尾}：每个$S_{H,s}$落在$(-2,2)$内且互不相同，$\sum_sS_{H,s}=-1$；
       最后每个$\omega_P^t$的实部平方加虚部平方等于$1$；
\end{enumerate}

另有一条针对块多项式的验算；第$r$层第$s$个粗块的成员满足
\[
 \prod_{b=0}^{L_r-1}\Bigl(y-S_{r,s+D_{r-1}b}\Bigr)
 =y^{L_r}+\sum_{h=1}^{L_r}c_{r,h,s}\,y^{L_r-h},
\]
其中$s=0$的系数可以用计数直接写出来：
\[
 c_{r,h,0}=(-1)^{h}\,\Xi_{r-1}\Bigl(
 \sum_{0\le b_1<\cdots<b_{h}<L_r}
 \Psi_{\rho_{r,b_{h}}}\cdots\Psi_{\rho_{r,b_2}}(\rho_{r,b_1})\Bigr),
\]
其余粗块把式中每个$S_{r-1,s}$的下标整体换成$S_{r-1,[s+s_0]_{D_{r-1}}}$即可；
例如十三次第二层$s=0$的块多项式是$y^2-S_{1,0}y+S_{1,1}$，
正是那里的二次方程；十七次第三层$s=0$的块多项式是$y^2-S_{2,0}y+S_{2,1}$；
把恢复出来的成员代回这张多项式，是一次独立性很强的验算。
\subsection{与半阶法的换算}
半阶法把全部$H$个实迹一次变换：
\[
 Z^{(H)}_k=\sum_{s=0}^{H-1}\omega_H^{ks}S_{H,s},\qquad 0\le k<H,
\]
其中$\omega_H$是本原$H$次单位根；按第$g$层的分块方式把求和下标拆开，
并记第$r$层第$s$个粗块的局部预解式为
\[
 K_{r,s,j}=\sum_{b=0}^{L_r-1}\omega_{L_r}^{jb}S_{r,s+D_{r-1}b},
\]
就得到
\[
 Z^{(H)}_k=\sum_{s=0}^{D_{g-1}-1}\omega_H^{ks}\,K_{g,s,[k]_{L_g}}.
\]
即整体预解式由各粗块内的局部预解式拼成，而各个粗块之间不必各自开方，
用恢复公式彼此表达即可；
局部预解式$K_{r,s,j}$属于「比值」一类，按周期下降法的用字规则用大写拉丁字母$K$表示
（与实迹记号$W_q$区分开）；全阶法与半阶法里的$J_r$、$K_r$就是这里的$J_{m,j}$、$K_{r,s,j}$，
只是下标取法不同；\par

另一半的换算走乘法恒等式；对$m,j$定义 Jacobi 和
\[
 J_{m,j}=\sum_{x=1}^{P-1}\omega_H^{\ m\,\delta_G(x)+j\,\delta_G(1-x)},
\]
指数可以先模$L$取余；于是对素数$L$有
\[
 \bigl(Z^{(H)}_j\bigr)^{L}=\Bigl(\prod_{m=1}^{L-1}J_{mj,j}\Bigr)Z^{(H)}_{Lj}\quad(Lj<H),
 \qquad
 \bigl(Z^{(H)}_j\bigr)^{L}=P\prod_{m=1}^{L-2}J_{mj,j}\quad(Lj=H).
\]
十三次的$H=6$给出$\bigl(Z^{(H)}_2\bigr)^3=13J_{2,2}$、$\bigl(Z^{(H)}_1\bigr)^2=J_{1,1}Z^{(H)}_2$，
其中$J_{2,2}=\sum_{x=1}^{12}\omega_3^{\ \delta_6(x)+\delta_6(1-x)}=2+6\omega_3+3\omega_3^2=-1+3\omega_3$，
于是$\bigl(Z^{(H)}_2\bigr)^3=13J_{2,2}=13(-1+3\omega_3)$，与十三次第一层的$F_1$完全相等：
先开三次根得到$Z^{(H)}_2$，再开二次根得到$Z^{(H)}_1$，与那两层开方的顺序一致；
十七次的$H=8$给出$\bigl(Z^{(H)}_4\bigr)^2=17$、$\bigl(Z^{(H)}_2\bigr)^2=J_{2,2}Z^{(H)}_4$、$\bigl(Z^{(H)}_1\bigr)^2=J_{1,1}Z^{(H)}_2$，
与该算例第一层的$F_1=17$也一致；\par

两个体系应当各用各的系数：完整 Gauss 和的模平方公式不能搬到局部预解式上，
例如十三次第二层的$S_{2,0}+S_{2,3}=S_{1,0}$一般不等于$P$，
而是一般的上一层块和的代数数。
\subsection{低次单位根的取值表}
$L_r\le5$时把本层辅助单位根展开成根式，直接查用：
\begin{align*}
 \omega_2&=-1,\\
 \omega_3&=\frac{-1+\sqrt3\,\ii}{2},&
 \omega_3^2&=\frac{-1-\sqrt3\,\ii}{2},\\
 \omega_5&=\frac{-1+\sqrt5}{4}+\frac{\sqrt{10+2\sqrt5}}{4}\,\ii,&
 \omega_5^2&=\frac{-1-\sqrt5}{4}+\frac{\sqrt{10-2\sqrt5}}{4}\,\ii,\\
 \omega_5^3&=\frac{-1-\sqrt5}{4}-\frac{\sqrt{10-2\sqrt5}}{4}\,\ii,&
 \omega_5^4&=\frac{-1+\sqrt5}{4}-\frac{\sqrt{10+2\sqrt5}}{4}\,\ii
\end{align*}
这四组值可以逐条代回$\omega_n^n=1$与$1+\omega_n+\cdots+\omega_n^{n-1}=0$验算；
$L_r\ge7$时不必另立公式：把本文的整套算法对素数$L_r$自身再走一遍，
就用$L_r$的实迹写出$\omega_{L_r}^{t}$，而$L_r<P$，
所以这次递归的规模严格更小，必定终止。
\subsection{三条路线的比较}
把三条路线放在一起比较，可以看出各自的位置；
全阶法与半阶法的公式长度最短，思路最直接，
代价是把困难集中在一个高次根号里，而这个根号没有被开出来；
周期下降法的流程最长，每一层都要写陪集表、数折叠值、解一次线性方程组，
但它的每一步都落在已经求出的实数上，
最后得到的是每个单位根的实部与虚部的具体根式；\par

用一张表把三种选择摆在一起：
\begin{center}
\begin{tabularx}{\linewidth}{>{\raggedright\arraybackslash}p{0.155\linewidth}
 >{\raggedright\arraybackslash}X
 >{\raggedright\arraybackslash}X
 >{\raggedright\arraybackslash}X}
\toprule
项目 & 全阶法 & 半阶法 & 周期下降法 \\
\midrule
中间量 & $Q$个预解式$Z^{(Q)}_j$ & $H$个预解式$Z^{(H)}_k$ & 逐层$D_r$个块和$S_{r,s}$ \\
阶的关系 & $Q$ & $H=Q/2$ & 每层$L_r$，$H=\prod L_r$ \\
本层开方次数 & 一次$Q$次方 & 一次$H$次方 & 每层一次$L_r$次方 \\
系数的算法 & 完整 Jacobi 和 & 半阶 Jacobi 和 & 局部计数与有理矩阵 \\
分支控制 & 一个共同根 & 一个共同根
 & 每层一个共同根，加整层相容恢复 \\
输出 & 形式解 & 形式解 & 完整的实虚部根式解 \\
\bottomrule
\end{tabularx}
\end{center}
三种方法用的是同一族预解式、同一条反演公式，
区别只在「对谁做 Fourier」：全阶法对整个$Q$个非零剩余类做，
半阶法对$H$个实迹做，周期下降法把$H$拆成链、逐层对$L_r$个块做；
而且半阶法并不是另一套东西：它的预解式恰好是全阶法的偶数位子列，
$Z^{(H)}_k=Z^{(Q)}_{2k}$；
三十一次的例子能把这种差别说得很清楚：
$H=15=3\cdot5$，半阶法要给一个$15$次根号，
周期下降法则先求三周期层的分量，再在所得系数域上求五次根，
中间出现的都是三次与五次根式；\par

分解链的次序可以任选，例如$H=6$既可写成$3\cdot2$也可写成$2\cdot3$，
$H=15$既可写成$3\cdot5$也可写成$5\cdot3$；
两种写法都正确，但得到的系数与表达式长短不同，
某种固定次序并不保证对每个素数都最短；
当$H$本身是素数时（例如$P=23$，$H=11$），
实周期域与有理数之间没有非平凡的中间层可继续分拆，
这时周期下降法退化成只有一层，与半阶法给出的东西是同一个；
另一方面，当$H$是$2$的幂时（$P=17,257,\dots$），
每一层都是二次扩张，全部结果都能写成纯嵌套二次根式，
这正是正多边形可以尺规作图的全部情形。
\pagebreak

\section{嵌套根式可开方性的判定}
本章的判定与算法是应用在全阶法与半阶法之中的；这两条路线各自留下一个高次根号：
全阶法是开方量在$\omega_Q$处的值的$Q$次根，半阶法是它在$\omega_H$处的值的$H$次根；
那个被开方数本身是若干层根式的嵌套，公式到此为止并没有把它写成一个数；
要真正落地，必须回答：这个根号能不能开出来、开出来是什么；
本章给出这一问题的两种含义（D1 内去嵌套与 D2 深度约化）、必要条件、完备判据与六步手算流程，
并在六个算例上走完全程。

\subsection{问题的提出与两种「可去」的定义}
\label{sec:d1d2}

给定外层根式 $\sqrt[n]{\alpha}$，其中 $\alpha$ 是若干内层根式的加减组合，
需要回答两个问题：外层根号能否去掉（判定），以及去掉之后表达式是什么（求解）；
全部判定归结为两个域内在的问句，本节把它们精确化；\par

\begin{definition}[根式表达式与深度]
\label{def:depth}
有理数表达式深度为 $0$；若表达式深度 $\le t$，则其和、差、积深度 $\le t$；
若其深度为 $t$，则其 $n$ 次根（$n\ge2$）深度 $\le t+1$；
域 $\mathscr{K}_t$ 表示深度 $\le t$ 的根式所生成的域；
\end{definition}

\begin{definition}[内层域]
\label{def:inner}
给定外层根式 $\sqrt[n]{\alpha}$，其\textbf{内层域} $\mathscr{K}$ 是 $\alpha$ 的表达式中
全部内层根式生成的 $\mathbb{Q}$-域（有限数域）；$\sqrt[n]{\alpha}$ 的深度为 $t+1$，
其中 $t$ 为 $\alpha$ 的表达深度；
\end{definition}

\begin{definition}[两种去嵌套]
\label{def:two}
设 $\mathscr{K}$ 为内层域，$t$ 为 $\alpha$ 的表达深度；
\begin{enumerate}
\item \textbf{D1（内去嵌套）}：存在 $\beta\in\mathscr{K}$ 使 $\beta^n=\alpha$；
判据即 $\alpha\in\mathscr{K}^{*n}$（$\mathscr{K}^*$ 的 $n$ 次幂子群），不引入任何新根式；
\item \textbf{D2（深度约化）}：存在深度 $\le t$ 的表达式表示 $\sqrt[n]{\alpha}$ 的某个分支；
由定义~\ref{def:depth}，这等价于
\[
\sqrt[n]{\alpha}\in\mathscr{K}\bigl(\omega_n,\sqrt[n]{b_1},\dots,\sqrt[n]{b_r}\bigr),
\qquad b_j\in\mathbb{Q}^*,
\]
其中 $\omega_n=\sqrt[n]{1}$ 本身是底为有理数的一次根式，按定义合法；
两种情形的等价性由定理~\ref{thm:d2} 给出；
\end{enumerate}
\end{definition}

\begin{proposition}[内在性]
\label{prop:intrinsic}
D1 只依赖数域 $\mathscr{K}$ 与元素 $\alpha$，不依赖 $\mathscr{K}$ 的根式呈现方式；
D2 只依赖 $\mathscr{K},\alpha,n$；因此两者都是良定义的判定问题。
\end{proposition}

\subsubsection{记号}

\begin{itemize}
\item 单位根 $\omega_n=\exp(2\pi\ii/n)$，$\mathscr{U}_n=\langle\omega_n\rangle$；
特别地 $\omega_2=-1\in\mathbb{Q}$、$\omega_3=(-1+\sqrt{-3})/2$、$\omega_4=\ii$；
\item $v_{\mathfrak{p}}$ 为素理想 $\mathfrak{p}$ 上的赋值，$e(\mathfrak{p}/p)$ 为分歧指数，
$f(\mathfrak{p}/p)$ 为剩余域次数，$N_{\mathscr{K}/\mathbb{Q}}$ 为域范数，
$\mathscr{O}_{\mathscr{K}}$ 为整数环，$\operatorname{Res}$ 为结式，
$\operatorname{sqfree}$ 为无平方因子部分；
\item $d$ 记 $[\mathscr{K}:\mathbb{Q}]$，$n$ 记外层根式的次数；$\omega$ 表示任意单位根（$\omega\in\mathscr{U}_n$），$\omega_n$ 专指本原 $n$ 次单位根；
字母 $\ii$ 只表示虚数单位，求和、求积与生成元的下标用 $r,s,t$ 等，不用这个字母本身。
\end{itemize}
\subsection{预备知识：数域、嵌入、范数、赋值与单位}
\label{sec:prelim}

\subsubsection{数域、折叠与表示唯一性}

\begin{definition}[数域]
\label{def:numberfield}
数域是包含所有有理数、且对加、减、乘、除（除数非零）封闭的数集；
若 $\theta$ 满足 $\theta^m=a\in\mathbb{Q}$，则 $\mathbb{Q}(\theta)=\{c_0+c_1\theta+\dots+c_{m-1}\theta^{m-1}:c_j\in\mathbb{Q}\}$；
更一般地，$\mathbb{Q}(\sqrt3,x)$（$x^5=2$）的元素形如
\[
\alpha=(a+bx+cx^2+dx^3+ex^4)+(a'+b'x+\dots+e'x^4)\sqrt3 .
\]
\end{definition}

\textbf{折叠规则}：若 $\theta^m=a$，则对任意 $q\ge0$ 有 $\theta^{qm+r}=(\theta^m)^q\theta^r=a^q\theta^r$；
\emph{证明}：$\theta^{qm}=(\theta^m)^q=a^q$；\qed

这条规则是把高次幂压回标准形状的机制；判定之前必须先把 $\alpha$ 写进它能属于的
\emph{最小域}（能合并的根式先合并，如 $\sqrt2\cdot\sqrt3=\sqrt6$、$\sqrt8=2\sqrt2$）；\par

\begin{proposition}[表示唯一]
\label{prop:unique}
若 $a+b\sqrt2=a'+b'\sqrt2$（$a,a',b,b'\in\mathbb{Q}$），则 $a=a'$，$b=b'$；
\end{proposition}

\begin{proof}
若 $b\ne b'$，移项得 $\sqrt2=(a-a')/(b'-b)\in\mathbb{Q}$，与 $\sqrt2$ 无理矛盾
（$p^2=2q^2$ 迫使 $p,q$ 同偶，与最简分数矛盾）。
\end{proof}

\subsubsection{嵌入（影子）与范数}

\begin{definition}[影子]
\label{def:shadow}
$\mathscr{K}$ 的一个\textbf{影子}是一个保持加法与乘法、固定所有有理数的映射
$\sigma:\mathscr{K}\to\mathbb{C}$，即一个 $\mathbb{Q}$-嵌入；
直观上它把每个根号换成同一个根号的另一个分支；例如 $\mathbb{Q}(\sqrt2)$ 有两个影子：
恒等映射与 $\sigma(\sqrt2)=-\sqrt2$；
\end{definition}

\begin{lemma}[影子尊重运算]
\label{lem:shadow-hom}
$\sigma(u+v)=\sigma(u)+\sigma(v)$，$\sigma(uv)=\sigma(u)\sigma(v)$；
\end{lemma}

\begin{proof}
这正是「保持域运算」的定义本身；对「逐根式替换」式的影子，按分配律逐项验证；
\end{proof}

\begin{lemma}[影子的个数与作用]
\label{lem:shadow-count}
设 $\theta^m=a\in\mathbb{Q}$ 且 $t^m-a$ 在 $\mathbb{Q}$ 上不可约（例如 Eisenstein 判别法适用：
素数 $p\mid a$、$p^2\nmid a$、$p$ 不整除首项系数）；则 $[\mathbb{Q}(\theta):\mathbb{Q}]=m$，
$\mathbb{Q}(\theta)$ 恰有 $m$ 个影子，由 $\theta\mapsto\omega_m^j\theta$（$j=0,\dots,m-1$）给出；
\end{lemma}

\begin{proof}
$\sigma(\theta)$ 必是极小多项式 $t^m-a$ 的根，而 $t^m-a$ 的根恰为 $\omega_m^j\theta$；
反之每个根都给出一个影子（标准事实）；影子由它在生成元上的取值决定，故恰有 $m$ 个；
\end{proof}

\begin{definition}[范数]
\label{def:norm}
$N_{\mathscr{K}/\mathbb{Q}}(\alpha)=\prod_\sigma\sigma(\alpha)$，乘积取遍 $\mathscr{K}$ 的全部影子；
\end{definition}

\begin{lemma}
\label{lem:norm-rational}
$N_{\mathscr{K}/\mathbb{Q}}(\alpha)\in\mathbb{Q}$；二次情形直接验证：
$N(a+b\sqrt d)=(a+b\sqrt d)(a-b\sqrt d)=a^2-db^2\in\mathbb{Q}$；
\end{lemma}

\begin{lemma}[乘性]
\label{lem:norm-mult}
$N(\alpha\beta)=N(\alpha)N(\beta)$；特别地 $N(\alpha^n)=N(\alpha)^n$；
\end{lemma}

\begin{proof}
逐影子用引理~\ref{lem:shadow-hom}：$\sigma(\alpha\beta)=\sigma(\alpha)\sigma(\beta)$，
两组乘积逐项相等；
\end{proof}

\begin{theorem}[两项范数公式]
\label{thm:norm-two-term}
设 $\theta^m=a\in\mathbb{Q}$、$t^m-a$ 不可约，$\alpha=A+B\theta$；则
\[
N_{\mathbb{Q}(\theta)/\mathbb{Q}}(\alpha)=A^m+(-1)^{m+1}aB^m .
\]
\end{theorem}

\begin{proof}
由引理~\ref{lem:shadow-count}，$\sigma_j(\alpha)=A+B\omega_m^j\theta$，故
\[
N(\alpha)=\prod_{j=0}^{m-1}(A+B\omega_m^j\theta)=A^m-(-B\theta)^m=A^m-(-1)^maB^m,
\]
其中用了 $\prod_j(t-\omega_m^j)=t^m-1$（两边都是首一 $m$ 次多项式且根相同）；
\end{proof}

校准例：$N(\sqrt2)=0^2+(-1)^3\cdot2\cdot1=-2$；$N(2-\sqrt3)=4-3=1$；$N(\sqrt3)=-3$；\par

\textbf{多层收缩}：$\mathscr{K}=\mathbb{Q}(\sqrt3,x)$ 时，先把 $\alpha$ 写成
$A(x)+B(x)\sqrt3$，用 $N_{\mathscr{K}/\mathbb{Q}(x)}(\alpha)=A^2-3B^2$，再对 $x$ 收缩到 $\mathbb{Q}$；
若 $\deg\le2$ 而 $m=3$，则用 $\omega$-恒等式
\[
N(A+B\theta+C\theta^2)=A^3+B^3a+C^3a^2-3ABCa\qquad(\theta^3=a);
\]
最后一步（收缩到 $\mathbb{Q}$）也可用结式 $\operatorname{Res}(f_{\min},\cdot\,)$ 完成。\par

\subsubsection{赋值、素理想与无平方因子部分}

$v_p(m)$ 表示素数 $p$ 在整数 $m$ 中的幂次（$72=2^3\cdot3^2$：$v_2=3$，$v_3=2$），
$v_p(a/b)=v_p(a)-v_p(b)$，满足 $v_p(xy)=v_p(x)+v_p(y)$、$v_p(x^n)=n\,v_p(x)$；
\textbf{无平方因子部分} $\operatorname{sqfree}(m)$ 是把 $m$ 中的偶次幂因子削掉所得的整数
（$126=3^2\cdot14\Rightarrow\operatorname{sqfree}(126)=14$），于是
\[
p/q\in\mathbb{Q}^2\iff p,\ q\ \text{都是完全平方数}.
\]

在 $\mathscr{K}$ 里，普通素数 $p$ 会分裂；例如在 $\mathbb{Q}(\sqrt3,x)$（$x^5=2$）中，
\[
\mathfrak{q}_1=(3,\ \sqrt3,\ x+1)
\]
表示「使 $\sqrt3\equiv0$、$x\equiv-1\pmod{\mathfrak{q}_1}$」的广义素数；
检验其合法性：把 $\sqrt3\mapsto0$、$x\mapsto-1$ 代入关系式 $\sqrt3^2=3$、$x^5=2$，
得 $0=3\equiv0$ 与 $(-1)^5=-1\equiv2\pmod3$；商环 $\mathscr{K}/\mathfrak{q}_1\cong\mathbb{F}_3$；\par

\begin{definition}[素理想赋值]
\label{def:vq}
$v_{\mathfrak{q}}(\alpha)$ 是使 $\alpha\in\mathfrak{q}^k$ 的最大整数 $k$，
即 $\alpha$ 被 $\mathfrak{q}$ 整除的次数；分歧指数 $e(\mathfrak{q}/p)=v_{\mathfrak{q}}(p)$
表示 $p$ 在 $\mathfrak{q}$ 处被折成的层数；剩余域次数 $f(\mathfrak{q}/p)$ 是
$\mathscr{K}/\mathfrak{q}$ 作为 $\mathbb{F}_p$ 扩域的次数；恒等式
$[\mathscr{K}:\mathbb{Q}]=\sum_{\mathfrak{q}\mid p}ef$；
\end{definition}

\begin{lemma}[范数--赋值公式]
\label{lem:norm-val}
$v_p\bigl(N_{\mathscr{K}/\mathbb{Q}}(\alpha)\bigr)=\sum_{\mathfrak{q}\mid p}f(\mathfrak{q}/p)\,v_{\mathfrak{q}}(\alpha)$；
\end{lemma}

校准例：$\mathscr{K}=\mathbb{Q}(\sqrt3)$，$\alpha=\sqrt3$：$N=-3$，$v_3=1$；
右端只有 $\mathfrak{q}=(\sqrt3)$，$f=1$，$v_{\mathfrak{q}}(\sqrt3)=1$，和为 $1$，
其中 $e$ 已经藏在 $v_{\mathfrak{q}}(3)=e$ 里。\par

\subsubsection{单位与单位根}

\begin{definition}[单位]
\label{def:unit}
代数整数（极小多项式首一整系数的元素）$\varepsilon$ 称为\textbf{单位}，
若 $N(\varepsilon)=\pm1$——等价于 $\varepsilon^{-1}$ 仍是代数整数；
\end{definition}

例如 $\varepsilon=1+\sqrt[3]2+\sqrt[3]4$ 满足 $N(\varepsilon)=1$，且 $\varepsilon^{-1}=\sqrt[3]2-1$；
秩 $1$ 的域中所有单位恰为 $\pm$ 基本单位 $\varepsilon$ 的幂；\par

单位根 $\omega_n=\exp(2\pi\ii/n)$ 已见上；$\omega_2=-1\in\mathbb{Q}$，
故 $n=2$ 时没有任何单位根带来的麻烦，这是平方根判定特别干净的原因。
\subsection{必要条件与否证证书的来源}
\label{sec:necessary}

\begin{theorem}[$n$ 次幂的三个必要条件]
\label{thm:necessary}
设 $\alpha=\beta^n$，$\beta\in\mathscr{K}$，$[\mathscr{K}:\mathbb{Q}]=d$；则
\begin{enumerate}
\item \textbf{（赋值条件）}对 $\mathscr{K}$ 的每个素理想 $\mathfrak{p}$：$v_{\mathfrak{p}}(\alpha)\equiv0\pmod n$；
\item \textbf{（范数条件）}$|N_{\mathscr{K}/\mathbb{Q}}(\alpha)|=q^n$，$q\in\mathbb{Q}$；
等价地 $|N(\alpha)|=|N\mathfrak{b}|^n$，其中 $\mathfrak{b}=(\beta)$ 是主理想；
\item \textbf{（全正性）}$n$ 为偶数时，$\sigma(\alpha)\ge0$ 对每个实嵌入
$\sigma:\mathscr{K}\hookrightarrow\mathbb{R}$；
\end{enumerate}
\end{theorem}

\begin{proof}
(1) $(\alpha)=(\beta)^n$，逐素理想比较指数即得；
(2) 对 (1) 取域范数：$\bigl(N(\alpha)\bigr)=N\bigl((\alpha)\bigr)=N\bigl((\beta)\bigr)^n$，
而主理想的域范数由生成元的绝对值给出；
(3) $\sigma(\alpha)=\sigma(\beta)^n\ge0$；
\end{proof}

手算中最常用的只有范数条件一项，它可由结式（单生成元情形对
$\operatorname{Res}(f,x)$ 的嵌套）或乘法矩阵的行列式\textbf{精确}计算，
不依赖任何数值近似；\par

\begin{proposition}[\textbf{手算用法与警示}]
算出 $|N(\alpha)|$ 并作素因子分解：只要有一个素数的指数不被 $n$ 整除，就有
$\alpha\notin\mathscr{K}^{*n}$，D1 判死——这是证明而非启发式；
反之，范数条件通过\textbf{不蕴含}可去：例~\ref{ex:two-sqrt3} 的 $N(2-\sqrt3)=1$、
例~\ref{ex:ramanujan} 的 $N(x-1)=-1$ 都通过，但两者的 D1 都失败；
\end{proposition}

\textbf{注}；三个条件都只是\textbf{必要}而非充分，因此它们只能产出否证证书；
完备判定需要素理想赋值、类群与单位群（定理~\ref{thm:complete}）。
\subsection{两个主定理}
\label{sec:main}

\subsubsection{两层大公式}

\begin{theorem}[二次层的封闭公式]
\label{thm:two-level}
设 $\mathscr{F}$ 的特征 $\ne2$，$a\in \mathscr{F}$ 非平方，$\mathscr{K}=\mathscr{F}(\sqrt a)$，
$\alpha=a_0+a_1\sqrt a$（$a_j\in \mathscr{F}$），令 $s^2=a_0^2-a_1^2a$；则
\begin{enumerate}
\item \textbf{（D2）}$\sqrt{\alpha}$ 可深度约化 $\iff s\in \mathscr{F}$，此时
\[
\sqrt{\alpha}=\sqrt{\tfrac{a_0+s}{2}}\pm\sqrt{\tfrac{a_0-s}{2}},
\]
其中 $\pm$ 取使交叉项等于 $a_1\sqrt a$ 的符号；
\item \textbf{（D1）}若再有 $\tfrac{a_0+s}{2}\in \mathscr{F}^2$ 且 $\tfrac{a_0-s}{2}\in a\cdot \mathscr{F}^2$，
则得 D1；且此条件对 D1 是\textbf{完备}的；
\end{enumerate}
\end{theorem}

\begin{proof}
设 $\beta=u+v\sqrt a$（$u,v\in \mathscr{F}$）；$\beta^2=\alpha$ 等价于
\begin{equation}
\label{eq:twolvl}
u^2+av^2=a_0,\qquad 2uv=a_1 .
\end{equation}
\emph{必要性}：由~\eqref{eq:twolvl}，
$(u^2-av^2)^2=(u^2+av^2)^2-4au^2v^2=a_0^2-a_1^2a=s^2$，
而 $u^2-av^2\in \mathscr{F}$，故 $s\in \mathscr{F}$；与~\eqref{eq:twolvl} 第一式联立解出
$u^2=\tfrac{a_0+s'}2$、$av^2=\tfrac{a_0-s'}2$，其中 $s'=u^2-av^2=\pm s$：
这正是 D1 的两个条件；同时 $u=\pm\sqrt{\tfrac{a_0+s'}2}$、$v\sqrt a=\pm\sqrt{\tfrac{a_0-s'}2}$，
两个新根号的底都在 $\mathscr{F}$ 内——D2 表示自动出现；
\emph{充分性}：展开
$\bigl(\sqrt{\tfrac{a_0+s}{2}}\pm\sqrt{\tfrac{a_0-s}{2}}\bigr)^2
  =a_0\pm2\sqrt{\tfrac{a_0^2-s^2}4}=a_0\pm|a_1|\sqrt a$，
按 $a_1$ 的符号选取 $\pm$；D1 情形把 $u,v$ 回代~\eqref{eq:twolvl} 逐条验证；
\end{proof}

$\mathscr{F}=\mathbb{Q}$、$a$ 无平方因子时，D1 就是经典规则「$a_0^2-a_1^2a$ 为完全平方」；
两层题在这一步可同时得到 D1 的判定与 D2 的答案：算 $s^2=a_0^2-a_1^2a$，
看 $s$ 是否为有理数（否则两层题在此判死），再过两道平方性检查定 D1。\par

\subsubsection{Landau 单项式引理}

\begin{theorem}[Landau 单项式引理]
\label{thm:landau}
设 $\operatorname{char}\mathscr{F}=0$，$\mathscr{U}_n\subseteq \mathscr{F}$，$\mathscr{K}=\mathscr{F}(\theta_1,\dots,\theta_m)$，
$\theta_j^n=a_j\in \mathscr{F}$；若 $\beta\in\mathscr{K}$ 且 $\beta^n=\alpha\in \mathscr{F}$，
则存在 $\gamma\in \mathscr{F}$、$e_r\in\{0,\dots,n-1\}$ 与 $\omega\in\mathscr{U}_n$ 使
\[
\beta=\omega\cdot\gamma\cdot\prod_r\theta_r^{e_r},
\qquad\text{从而}\qquad
\alpha=\gamma^n\prod_r a_r^{e_r}.
\]
\end{theorem}

\begin{proof}
分五步；
\begin{enumerate}
\item 每个 $\mathscr{F}(\theta_j)/\mathscr{F}$ 是循环 Kummer 扩张：$\theta_j$ 的共轭 $\omega^m\theta_j$
全在 $\mathscr{F}(\theta_j)$ 内（因 $\mathscr{U}_n\subseteq \mathscr{F}$）；其复合 $\mathscr{K}/\mathscr{F}$ 是
Galois、交换、指数整除 $n$ 的扩张；
\item （Kummer 对应，[Washington, Ch.~14] 型标准结果）对子群
$\mathscr{D}=\langle\bar a_1,\dots,\bar a_m\rangle\le \mathscr{F}^*/\mathscr{F}^{*n}$ 有
$\mathscr{K}=\mathscr{F}(\sqrt[n]{\mathscr{D}})$，且配对
\[
\operatorname{Gal}(\mathscr{K}/\mathscr{F})\times\mathscr{D}\to\mathscr{U}_n,\qquad
(\sigma,\delta)\mapsto\sigma(\sqrt[n]\delta)/\sqrt[n]\delta
\]
是完美的；中间域 $\mathscr{L}\mapsto\mathscr{D}_\mathscr{L}:=\{\delta\in\mathscr{D}:\mathscr{F}(\sqrt[n]\delta)\subseteq \mathscr{L}\}$
是双射，且 $\mathscr{L}=\mathscr{F}(\sqrt[n]{\mathscr{D}_\mathscr{L}})$；
\item 令 $\mathscr{L}=\mathscr{F}(\beta)$；由 $\beta^n=\alpha\in \mathscr{F}$ 与 $\mathscr{U}_n\subseteq \mathscr{F}$：$\beta$ 的
$\mathscr{F}$-共轭为 $\omega^m\beta\in \mathscr{L}$，故 $\mathscr{L}/\mathscr{F}$ 循环、次数 $d\mid n$；又 $\beta$ 与任一选定的
$\sqrt[n]\alpha$ 相差 $\mathscr{U}_n$ 中元素，故 $\mathscr{L}=\mathscr{F}(\sqrt[n]{\bar\alpha})$；
\item 由第 2 步的双射：$\mathscr{F}(\sqrt[n]{\bar\alpha})=\mathscr{F}(\sqrt[n]{\mathscr{D}_\mathscr{L}})$ 强制
$\langle\bar\alpha\rangle=\mathscr{D}_\mathscr{L}\le\mathscr{D}$（在 $\mathscr{F}^*/\mathscr{F}^{*n}$ 中），即
$\alpha=\gamma^n\prod_r a_r^{e_r}$，$\gamma\in \mathscr{F}$；
\item 令 $\beta_0=\gamma\prod_r\theta_r^{e_r}$：$\beta_0^n=\alpha=\beta^n$，
故 $\beta/\beta_0\in\mathscr{U}_n$；
\end{enumerate}
\end{proof}

引理~\ref{thm:landau} 的单项式结论不需要 $a_r$ 在 $\mathscr{F}^*/\mathscr{F}^{*n}$ 中独立；\par

\begin{proposition}[\textbf{D1 快速通道}]
在定理~\ref{thm:landau} 的设置下，对 $\alpha\in \mathscr{F}$：
\[
\sqrt[n]{\alpha}\in\mathscr{K}\iff\bar\alpha\in\mathscr{D}
\iff\exists\,\gamma\in \mathscr{F},\ e_r:\ \alpha=\gamma^n\prod_r a_r^{e_r},
\]
且命中时 $\sqrt[n]{\alpha}=\omega\,\gamma\prod_r\theta_r^{e_r}$ 是\textbf{单项式}；
搜索空间为 $(\mathbb{Z}/n)^m$，可由定理~\ref{thm:necessary} 的赋值条件剪枝；
若 $\alpha$ 位于更低层，判定逐层下推；
\end{proposition}

多个根式相加时，外层内的项只来自单根式交叉与两两乘积交叉：设
$\alpha=c_0+\sum_jc_j\sqrt{d_j}$（$n=2$），猜 $\beta=q+\sum_j s_j\sqrt{b_j}$（$s_j=\pm1$），
则 $\beta^2$ 的根式项恰为 $2q\sqrt b$ 与 $2s's''\sqrt{bb'}$；用根式间的乘除关系
（如 $\sqrt{15}=\sqrt6\sqrt{10}/2$）把项分类为「单」与「对」，由「对」得
$bb'=(c'/2)^2d'$，由「单」得 $2q\sqrt b=c'\sqrt{d'}$，解两两乘积组
（三个未知时 $b_1^2=P_{12}P_{13}/P_{23}$ 等，$b$ 须为正有理数），
再校验 $q^2=c_0-\sum b\in\mathbb{Q}^2$ 与各项符号；
若诸 $\sqrt b$ 全在 $\mathscr{K}$ 内则为 D1，否则为 D2（按引理~\ref{thm:landau}
可化为单项式 $\gamma\sqrt c$，$\gamma\in\mathscr{K}$，$c\in\mathbb{Q}$）；
两项情形即退化为定理~\ref{thm:two-level}；最后必须完整展开 $\beta^2$ 验证。\par

\subsubsection{在 $\mathscr{K}(\omega_n)$ 中检验的危险与挽救}

\begin{proposition}[$\omega$-陷阱]
\label{prop:trap}
在扩域 $\mathscr{K}(\omega_n)$ 中检验幂性会产生假阳性：
\[
16=(1+\ii)^8\in\mathbb{Q}(\omega_8)^{*8},\qquad\text{但 }16\notin\mathbb{Q}^{*8}
\]
（主实根 $2^{1/4}\notin\mathbb{Q}$）；同理
$-4=(1+\ii)^4\in\mathbb{Q}(\ii)^{*4}$，而 $-4\notin\mathbb{Q}^{*4}$；
即「在 $\mathscr{K}(\omega_n)$ 中是 $n$ 次幂」\textbf{不蕴含}「在 $\mathscr{K}$ 中是 $n$ 次幂」；
\end{proposition}

\begin{proposition}[挽救条件]
\label{prop:rescue}
若 $\gcd(n,\varphi(n))=1$（即 $n=2$ 或 $n$ 为素数），则
\[
\alpha\in\mathscr{K}(\omega_n)^{*n}\implies\alpha\in\mathscr{K}^{*n}.
\]
\end{proposition}

\begin{proof}
设 $\beta\in\mathscr{K}(\omega_n)$，$\beta^n=\alpha\in\mathscr{K}$；
令 $d=[\mathscr{K}(\beta):\mathscr{K}]$；因 $\mathscr{K}(\beta)\subseteq\mathscr{K}(\omega_n)$
且 $\mathbb{Q}(\omega_n)/\mathbb{Q}$ 是 Galois 扩张，
\[
d\ \big|\ [\mathscr{K}(\omega_n):\mathscr{K}]
=\frac{\varphi(n)}{[\mathscr{K}\cap\mathbb{Q}(\omega_n):\mathbb{Q}]}\ \big|\ \varphi(n),
\]
故 $\gcd(n,d)=1$；令 $\kappa=N_{\mathscr{K}(\beta)/\mathscr{K}}(\beta)\in\mathscr{K}$：
$\kappa^n=N(\beta^n)=\alpha^d$；取 $r,s\in\mathbb{Z}$ 使 $rn+sd=1$，则
\[
\alpha=\alpha^{rn}\cdot\alpha^{sd}=(\alpha^r)^n(\kappa^n)^s=(\alpha^r\kappa^s)^n\in\mathscr{K}^{*n}.
\]
\end{proof}

因此对素数 $n$，「在 $\mathscr{K}(\omega_n)$ 中检验」是合法的 D1 判据；
对复合 $n$（如 $4,8,9$）必须回到域 $\mathscr{K}$ 本身（定理~\ref{thm:complete}），
$\omega_n$ 只能作为\textbf{构造辅助}使用，最终答案须在 $\mathscr{K}$ 中精确复核。
\subsection{完备判定：主定理与主算法}
\label{sec:complete}

\begin{theorem}[D1 的完备判定]
\label{thm:complete}
设 $\mathscr{K}$ 为数域，$\alpha\in\mathscr{K}^*$，$n\ge2$；则 $\alpha\in\mathscr{K}^{*n}$
当且仅当：
\begin{enumerate}
\item[(a)] $v_{\mathfrak{p}}(\alpha)\equiv0\pmod n$ 对所有素理想 $\mathfrak{p}$；
\item[(b)] 分式理想
$\displaystyle\mathfrak{b}:=\prod_{\mathfrak{p}}\mathfrak{p}^{\,v_{\mathfrak{p}}(\alpha)/n}$
是主理想，取生成元 $\gamma$ 使 $(\gamma)=\mathfrak{b}$；
\item[(c)] 单位 $u:=\gamma^n/\alpha\in\mathscr{O}_{\mathscr{K}}^*$ 是单位群中的 $n$ 次幂：
$u=w^n$，$w\in\mathscr{O}_{\mathscr{K}}^*$；
\end{enumerate}
此时 $\beta:=\gamma/w$ 满足 $\beta^n=\alpha$；
\end{theorem}

\begin{proof}
($\Leftarrow$) $\beta^n=\gamma^n/w^n=\gamma^n/u=\alpha$；
($\Rightarrow$) 设 $\beta^n=\alpha$：则 $(\beta)^n=(\alpha)$，逐素理想给出
$v_{\mathfrak{p}}(\alpha)=n\,v_{\mathfrak{p}}(\beta)$，故 (a) 成立且 $\mathfrak{b}=(\beta)$ 是主理想；
取 $\gamma=\beta$ 则 $u=1=1^n$，(b)(c) 成立；
\end{proof}

\begin{proposition}[可判定性与复杂度]
\label{prop:decidable}
D1 可有效判定：(a) 需要理想分解（有效）；(b) 需要类群结构与主理想检验
（类群有限交换、可计算，Buchmann 型亚指数算法）；(c) 需要 Dirichlet 单位定理的结构
$\mathscr{O}_{\mathscr{K}}^*\cong\mathscr{U}_{\mathscr{K}}\times\mathbb{Z}^r$（基本单位组与调节子可有效计算，
$u$ 的坐标由对数嵌入的格中线性代数读出后精确复核）；一般复杂度受类群与单位群计算支配
（亚指数）；对固定呈现形状的 Kummer 层，Landau 给出了多项式时间算法 [Landau 1992]；
实用实现不应自研 (b)(c)，而应调用 Pari/GP 的 \texttt{bnfinit}、\texttt{bnfisprincipal}、
\texttt{bnfisunit} 或 SageMath/Magma 的等价物。
\end{proposition}

\subsubsection{D2 的完备判据}

\begin{theorem}[D2 判据]
\label{thm:d2}
设内层域 $\mathscr{K}$ 的表达深度为 $t\ge1$，$\alpha\in\mathscr{K}^*$，$n\ge2$；
则 $\sqrt[n]{\alpha}$ 可用深度 $\le t$ 的根式表示当且仅当
\[
\exists\,c\in\mathbb{Q}^*,\ \gamma\in\mathscr{K}(\omega_n)\quad\text{使}\quad
\alpha=\gamma^n c .
\]
\end{theorem}

\begin{proof}
($\Leftarrow$) 取 $c$ 为新根式的底：$\beta:=\gamma\cdot\sqrt[n]{c}$；则 $\beta^n=\alpha$，
且 $\beta\in\mathscr{K}(\omega_n,\sqrt[n]{c})$：$\omega_n=\sqrt[n]{1}$ 与 $\sqrt[n]{c}$
都是底为有理数的一次根式（深度 $1\le t$），$\gamma\in\mathscr{K}(\omega_n)$ 深度 $\le t$，
故 $\beta$ 深度 $\le t$；

($\Rightarrow$) 设 $\beta\in\mathscr{K}(\omega_n,\sqrt[n]{b_1},\dots,\sqrt[n]{b_r})$
（$b_j\in\mathbb{Q}^*$）且 $\beta^n=\alpha\in\mathscr{K}\subseteq\mathscr{K}(\omega_n)$；
在定理~\ref{thm:landau} 中取底域 $\mathscr{F}=\mathscr{K}(\omega_n)$（此时 $\mathscr{U}_n\subseteq \mathscr{F}$）、
生成元 $\theta_j=\sqrt[n]{b_j}$（$\theta_j^n=b_j\in\mathbb{Q}\subseteq \mathscr{F}$），
得 $\beta=\omega\gamma\prod_r\theta_r^{f_r}$，$\gamma\in\mathscr{K}(\omega_n)$；
两边取 $n$ 次幂：$\alpha=\gamma^n\prod_jb_j^{f_j}=\gamma^n c$，
其中 $c:=\prod_jb_j^{f_j}\in\mathbb{Q}^*$；
\end{proof}

\begin{proposition}[\textbf{D2 判据的两个特例}]
$n=2$ 时 $\omega_2=-1\in\mathbb{Q}$，D2 判据化为
$\exists\,c\in\mathbb{Q}^*:\ \alpha/c\in\mathscr{K}^{*2}$；
$\gcd(n,\varphi(n))=1$ 时由命题~\ref{prop:rescue} 化为
$\exists\,c:\ \alpha/c\in\mathscr{K}^{*n}$；
一般 $n$ 的完备判定：枚举候选 $c$，对每个 $\alpha/c$ 在域 $\mathscr{K}(\omega_n)$ 上
执行定理~\ref{thm:complete}。
\end{proposition}

\subsubsection{候选 $c$ 的枚举}

设 $\alpha=\gamma^n c$；逐素理想读出两个约束；\par

\begin{enumerate}
\item \textbf{赋值约束}：$v_{\mathfrak{p}}(\alpha)=n\,v_{\mathfrak{p}}(\gamma)+e(\mathfrak{p}/p)\,v_p(c)$
（$p=\mathfrak{p}\cap\mathbb{Q}$），即
\[
e(\mathfrak{p}/p)\,v_p(c)\equiv v_{\mathfrak{p}}(\alpha)\pmod n
\qquad\text{对每个 }\mathfrak{p};
\]
\item \textbf{范数约束}：$N(\alpha)=N(\gamma)^n\,c^{[\mathscr{K}:\mathbb{Q}]}$，即
\[
e_p\bigl(|N(\alpha)|\bigr)-[\mathscr{K}:\mathbb{Q}]\,v_p(c)\equiv0\pmod n .
\]
\end{enumerate}

当 $\gcd([\mathscr{K}:\mathbb{Q}],n)=1$ 时，范数约束\textbf{单独}确定 $c$ 的 $n$ 次幂剩余类：
\[
\boxed{\ c\ \cong\ \prod_p p^{\,e_p(|N(\alpha)|)\,[\mathscr{K}:\mathbb{Q}]^{-1}\bmod n}
\pmod{\mathbb{Q}^{*n}}\ }
\]
特别地 $n=2$ 且 $[\mathscr{K}:\mathbb{Q}]$ 为奇数时，$c$ 的无平方因子部分恰为
$|N(\alpha)|$ 的无平方因子部分；例~\ref{ex:seventh}：
$|N|=126=2\cdot3^2\cdot7$、$[\mathscr{K}:\mathbb{Q}]=7$，故 $c\cong2\cdot7=14$；
例~\ref{ex:two-sqrt3}：$[\mathscr{K}:\mathbb{Q}]=2$ 与 $n=2$ 不互素，范数不约束（$N=1$），
$c=2$ 由赋值与单位结构给出；若 $n\mid[\mathscr{K}:\mathbb{Q}]$，则 $c^{[\mathscr{K}:\mathbb{Q}]}$
自动是 $n$ 次幂，范数完全失效，此时只能试小底 $c=1,p^{\pm1},p^{\pm2},\dots$；\par

检验 $\alpha/c\in\mathscr{K}^{*n}$ 的系统化办法：令 $\delta=\alpha/c$，取其极小多项式 $R$，
构造 $R(t^n)=N_{\mathscr{K}/\mathbb{Q}}(t^n-\delta)$；当 $n$ 为素数（Capelli）时，
$R(t^n)$ 在 $\mathbb{Q}$ 上不可约 $\implies\delta\notin\mathscr{K}^{*n}$；
可约 $\implies\delta\in\mathscr{K}^{*n}$，且 $\gamma$ 可从因子里恢复
（数值共轭值加范德蒙德或逆 DFT，再做有理识别）。\par

\subsubsection{候选 $c$ 与生成元的相依}

判据 $\bar\alpha\in\langle\bar a_1,\dots,\bar a_m\rangle$ 是在\textbf{子群}中搜索，
允许 $a_r$ 相依；例：$\sqrt{7+2\sqrt{10}}$，最小域 $\mathbb{Q}(\sqrt{10})$，$n=2$，
$7+2\sqrt{10}\notin\mathbb{Q}(\sqrt{10})^2$；但 $\bar{10}=\bar2\cdot\bar5$ 相依，使单项式
\[
\frac{2+\sqrt{10}}{\sqrt2}=\sqrt2+\sqrt5
\]
成为解：$\bigl(\frac{2+\sqrt{10}}{\sqrt2}\bigr)^2=7+2\sqrt{10}$，
即 $c=1/2$ 与 $\gamma=2+\sqrt{10}\in\mathbb{Q}(\sqrt{10})$；
若把生成元集错当独立处理（只试各 $a_r$ 与它们的乘积），会漏掉这个解。
\subsection{六步手算法}
\label{sec:manual}

\subsubsection{流程总览}

\begin{center}
\begin{tikzpicture}[node distance=4.6mm and 10mm,>=stealth,
  box/.style={draw,rounded corners,align=center,inner sep=2.4pt,font=\small},
  lab/.style={font=\scriptsize,inner sep=1.5pt}]
\node[box] (s0) {第 0 步 定域定向};
\node[box,right=of s0] (s1) {第 1 步 范数快筛};
\node[box,right=of s1] (s2) {第 2 步 两层特判（大公式）};
\node[box,below=11mm of s1] (s3) {第 3 步 D1 正面搜索（单项式枚举 / 系数比较）};
\node[box,below=of s3] (s4) {第 4 步 D2 搜索（定候选 $c$，检验 $\alpha/c$ 的幂性）};
\node[box,below=of s4] (s5) {第 5 步 进阶否证（赋值 / 单位 / 次数）};
\node[box,below=of s5] (s6) {第 6 步 精确验证（必做）$\to$ 输出 $\beta=\gamma\sqrt[n]{c}$};
% 主干：第 0 → 第 1 → 第 2 步
\draw[->] (s0)--(s1);
\draw[->] (s1)--node[lab,above]{通过}(s2);
% 第 1 步「否」：从上方向右的弧线，不压任何方框
\draw[->] (s1.north) to[out=100,in=80] node[lab,above]{否：D1 不可去，停}(s2.north);
% 主干：第 3 → 第 4 → 第 5 → 第 6 步
\draw[->] (s3)--(s4);
\draw[->] (s4)--(s5);
\draw[->] (s5)--(s6);
% 第 2 步「未覆盖」：先下折，再横向左行，最后自上方进入第 3 步
\draw[->] (s2.south) -- ++(0,-5.5mm) -| (s3.north);
\node[lab,above] at ($(s2.south |- s3.north)+(0,5.5mm)$) [xshift=-14mm] {未覆盖};
% 第 2 步「命中」：右侧独立走廊，下行后自右方进入第 6 步
\draw[->] (s2.east) -- ++(13mm,0) |- (s6.east);
\node[lab,right] at ($(s2.east)+(13mm,0)+(0,-9mm)$) {命中：验证，停};
\end{tikzpicture}
\end{center}

\subsubsection{第 0 步：定域与定向}

把 $\alpha$ 写成内层生成元的规范形 $\alpha=\sum_Sc_S\prod_r\theta_r^{a_r}$
（$0\le a_r<e_r$，用折叠规则降幂）；合并到最小子域，找出 $\alpha$ 真正属于的最小域——
判定永远在最小域上做，这一步直接决定后面用哪条公式；逐层数出 $[\mathscr{K}:\mathbb{Q}]$，
每加一个生成元前先确认它的底\textbf{不是}当前底域的幂（塌缩防护，见
第~\ref{sec:manual} 节末 T6）；$n$ 为偶数时确认 $\alpha$ 的所有实影子 $\ge0$；
出口：最小域 $\mathscr{K}$、生成元表、$[\mathscr{K}:\mathbb{Q}]$、$n$。\par

\subsubsection{第 1 步：范数快筛}

逐层收缩计算 $N(\alpha)$（工具 T1），直到得到整数；$|N(\alpha)|$ 必须恰好是有理数的
$n$ 次幂：否则 D1 不可去，停，这就是证书；通过则把 $|N(\alpha)|$ 的素因子分解留好，
第 4 步定候选 $c$ 要用；手算要点是先折叠成 $\alpha=A+B\theta$ 型（$\deg\le1$）再用 T1 的
封闭式；多生成元逐层做（如 $\mathbb{Q}(\sqrt3,x)$ 先对 $\sqrt3$ 共轭配对相乘，再对 $x$ 收缩）。\par

\subsubsection{第 2 步：两层特判}

经第 0 步合并后 $\alpha=a_0+a_1\sqrt a$（$a$ 无平方因子，$a_0,a_1\in\mathbb{Q}$）时，
算 $s^2=a_0^2-a_1^2a$ 并用定理~\ref{thm:two-level}：D1 需再过两道平方性检查，D2 由公式
直接给出全部答案，两个新根式的底 $\tfrac{a_0\pm s}2\in\mathbb{Q}$ 深度为 $1$，嵌套消失。\par

\subsubsection{第 3 步：D1 正面搜索}

\textbf{情形 A：$\alpha$ 落在某层底域 $\mathscr{F}$ 中}，该层生成元 $\theta_r$ 满足
$\theta_r^n=a_r\in \mathscr{F}$：枚举 $e=(e_1,\dots,e_m)\in\{0,\dots,n-1\}^m$（先试 $e=0$，
用第 1 步的 $|N(\alpha)|$ 分解剪枝）；检验 $\alpha\prod_ra_r^{-e_r}\in \mathscr{F}^{*n}$
（$\mathscr{F}=\mathbb{Q}$ 直接开方验算，$\mathscr{F}$ 是二次域用 T2，更深则递归）；命中则取候选
$\beta=\gamma\prod_r\theta_r^{e_r}$，转第 6 步；\par

\textbf{情形 B：$\alpha$ 真含顶层生成元}（$\alpha\notin$ 底域）：设
$\beta=\sum_Sb_S\theta^S$（系数 $b_S$ 取自底域），展开 $\beta^n$、折叠、与 $\alpha$
逐系数相等得系数方程组，解之；无解则本层 D1 否；$n=2$、单生成元时 T2 就是这个方程组的解。\par

\subsubsection{第 4 步：D2 搜索}

\textbf{4.1 定候选 $c\in\mathbb{Q}^*$}：用第 1 步存下的 $|N(\alpha)|=\prod_pp^{f_p}$；
若 $\gcd([\mathscr{K}:\mathbb{Q}],n)=1$，用上文同余式
$c\cong\prod_pp^{f_p[\mathscr{K}:\mathbb{Q}]^{-1}\bmod n}\pmod{\mathbb{Q}^{*n}}$；
特别地 $n=2$、$[\mathscr{K}:\mathbb{Q}]$ 奇时 $c=\operatorname{sqfree}(|N(\alpha)|)$；
若 $n\mid[\mathscr{K}:\mathbb{Q}]$，范数不约束，试小底 $c=1,p^{\pm1},p^{\pm2},\dots$，
逐个进 4.2；\par

\textbf{4.2 检验 $\alpha/c\in\mathscr{K}^{*n}$}：用第 3 步情形 A 的手段；机械替代是
$R(t^n)$ 的可约性判据与因式恢复；恢复 $\gamma$ 的通用配方：数值算出所有影子的目标值，
解范德蒙德方程组读出系数，有理逼近，最后精确展开验证；\par

\textbf{4.3 命中} $\gamma^n=\alpha/c$ 则结果 $\beta=\gamma\sqrt[n]{c}$，新根式的底
$c\in\mathbb{Q}$ 深度为 $1$。\par

\subsubsection{第 5 步：进阶否证}

\begin{enumerate}
\item \textbf{素理想赋值}：取素数 $p$（优先 $|N(\alpha)|$ 的素因子），求 $\theta$ 的极小多项式
$\bmod p$ 的根 $r$，得 $\mathfrak{p}=(p,\theta-r)$；若 $\alpha\equiv0\pmod{\mathfrak{p}}$，
用泰勒展开定重数 $v_{\mathfrak{p}}(\alpha)$；$n\nmid v_{\mathfrak{p}}(\alpha)$ 则不可去，停；
对 D2 还要做\textbf{交叉一致性}：把
$v_{\mathfrak{p}}(\alpha)=n\,v_p(\gamma)+e(\mathfrak{p}/p)v_p(c)$ 与第 4 步的范数约束联立，
同余矛盾则连 D2 也没有；
\item \textbf{单位论证}（秩 $1$ 的域，如纯三次域）：$\alpha$ 是单位时写成
$\alpha=\pm\varepsilon^m$（$\varepsilon$ 的基本性用 Minkowski 盒穷举证），
$n\nmid m$ 则不可去；
\item \textbf{次数检验}：4.2 的 $R(t^n)$ 不可约性，反着用作否证证书。
\end{enumerate}

\subsubsection{第 6 步：精确验证（必做）}

把 $\beta^n$ 在域里完整展开（逐层用 $\theta_r^{e_r}=a_r$ 折叠），逐系数确认等于 $\alpha$；
再数值抽查主分支以固定 $\pm$ 与 $\omega_n^j$ 分支（$n$ 个根相差 $\omega_n^j$ 倍）；
对不上则查错或换分支；这一步算数，其余各步都可能错。\par

\subsubsection{工具箱}

\textbf{T1（范数收缩）} $\theta^{n'}=a$，$\alpha=A+B\theta$（先折叠到 $\deg\le1$）：
\[
N_{\mathscr{F}(\theta)/\mathscr{F}}(\alpha)=A^{n'}+(-1)^{n'+1}aB^{n'}.
\]
特例：$n'=2$ 时 $A^2-aB^2$，$n'=3$ 时 $A^3+aB^3$；$\deg2$、$n'=3$ 用
$\omega$-恒等式（见上文）；多生成元逐层收缩，最后一步到 $\mathbb{Q}$ 可用结式；\par

\textbf{T2（二次域系数比较）} $(u+v\sqrt d)^2=\alpha_0+\alpha_1\sqrt d$：
\[
2uv=\alpha_1,\qquad v^2=\frac{\alpha_0\pm s}{2d},\quad s^2=\alpha_0^2-d\alpha_1^2 .
\]
$u,v\in\mathbb{Q}$ 有解当且仅当 $s\in\mathbb{Q}$ 且
$\tfrac{\alpha_0+s}2,\tfrac{\alpha_0-s}{2d}\in\mathbb{Q}^2$——
这就是第 2 步 D1 判据的来历；\par

\textbf{T3（折叠）} $\theta_r^{q}=a_r^{\lfloor q/e_r\rfloor}\theta_r^{\,q\bmod e_r}$；\par

\textbf{T4（候选 $c$ 速查）} $n=2$ 且 $[\mathscr{K}:\mathbb{Q}]$ 奇时
$c=\operatorname{sqfree}(|N(\alpha)|)$；一般用 4.1 的同余式；\par

\textbf{T5（有理识别）} 高精度数值 $\to$ 连分数或 \texttt{limit\_denominator}
$\to$ 猜有理数 $\to$ 精确展开验证；只用于「找」，不用于「判」；\par

\textbf{T6（塌缩检查）} 每加一层前验证底不是当前底域的幂：
$5+2\sqrt6\notin\mathbb{Q}(\sqrt6)^2$，但 $(\sqrt2+\sqrt3)^2=5+2\sqrt6$；
同理，向底域添加 $\omega_n$ 或更深的根式可能使某层根式\textbf{塌缩}（其底成为当前底域的幂），
因此层化算法在每一层必须用\textbf{当前}底域重新检验底的幂性，不能只在最初呈现上判定一次；\par

\textbf{T7（多项配平法）} 见上文多项根式相加的配平步骤；注意 $[\mathscr{K}:\mathbb{Q}]$ 与 $n$
不互素时范数读不出 $c$，T7 是此时的主力；两项情形退化为 T2 大公式。\par

\subsubsection{一页速查表}

\begin{center}
\begin{tabularx}{\textwidth}{@{}lX@{}}
\toprule
看到的形状 & 直接走 \\
\midrule
$\sqrt{a_0+a_1\sqrt a}$ & 第 0 步合并 $\to$ 第 2 步大公式（一步出 D1 判定与 D2 答案） \\
外层内多个根号相加 $\sqrt{c_0+\sum_jc_j\sqrt{d_j}}$ & 第 1 步范数 $\to$ T7 多项配平 \\
$\sqrt[n]{\text{多项式}(x)}$，$x^m=a\in\mathbb{Q}$ & 第 1 步范数（T1）$\to$ 否则第 3、4 步 \\
范数是 $n$ 次幂但仍可疑 & 第 5 步：赋值、单位、次数 \\
想要漂亮答案 & 第 4 步：$c=\operatorname{sqfree}(|N|)$（$n=2$、$[\mathscr{K}:\mathbb{Q}]$ 奇）；不互素时用 T7 试探 \\
任何候选 & 第 6 步展开验证，不过就重来 \\
\bottomrule
\end{tabularx}
\end{center}

当 $[\mathscr{K}:\mathbb{Q}]\ge5$，或生成元不少于 $2$ 个而第 2 步不适用时，系数方程组与
$R(t^n)$ 的分解会迅速超出手算规模：此时改用「数值共轭值加逆 DFT 加有理识别加精确展开」
的机械流程，或调用 Pari/GP 的 \texttt{bnfinit} 体系（命题~\ref{prop:decidable}）；
完备判定的 (b)(c) 两步本质上是类群与单位群计算，不要手算。
\subsection{六个算例的完整裁决}
\label{sec:examples}

以下每条结论都配有可复核的证书：既有域内的精确等式，也有赋值或范数的否证证书。\par

\subsubsection{$\sqrt{3+2\sqrt2}$：D1 成立}
\label{ex:first}

\emph{第 0 步}：$\alpha=3+2\sqrt2$ 的最小域 $\mathscr{K}=\mathbb{Q}(\sqrt2)$，$n=2$；
\emph{第 1 步}：$N=3^2-2\cdot2^2=1$；
\emph{第 2 步}：$s^2=9-8=1$，取 $s=-1$：$\tfrac{3-1}2=1\in\mathbb{Q}^2$，
$\tfrac{3+1}{2\cdot2}=1\in\mathbb{Q}^2$，故 D1 通过，$u=v=1$，$2uv=2$ 与 $a_1=2$ 相符；
\emph{第 6 步}：$(1+\sqrt2)^2=1+2+2\sqrt2=3+2\sqrt2$。故
\[
\sqrt{3+2\sqrt2}=1+\sqrt2 .
\]

\subsubsection{$\sqrt{2-\sqrt3}$：D1 失败，D2 成功}
\label{ex:two-sqrt3}

第 0 步：$\mathscr{K}=\mathbb{Q}(\sqrt3)$，$n=2$；第 1 步：$N=4-3=1$（通过，但无信息）；
第 2 步：$s^2=4-3=1$；D1 检查：$s=1$ 时 $\tfrac{2+1}2=\tfrac32\notin\mathbb{Q}^2$；
$s=-1$ 时 $\tfrac{2-1}2=\tfrac12\notin\mathbb{Q}^2$，故 \textbf{D1 失败}；
等价地：$2-\sqrt3=(2+\sqrt3)^{-1}$ 是基本单位 $2+\sqrt3$ 的 $-1$ 次幂，
指数不被 $2$ 整除，因而不可能是平方；\par

D2 构造：$s=1$，取负号使交叉项为 $-\sqrt3$，
\[
\sqrt{2-\sqrt3}=\sqrt{\tfrac32}-\sqrt{\tfrac12}=\frac{\sqrt6-\sqrt2}{2},
\]
即 $c=2$：$\bigl(\tfrac{\sqrt3-1}{2}\bigr)^2=\tfrac{2-\sqrt3}{2}$，而
$\sqrt2\cdot\tfrac{\sqrt3-1}2=\tfrac{\sqrt6-\sqrt2}2$；
第 6 步：$\bigl(\sqrt{\tfrac32}-\sqrt{\tfrac12}\bigr)^2
  =\tfrac32+\tfrac12-2\sqrt{\tfrac34}=2-\sqrt3$。\par

\subsubsection{$\sqrt{2-2^{1/7}}$：D1 失败，D2 成功}
\label{ex:seventh}

第 0 步：$x=2^{1/7}$，$\mathscr{K}=\mathbb{Q}(x)$，$[\mathscr{K}:\mathbb{Q}]=7$
（$t^7-2$ 在 $2$ 处 Eisenstein）；
第 1 步（定理~\ref{thm:norm-two-term}，$m=7$，$A=2$，$B=-1$，$a=2$）：
\[
N(2-x)=2^7+(-1)^8\cdot2\cdot(-1)^7=128-2=126=2\cdot3^2\cdot7 .
\]
指数 $1$ 与 $1$ 不被 $2$ 整除，$126$ 不是平方，故 \textbf{D1 失败}；
又 $[\mathscr{K}:\mathbb{Q}]=7$ 是奇数，$\mathscr{K}$ 无二次子域，$\sqrt{14}\notin\mathscr{K}$；\par

第 4.1 步：$n=2$ 与 $[\mathscr{K}:\mathbb{Q}]=7$ 互素，$c\cong\operatorname{sqfree}(126)=2\cdot7=14$；
第 4.2 步：判定 $\delta=(2-x)/14\in\mathscr{K}^{*2}$ 归结为 $14$ 次多项式
$P(t)=R(t^2)$（$R$ 为 $\delta$ 的极小多项式）在 $\mathbb{Q}$ 上的因式分解
（Capelli，$n=2$：$R(t^2)$ 可约当且仅当 $\delta\in\mathscr{K}^2$）；
机器分解给出型 $[7,7]$，再由 $P$ 的七次因子作 $60$ 位以上数值逆 DFT 恢复系数、
逐级有理逼近（分母上界 $10^6$ 内即命中，实际分母只有 $7,14$），最后在
$\mathbb{Q}[t]/(t^7-2)$ 中逐系数精确验证
\[
\gamma=\frac{1+x^2-x^3-\tfrac12x^4+x^5}{7}\in\mathbb{Q}(x),
\qquad \gamma^2=\frac{2-x}{14}.
\]

手算复核：记 $s=1+x^2-x^3-\tfrac12x^4+x^5$，逐对展开 $s^2$——常数项 $1+4+2=7$，
$x$ 项 $\tfrac12x-4x=-\tfrac72x$，$x^2$ 至 $x^6$ 项全部相消，即
$s^2=7-\tfrac72x=\tfrac72(2-x)$，故 $\gamma^2=s^2/49=(2-x)/14$；\par

第 6 步：完整卷积表（$x^7=2$；$a_0=1$，$a_2=1$，$a_3=-1$，$a_4=-\tfrac12$，$a_5=1$，
其余为 $0$；$c_k=\sum_{r+t=k}a_ra_t$）：\par

\begin{center}
\begin{tabular}{@{}cccc@{}}
\toprule
$k$ & $c_k$ & 计算明细 & 折叠后 \\
\midrule
$0$ & $1$ & $1\cdot1$ & 常数 $+1$ \\
$1$ & $0$ & 无 $x$ 项 & --- \\
$2$ & $2$ & $2a_0a_2=2$ & $x^2+2$ \\
$3$ & $-2$ & $2a_0a_3=-2$ & $x^3-2$ \\
$4$ & $0$ & $2a_0a_4+a_2^2=-1+1$ & $0$ \\
$5$ & $0$ & $2a_0a_5+2a_2a_3=2-2$ & $0$ \\
$6$ & $0$ & $2a_2a_4+a_3^2=-1+1$ & $0$ \\
$7$ & $3$ & $2a_2a_5+2a_3a_4=2+1$ & $3x^7=6$，常数 $+6$ \\
$8$ & $-\tfrac74$ & $2a_3a_5+a_4^2=-2+\tfrac14$ & $-\tfrac74x^8=-\tfrac72x$ \\
$9$ & $-1$ & $2a_4a_5=-1$ & $-x^9=-2x^2$ \\
$10$ & $1$ & $a_5^2=1$ & $x^{10}=2x^3$ \\
\bottomrule
\end{tabular}
\end{center}

收集各项：常数 $1+6=7$；$x$ 项 $-\tfrac72$；$x^2$ 项 $2-2=0$；$x^3$ 项 $-2+2=0$；
$x^4,x^5,x^6$ 项为 $0$；故 $s^2=\tfrac72(2-x)$，$\gamma^2=s^2/49=(2-x)/14$；于是
\[
\boxed{\ \sqrt{2-2^{1/7}}\ =\ \frac{\sqrt{14}}{7}
\left(1+2^{2/7}-2^{3/7}-\frac{2^{4/7}}{2}+2^{5/7}\right)\ }
\]
深度由 $2$ 降为 $1$；此处「记忆中 $\sqrt{2-2^{1/7}}$ 可以化掉外层根号」是\textbf{对的，
但只在 D2 意义下}：D1 判定失败（$N=126$ 非平方），D2 判定成功，
新根号 $\sqrt{14}$ 恰是 $|N(\alpha)|=126=2\cdot3^2\cdot7$ 的无平方因子部分。\par

\subsubsection{$\sqrt[3]{2-\sqrt3+2\cdot2^{1/5}}$：D1 与 D2 双双失败}
\label{ex:double-no-1}

第 0 步：$x=2^{1/5}$，$\mathscr{K}=\mathbb{Q}(\sqrt3,x)$，
$[\mathscr{K}:\mathbb{Q}]=10$（次数 $2$ 与 $5$ 互素，塔式相乘），$n=3$，
$\alpha=2-\sqrt3+2x=2(1+x)-\sqrt3$；\par

\emph{第 1 步（范数，两层收缩）}；对 $\sqrt3$（$\alpha=(2+2x)+(-1)\sqrt3$）：
$N_{\mathscr{K}/\mathbb{Q}(x)}(\alpha)=(2+2x)^2-3=4x^2+8x+1$；
对 $x$：
\[
N(\alpha)=\operatorname{Res}(t^5-2,\ 4t^2+8t+1)
=4^5\bigl(s_+^5-2\bigr)\bigl(s_-^5-2\bigr),
\]
其中 $s_\pm=(-2\pm\sqrt3)/2$ 是 $4t^2+8t+1$ 的两根；幂和递推
（$s_++s_-=-2$，$s_+s_-=\tfrac14$）：
$P_0=2$，$P_1=-2$，$P_2=\tfrac72$，$P_3=-\tfrac{13}2$，$P_4=\tfrac{97}8$，$P_5=-\tfrac{181}8$，于是
\[
(s_+^5-2)(s_-^5-2)=\tfrac1{1024}-2\bigl(-\tfrac{181}8\bigr)+4=\tfrac{50433}{1024},
\qquad
N(\alpha)=4^5\cdot\frac{50433}{1024}=50433=3\cdot16811 ,
\]
其中 $16811$ 为素数；$v_3=1$ 不被 $3$ 整除，故 \textbf{D1 失败}；\par

\emph{第 5 步（D2 否证）}；第 4.1 步先给出候选：$\gcd(10,3)=1$，
$c\cong3^{1}\cdot16811^{1}\pmod{\mathbb{Q}^{*3}}$，即 $c\cong50433$；取
$t^5-2\equiv(t+1)\mu_{10}\pmod3$，其中 $\mu_{10}=t^4-t^3+t^2-t+1$ 在
$\mathbb{F}_3$ 上不可约，故 $\mathfrak{q}_1=(3,\sqrt3,x+1)$ 是 $f=1$ 的素理想，
且 $e(\mathfrak{q}_1/3)=2$（$\mathfrak{q}_1$ 在 $\mathfrak{p}=(\sqrt3)$ 上非分歧）；
$\alpha=2(1+x)-\sqrt3\in\mathfrak{q}_1$（两项都落入），而
$v_3\bigl(N(\alpha)\bigr)=\sum f\,v_{\mathfrak{q}}(\alpha)=1$ 强制
$v_{\mathfrak{q}_1}(\alpha)=1$ 恰好一次；\par

若 $\gamma\in\mathscr{K}$：D2 需 $\alpha=\gamma^3c$，在 $\mathfrak{q}_1$ 上
$1=3v_{\mathfrak{q}_1}(\gamma)+v_{\mathfrak{q}_1}(c)=3v+2v_3(c)$，
故 $v_3(c)\equiv2\pmod3$；再由范数 $50433=N(\gamma)^3c^{10}$ 取 $v_3$：
$1=3v_3(N\gamma)+10v_3(c)$，故 $v_3(c)\equiv1\pmod3$；\textbf{矛盾}；\par

若 $\gamma\in\mathscr{K}(\omega_3)$：首先 $\omega_3\notin\mathscr{K}$——$\mathscr{K}$ 有实嵌入
$\sqrt3\mapsto\sqrt3$、$x\mapsto2^{1/5}$，其像全实，而 $\omega_3$ 的像只能是虚的
$\omega_3^{\pm1}$，矛盾——故 $[\mathscr{K}(\omega_3):\mathscr{K}]=2$、
$[\mathscr{K}(\omega_3):\mathbb{Q}]=20$，范数方程变为
$N(\alpha)^2=N_{\mathscr{K}(\omega_3)}(\gamma)^3c^{20}$，取 $v_3$ 得
$2=3g+20v_3(c)$，即 $v_3(c)\equiv1\pmod3$；另一方面在 $\mathfrak{q}_1$ 上方的素理想
$\mathfrak{Q}$ 处 $v_{\mathfrak{Q}}(\alpha)=1$，而
$v_{\mathfrak{Q}}(c)=e(\mathfrak{Q}/3)v_3(c)=2v_3(c)$（复合域 $\mathbb{Q}_3(\sqrt3,\ii)$ 在 $3$ 上
$e=2$、$f=2$，故 $e(\mathfrak{Q}/\mathfrak{q}_1)=1$），得
$1=3v+2v_3(c)$，即 $v_3(c)\equiv2\pmod3$，仍\textbf{矛盾}；\par

两种归属都已封死，故\textbf{不存在任何 $c$，D2 失败}；这是完整证明，不是启发式。\par

\subsubsection{$\sqrt{2+\sqrt3+\sqrt[3]{3+\sqrt5}}$：D1 与 D2 双双失败}
\label{ex:double-no-2}

第 0 步：$\theta^3=3+\sqrt5$，$\mathscr{F}_0=\mathbb{Q}(\sqrt3,\sqrt5)$，$[\mathscr{F}_0:\mathbb{Q}]=4$，
$[\mathscr{K}:\mathbb{Q}]=12$；事实上 $3+\sqrt5$ 不是 $\mathscr{F}_0$ 的立方：否则
$N_{\mathscr{F}_0/\mathbb{Q}}(\delta)^3=N(3+\sqrt5)=(3+\sqrt5)(3-\sqrt5)=4$，
而 $N(\delta)^3=4$ 无有理数解；\par

\emph{第 1 步（范数，三层收缩）}；对 $\theta$
（$N_{\mathscr{K}/\mathscr{F}_0}(a_0+\theta)=a_0^3+\theta^3$，因
$\prod_j(A+\omega_3^j\theta)=A^3+\theta^3$）：
\[
N_{\mathscr{K}/\mathscr{F}_0}(\alpha)=(2+\sqrt3)^3+(3+\sqrt5)=(26+15\sqrt3)+(3+\sqrt5)
=29+15\sqrt3+\sqrt5 .
\]
对 $\sqrt5$：$(29+15\sqrt3)^2-5=1516+870\sqrt3-5=1511+870\sqrt3$；
对 $\sqrt3$：
\[
1511^2-3\cdot870^2=2283121-2270700=12421 .
\]
$12421$ 是\textbf{素数}（逐一试除 $2,3,5,\dots,109$，$\sqrt{12421}\approx111.4$，全不尽），
故非平方，\textbf{D1 失败}；\par

\emph{第 5 步（D2 否证）}；$|N(29+15\sqrt3+\sqrt5)|=12421=p$ 为素数，故该元素生成
$\mathscr{F}_0$ 中范数 $p$、$f=1$ 的素理想 $\mathfrak{B}$；$\theta_0=-(2+\sqrt3)$ 是
$t^3-(3+\sqrt5)$ 在 $\mathfrak{B}$ 上的根（因 $29+15\sqrt3+\sqrt5\equiv0$），且为单根
（$3\theta_0^2=3(7+4\sqrt3)$ 的范数为 $9$，与 $p$ 互素），又 $3\mid p-1$ 故 $\omega_3$
落在剩余域中，且 $3+\sqrt5\not\equiv0$（$N(3+\sqrt5)=4$），故 $t^3-(3+\sqrt5)$ 在
$\mathfrak{B}$ 上\textbf{完全分裂}，得三个素理想 $\mathfrak{Q}_1,\mathfrak{Q}_2,\mathfrak{Q}_3$，
$e=f=1$，而 $\alpha=\theta+(2+\sqrt3)=\theta-\theta_0$ 在每个 $\mathfrak{Q}_j$ 上赋值恰为 $1$；
$\sqrt3$-共轭素理想 $\mathfrak{B}'$（$\ne\mathfrak{B}$，否则 $30\sqrt3\in\mathfrak{B}$，
但 $p\nmid N(30\sqrt3)=30^4\cdot9$）同理给出三个；合计六个 $f=1$ 的素理想上
$v(\alpha)=1$ 为奇，D2 需要 $v_p(c)$ 为奇；但 $N(\alpha)=N(\beta)^2c^{12}$ 给出
$e_p\bigl(N(\alpha)\bigr)=1=2g+12v_p(c)$，右边恒为偶——\textbf{矛盾，D2 失败}；
（$[\mathscr{K}:\mathbb{Q}]=12$ 为偶数是关键，与命题~\ref{prop:rescue} 的条件恰成对偶。）\par

\subsubsection{$\sqrt[3]{\sqrt[3]2-1}$（Ramanujan）：D1 失败，D2 成功}
\label{ex:ramanujan}

第 0 步：$x=\sqrt[3]2$，$\mathscr{K}=\mathbb{Q}(x)$，$n=3$，$\alpha=x-1$；
第 1 步：$N(x-1)=x^3-1=1$，$|N|=1$，无信息（$\alpha$ 是单位）；\par

\emph{D1 否证}；$(1+x+x^2)(x-1)=x^3-1=1$，即 $x-1=\varepsilon^{-1}$，
$\varepsilon=1+x+x^2$ 是 $\mathbb{Q}(\sqrt[3]2)$ 的基本单位：在 Minkowski 型界
$|u_1|<\varepsilon<3.85$、$|u_2|<1$（等价地 $|a|<1.95$，$|b|<2.78$，$|c|<2.21$）的盒内
穷举，$(1,\varepsilon)$ 中没有 $|N|=1$ 的单位，故任何单位都是 $\pm\varepsilon^m$；
若 $x-1=\gamma^3$，则 $\gamma$ 是单位（$N(\gamma)^3=-1$），$\gamma=\pm\varepsilon^m$，
于是 $\pm\varepsilon^{3m}=\varepsilon^{-1}$ 迫使 $3m=-1$，无整数解，故 \textbf{D1 失败}；
独立复核：$\gamma$ 的极小多项式 $(t^3+1)^3-2$ 在 $\mathbb{Q}$ 上不可约，
$[\mathbb{Q}(\gamma):\mathbb{Q}]=9>3$；\par

\emph{D2 构造}；$n=3=[\mathscr{K}:\mathbb{Q}]$，范数失效，试小底；钥匙恒等式：
\[
(1-x+x^2)(1+x)=1+x^3=3,\qquad (1+x)^3=3(1+x+x^2)=3\varepsilon,
\]
故
\[
(1-x+x^2)^3=\frac{27}{(1+x)^3}=\frac{27}{3\varepsilon}=9\,\varepsilon^{-1}=9(x-1).
\]
取 $\gamma=1-x+x^2\in\mathscr{K}$、$c=\tfrac19$，则 $\alpha=\gamma^3c$；\par

第 6 步：$\sqrt[3]{x-1}=(1-x+x^2)\sqrt[3]{\tfrac19}$，且
$\frac{(1-x+x^2)^3}9=\frac{9(x-1)}9=x-1$；故
\[
\sqrt[3]{\sqrt[3]2-1}=\frac{1-\sqrt[3]2+\sqrt[3]4}{\sqrt[3]9}
=\sqrt[3]{\tfrac19}-\sqrt[3]{\tfrac29}+\sqrt[3]{\tfrac49},
\]
这正是 Landau 单项式形状 $\gamma\theta^{-1}$（$\theta=\sqrt[3]9$）。
\subsection{复数参与进来时的情形}
\label{sec:complex}

全部理论逐字适用，数域不必实，变化只有三点；\par

\begin{enumerate}
\item \textbf{全正性条件消失}：定理~\ref{thm:necessary}(3) 只对实嵌入有意义；
\item \textbf{$\omega_n$ 部分内置于底域}：$\mathbb{Q}(\ii)\supset\mathscr{U}_4$、
$\mathbb{Q}(\omega_3)\supset\mathscr{U}_3$ 等，Kummer 快速通道（引理~\ref{thm:landau}）的合法化更容易；
\item \textbf{去嵌套的可能性变大}：$\sqrt{2\ii}=1+\ii$（$(1+\ii)^2=2\ii$，D1 于
$\mathbb{Q}(\ii)$）；$\sqrt[4]{-4}=1+\ii$（D1 于 $\mathbb{Q}(\ii)$，而 D1 于 $\mathbb{Q}$ 失败）；
\end{enumerate}

\textbf{反向陷阱}：实 $\alpha>0$ 可以有复去嵌套而无实去嵌套——$\sqrt[8]{16}=1+\ii\in\mathbb{Q}(\ii)$，
而主实根 $2^{1/4}\notin\mathbb{Q}$；这正是命题~\ref{prop:rescue} 的 $\gcd(n,\varphi(n))$
条件在起作用，也再次说明「在 $\mathscr{K}(\omega_n)$ 中是幂」不蕴含「在 $\mathscr{K}$ 中是幂」。
\subsection{补充算例：几个根号相加的情形}

\textbf{例 G} $\sqrt{6+2\sqrt2+2\sqrt3+2\sqrt6}$；
第 0 步：$\sqrt6=\sqrt2\sqrt3$，最小域 $\mathscr{K}=\mathbb{Q}(\sqrt2,\sqrt3)$，
$[\mathscr{K}:\mathbb{Q}]=4$，四个实影子下 radicand 全非负；
第 1 步：对 $\sqrt2$ 取影子得 $(6+2\sqrt3)^2-(2\sqrt2+2\sqrt6)^2=16+8\sqrt3$，
再对 $\sqrt3$ 得 $16-8\sqrt3$，故 $N=256-192=64=8^2$（通过，但不保证可去）；
第 3 步（T7 配平加系数比较）：设 $\beta=q+r\sqrt2+s\sqrt3+t\sqrt6$，则
\[
q^2+2r^2+3s^2+6t^2=6,\quad qr+3st=1,\quad qs+2rt=1,\quad qt+rs=1 .
\]
对称试探 $t=0$：$qr=qs=rs=1$ 给出 $r=s=\pm1$、$q=\pm1$，代回 $1+2+3=6$ 成立；
第 6 步：$(1+\sqrt2+\sqrt3)^2=6+2\sqrt2+2\sqrt3+2\sqrt6$，故
\[
\sqrt{6+2\sqrt2+2\sqrt3+2\sqrt6}=1+\sqrt2+\sqrt3\qquad(\text{D1}).
\]

\textbf{例 H} $\sqrt{10+2\sqrt6+2\sqrt{10}+2\sqrt{15}}$；
第 0 步：$\sqrt{15}=\sqrt6\sqrt{10}/2\in\mathbb{Q}(\sqrt6,\sqrt{10})$，
$[\mathscr{K}:\mathbb{Q}]=4$（$\sqrt6\sqrt{10}=2\sqrt{15}$）；
第 1 步：$(10+2\sqrt{10})^2-(2\sqrt6+2\sqrt{15})^2=56+16\sqrt{10}$，
$56^2-16^2\cdot10=576=24^2$（通过）；注意 $[\mathscr{K}:\mathbb{Q}]=4$ 与 $n=2$ 不互素，
范数读不出候选 $c$，T7 主力登场；
第 3 步（T7 配平）：三项全部来自两两乘积（$q=0$）：$uv=6$，$uw=10$，$vw=15$，
故 $uvw=30$，$(u,v,w)=(2,3,5)$；和校验 $2+3+5=10=c_0$ 成立，得
$\beta=\sqrt2+\sqrt3+\sqrt5$；
D1 与 D2 判定：$\mathbb{Q}(\sqrt6,\sqrt{10})$ 对应平方类
$\{1,6,10,15\}\subseteq\mathbb{Q}^*/{\square}$，而 $2\notin\{1,6,10,15\}$，
故 $\beta\notin\mathscr{K}$，\textbf{D1 失败}（$\mathscr{K}$ 内平方根在 $\pm$ 下唯一，确无解）；
单项式形式 $\beta=\dfrac{2+\sqrt6+\sqrt{10}}{\sqrt2}$，$\gamma=2+\sqrt6+\sqrt{10}\in\mathscr{K}$，
即 D2 判据中的 $c=\tfrac12$：$2\alpha=(2+\sqrt6+\sqrt{10})^2$。
第 6 步：$(\sqrt2+\sqrt3+\sqrt5)^2=10+2\sqrt6+2\sqrt{10}+2\sqrt{15}$，故
\[
\sqrt{10+2\sqrt6+2\sqrt{10}+2\sqrt{15}}=\sqrt2+\sqrt3+\sqrt5
=\frac{2+\sqrt6+\sqrt{10}}{\sqrt2}\qquad(\text{D2}).
\]

\subsection{三个容易被忽略的技术要点}
\label{sec:closing}

\begin{enumerate}
\item 引理~\ref{thm:landau} 的结论\textbf{不需要} $a_r$ 在 $\mathscr{F}^*/\mathscr{F}^{*n}$ 中独立：
判据是在子群 $\mathscr{D}$ 中搜索，允许生成元相依（$\sqrt{7+2\sqrt{10}}$ 即为例）；
\item 定理~\ref{thm:d2} 中 $\omega_n=\sqrt[n]{1}$ 是按定义合法的深度 $1$ 根式，
这是 D2 判据「允许新根式的底是有理数」能够写成 $\alpha=\gamma^nc$ 的关键；
\item 命题~\ref{prop:trap} 与第~\ref{sec:complex} 节的假阳性反例是检验任何实现的
第一组试金石：$\sqrt[8]{16}$、$\sqrt[4]{-4}$、$\sqrt{7+2\sqrt{10}}$、$\sqrt{5+2\sqrt6}$。
\end{enumerate}
\pagebreak

\section*{结论}
本文从正十七边形出发，完整地讨论了分圆方程的根式解：
第一章交代了问题的由来、四件预备工具与三条路线共用的套路；
第二章给出了全阶法与半阶法两条一次开方的路线，
并用五次、十七次、七次三个算例说明各自的计算量；
第三章给出了周期下降法，沿十三次与十七次两个算例逐层走到底，
得到每个单位根的实部与虚部；
第四章则回答了全阶法与半阶法一直留下的那个问题——
它们的公式里那个高次根号能否真正开出来，
判定归结为两个域内在的问句 D1 与 D2，求值则由两层封闭公式与六步手算法给出；\par

回顾全文，笔者以为本文最有意思的地方有两点：
其一，三条路线最终都归结到同一件事——
把开方量写成$\omega_P$的多项式的某个幂次根，
区别只在于把困难集中进一个高次根号，还是摊到一串低次根号上；
其二，周期下降法每一步所解的方程都可以只用数数写出来，
块乘法的计数系数$\textcyrillic{Б}$与由它排成的整数矩阵，
把 Gauß 周期方程直接交到手上，不必先解任何方程；\par

当然，本文的讨论仍有一些可以继续深入的方向：
例如$L_r\ge7$时辅助单位根递归构造的实际代价，
以及周期方程在更大规模算例上的形态，笔者尚未进行系统的研究；
另外，本文只讨论了素数次的单位根，
合数次的情形虽然按算术基本定理可以化归，
其逐步回代的具体写法也还没有展开；
至此，本文的讨论告一段落。\par

\pagebreak

\appendix

\section{附录}
\subsection{分圆部分的记号表}
表中$\pi$与$e$是通用常数记号（圆周率、自然对数底），沿用惯例，不占字体分工；
虚数单位一律写作$\ii$。

\textbf{基本量。}
\begin{center}
\begin{tabular}{c|>{\raggedright\arraybackslash}p{0.68\linewidth}}
\toprule
记号 & 含义 \\
\midrule
$P$ & 目标的奇素数，$P\ge3$ \\
$Q$ & $Q=P-1$，非零剩余类的个数 \\
$H$ & $H=Q/2=(P-1)/2$，实迹的个数 \\
$G$ & 模$P$的一个原根，逐个试$2,3,\dots$，取第一个使$1,G,\dots,G^{Q-1}$模$P$遍历全部非零剩余类者 \\
$G^a$ & 原根表的第$a$行，$0\le a<Q$ \\
$\delta_G(x)$ & $x$的编号，由$G^{\delta_G(x)}\equiv x\pmod P$与$0\le\delta_G(x)<Q$唯一确定；它是原根表的反查 \\
$[x]_m$ & $x$模$m$的最小非负剩余，$0\le[x]_m<m$ \\
$\iota_{\mathscr{P}}$ & 指示函数，整数条件$\mathscr{P}$成立时取$1$、否则取$0$ \\
$\omega_n^t$ & $e^{2\pi t\ii/n}$，本原$n$次单位根的第$t$个幂；阶$n\ge1$，幂次$0\le t<n$ \\
$\ii$ & 虚数单位 \\
\bottomrule
\end{tabular}
\end{center}

\textbf{三条路线共用的预解式与公共根式。}
预解式都记作$Z$，公共根式都记作$U$，只用上标的阶区分；
其中半阶预解式恰好是全阶预解式的偶数位子列。
\begin{center}
\begin{tabular}{c|>{\raggedright\arraybackslash}p{0.68\linewidth}}
\toprule
记号 & 含义 \\
\midrule
$Z^{(Q)}_j$ & 全阶预解式，阶为$Q$ \\
$Z^{(H)}_k$ & 半阶预解式，阶为$H$，且$Z^{(H)}_k=Z^{(Q)}_{2k}$ \\
$\omega_Q$ & 本原$Q$次单位根（全阶法的辅助根） \\
$\omega_H=\omega_Q^2$ & 本原$H$次单位根（半阶法的辅助根） \\
$J_r,\ B_j$ & 全阶法的 Jacobi 和与累乘系数 \\
$K_r,\ C_j$ & 半阶法的 Jacobi 和与累乘系数 \\
$E^{(Q)}_{P,G}$ & 全阶法的开方多项式 \\
$E^{(H)}_{P,G}$ & 半阶法的开方多项式 \\
$U^{(Q)}$ & 全阶法的公共根式，$\bigl(U^{(Q)}\bigr)^{Q}=E^{(Q)}_{P,G}(\omega_Q)$ \\
$U^{(H)}$ & 半阶法的公共根式，$\bigl(U^{(H)}\bigr)^{H}=E^{(H)}_{P,G}(\omega_H)$ \\
\bottomrule
\end{tabular}
\end{center}

\textbf{周期下降法：层结构与块。}
\begin{center}
\begin{tabular}{c|>{\raggedright\arraybackslash}p{0.68\linewidth}}
\toprule
记号 & 含义 \\
\midrule
$g$ & 分解链的长度，即素因子（带重数）的个数 \\
$L_r$ & 分解链的第$r$个素因子，$H=\prod_{a=1}^{g}L_a$，$D_r=\prod_{a=1}^{r}L_a$ \\
$D_r$ & 第$r$层的块数，$D_0=1$，$D_g=H$；$D_{r-1}$是上一层的块数 \\
$N_r$ & $N_r=Q/D_r$，第$r$层每块里的指数个数 \\
$\mathscr{C}_{r,s}$ & 第$r$层的第$s$个块（陪集），$x\in\mathscr{C}_{r,s}\iff\delta_G(x)\equiv s\pmod{D_r}$ \\
$S_{r,s}$ & 第$r$层第$s$块的和，$\sum_{x\in\mathscr{C}_{r,s}}\omega_P^{\,x}$ \\
$\textcyrillic{Б}_{r,s_1,s_2,s_3}$ & 块乘法的计数系数 \\
\bottomrule
\end{tabular}
\end{center}

\textbf{周期下降法：每层内部用到的量。}
\begin{center}
\begin{tabular}{c|>{\raggedright\arraybackslash}p{0.30\linewidth}|>{\raggedright\arraybackslash}p{0.36\linewidth}}
\toprule
记号 & 名称 & 定义式 \\
\midrule
$\omega_{L_r}$ & 本层的辅助单位根 & $\omega_{L_r}=\omega_{L_r}^{1}$ \\
$\chi_{r,j}$ & 权重表 & 在第$b$个细块上取$\omega_{L_r}^{\,jb}$，其余取$0$；$1\le j<L_r$ \\
$\varphi_r$ & 折叠表 & $\varphi_r=\Psi_{\chi_{r,1}}^{\,L_r}(\iota_{x=0})$ \\
$\xi_{r,j}$ & 交叉表 & $\xi_{r,j}=\Psi_{\chi_{r,1}}^{\,L_r-j}(\chi_{r,j})$，$1\le j<L_r$ \\
$\Psi$ & 折叠算子 & $\Psi_\eta(\zeta)(x)=\sum_{f=0}^{P-1}\zeta(f)\eta([x-f]_P)$，$\eta,\zeta$是任意表的占位符 \\
$\Xi_{r-1}$ & 求值泛函 & $\Xi_{r-1}(\zeta)=\zeta(0)+\sum_{s}\zeta(G^{\,s})S_{r-1,s}$ \\
$F_r$ & 本层的开方量 & $F_r=\Xi_{r-1}(\varphi_r)$，取值非零 \\
$T_{r,j}$ & 本层预解式的模 & $T_{r,j}=\Xi_{r-1}(\xi_{r,j})$，$1\le j<L_r$ \\
$Y_{r,j}$ & 系数比 & $Y_{r,j}=T_{r,j}/F_r$，且$Y_{r,1}=1$ \\
$R_r$ & 满足$R_r^{\,L_r}=F_r$的一个取定的根 & 复数，分支任取 \\
$X_{r,s}$ & 第$0$块里的$L_r$个新块和（反演输出） & $0\le s<L_r$ \\
\bottomrule
\end{tabular}
\end{center}
分量与块和之间差一个因子：第$j$个非平凡分量写成$Y_{r,j}R_r^{\,j}$，
也就是系数比乘上同一个根$R_r$的$j$次幂，不能再各自另取分支。\par

\textbf{周期下降法：恢复与收尾用到的量。}
\begin{center}
\begin{tabular}{c|>{\raggedright\arraybackslash}p{0.27\linewidth}|>{\raggedright\arraybackslash}p{0.38\linewidth}}
\toprule
记号 & 名称 & 定义式 \\
\midrule
$\rho_{r,b}$ & 第$b$个细块的指示表 & $\rho_{r,b}(x)=\iota_{x\in\mathscr{C}_{r,D_{r-1}b}}$，$0\le b<L_r$ \\
$\lambda_{r,h}$ & 幂表 & $\lambda_{r,h}=\Psi_{\rho_{r,0}}^{\,h}(\iota_{x=0})$，$0\le h<D_r$ \\
$\Theta_{r,s,h}$ & 恢复矩阵的元素 & $\lambda_{r,h}(G^s)-\lambda_{r,h}(0)$，$0\le s,h<D_r$ \\
$\theta_s$ & 恢复多项式 & 在$X_{r,0}$处取值，$S_{r,s}=\theta_s(X_{r,0})$ \\
$V_n$ & 终层递推多项式 & $V_0=2$，$V_1=y$，$V_{n+1}=yV_n-V_{n-1}$，$n\ge0$ \\
$u_s$ & 终层递推所需的下标 & $u_s=\min\bigl([G^s]_P,\,P-[G^s]_P\bigr)$，$1\le u_s\le H$ \\
$v$ & 最大实迹所在的下标 & 由最大值判定式定出，$0\le v<H$ \\
$W_q$ & 平移后的实迹 & $W_q=S_{H,[v+q]_H}$，$0\le q<H$ \\
$\alpha_t$ & 实部该查哪一行实迹 & $\alpha_t=[\delta_G(t_0)]_H$，$t_0\ne0$ \\
$\beta_t$ & 虚部该查哪一行实迹 & $\beta_t=[\delta_G([2t_0]_P)]_H$，$t_0\ne0$ \\
$\varepsilon_t$ & 虚部符号 & $\operatorname{sgn}\sin\dfrac{2\pi t_0}{P}$，取值$\pm1$ \\
\bottomrule
\end{tabular}
\end{center}
幂表$\lambda_{r,h}$是纯计数：$\lambda_{r,h}(x)$就是把$x$写成
$h$个$\mathscr{C}_{r,0}$中元素之和（可重复、计次序）的方法数。
恢复矩阵的转置作用很简单：对每个$h$成立
\[
 X_{r,0}^{\ h}=\sum_{s=0}^{D_r-1}\Theta_{r,s,h}\,S_{r,s},
\]
于是$D_r$个方程可以解出$D_r$个未知的块和，而$\theta_s$的系数就是解出来的系数。
矩阵$\Theta_r$的第$0$列（$h=0$）恒为$-1$，因为$\lambda_{r,0}=\iota_{x=0}$。
恢复矩阵属算子一类，用大写希腊$\Theta$，与预解式$Z$分属两个字母，不会混淆。
\subsection{全书字母分工表}
本表按字母档位汇总全书用字，供查字母占用与分工；
按对象列示的定义式与取值范围见附录中分圆部分的记号表与正文各节。\par

\subsubsection{字母分工总则}
\begin{center}
\begin{tabular}{c|>{\raggedright\arraybackslash}p{0.70\linewidth}}
\toprule
档位 & 用途 \\
\midrule
小写拉丁 & 函数的变量、求和与求积的中间变量、各种下标 \\
大写拉丁 & 过程中的一般数量：系数、常数、比值、块和、预解式 \\
花体拉丁 & 集合、代数结构、数域、整数条件 \\
镂空拉丁 & 约定俗成的数域与整数环（$\QQ$、$\ZZ$、$\NN$、$\mathbb F$、$\mathbb C$） \\
小写希腊 & 函数与各种表 \\
大写希腊 & 算子、泛函与矩阵 \\
哥特体 & 代数数论中的理想（素理想与分式理想） \\
西里尔 & 上述字母不足时的补充（本书记数系数用 $\textcyrillic{Б}$） \\
\bottomrule
\end{tabular}
\end{center}

\subsubsection{分圆部分}

\begin{center}
\begin{tabular}{c|>{\raggedright\arraybackslash}p{0.72\linewidth}}
\toprule
字母 & 含义（小写拉丁：函数的变量与各种下标） \\
\midrule
$a$ & 连乘哑下标；原根表的行号，$0\le a<Q$ \\
$b$ & 块内序号，$0\le b<L_r$ \\
$c,\ d$ & 系数（含 $c_s$；七次算例中的 $c+d\,\omega_H$） \\
$f$ & 折叠求和的第二行号，$0\le f<P$ \\
$g$ & 分解链长度，即素因子（带重数）的个数 \\
$h$ & 幂次指标，$0\le h<D_r$ \\
$j$ & 分量序号，$1\le j<L_r$ \\
$k$ & 求和指标，$0\le k<Q$ \\
$m$ & Jacobi 和的下标（与半阶法换算一节） \\
$n$ & 单位根的阶，$\omega_n^t$ \\
$q$ & 平移后的实迹下标，$0\le q<H$ \\
$r$ & 层号，$1\le r\le g$ \\
$s$ & 块号，$0\le s<D_r$ \\
$t$ & 幂次，$0\le t<P$ \\
$u$ & 终层递推所需的下标，$u_s=\min\bigl([G^s]_P,P-[G^s]_P\bigr)$ \\
$v$ & 最大实迹所在的下标 \\
$w$ & 开方量多项式的未定元 \\
$x$ & 行号，$0\le x<P$ \\
$y$ & 形式变量（递推多项式与分圆多项式的未定元） \\
$z$ & 块内序号，$0\le z<N_r$ \\
\bottomrule
\end{tabular}
\end{center}

\begin{center}
\begin{tabular}{c|>{\raggedright\arraybackslash}p{0.72\linewidth}}
\toprule
字母 & 含义（大写拉丁：一般数量） \\
\midrule
$P$ & 目标的奇素数 \\
$Q$ & $Q=P-1$，非零剩余类的个数 \\
$H$ & $H=Q/2$，实迹的个数 \\
$G$ & 模 $P$ 的一个原根 \\
$Z^{(Q)}_j,\ Z^{(H)}_k$ & 全阶、半阶预解式，$Z^{(H)}_k=Z^{(Q)}_{2k}$ \\
$J_r,\ B_j$ & 全阶法的 Jacobi 和与累乘系数 \\
$K_r,\ C_j$ & 半阶法的 Jacobi 和与累乘系数 \\
$E^{(Q)}_{P,G},\ E^{(H)}_{P,G}$ & 开方量：$\bigl(Z^{(Q)}_1\bigr)^Q=E^{(Q)}_{P,G}(\omega_Q)$，$\bigl(Z^{(H)}_1\bigr)^H=E^{(H)}_{P,G}(\omega_H)$ \\
$U^{(Q)},\ U^{(H)}$ & 公共根式，分别为 $E^{(Q)}_{P,G}(\omega_Q)$、$E^{(H)}_{P,G}(\omega_H)$ 的 $Q$、$H$ 次根 \\
$L_r$ & 分解链的第 $r$ 个素因子 \\
$D_r$ & 第 $r$ 层的块数，$D_r=\prod_{a=1}^{r}L_a$，$D_0=1$ \\
$N_r$ & $N_r=Q/D_r$，第 $r$ 层每块内的指数个数 \\
$S_{r,s}$ & 第 $r$ 层第 $s$ 块的和 \\
$F_r$ & 本层的开方量 \\
$T_{r,j}$ & 本层预解式的模 \\
$Y_{r,j}$ & 系数比，$Y_{r,j}=T_{r,j}/F_r$ \\
$R_r$ & 满足 $R_r^{\,L_r}=F_r$ 的一个取定的根 \\
$X_{r,s}$ & 第 $0$ 块里的 $L_r$ 个新块和 \\
$V_n$ & 终层递推多项式，$V_0=2$，$V_1=y$，$V_{n+1}=yV_n-V_{n-1}$ \\
$W_q$ & 平移后的实迹，$W_q=S_{H,[v+q]_H}$ \\
\bottomrule
\end{tabular}
\end{center}

\begin{center}
\begin{tabular}{c|>{\raggedright\arraybackslash}p{0.72\linewidth}}
\toprule
字母 & 含义（花体与镂空） \\
\midrule
$\mathscr{C}_{r,s}$ & 第 $r$ 层的第 $s$ 个块（陪集） \\
$\mathscr{G}_m$ & 模 $m$ 的乘法群 \\
$\mathscr{K}_{D_r}$ & 第 $r$ 层对应的数域 \\
$\mathscr{P}$ & 指示函数所依据的整数条件 \\
$\mathscr{S}$ & 集合（$\iota_{x\in\mathscr S}$） \\
$\NN,\ \QQ,\ \ZZ$ & 自然数、有理数域、整数环（镂空） \\
\bottomrule
\end{tabular}
\end{center}

\begin{center}
\begin{tabular}{c|>{\raggedright\arraybackslash}p{0.72\linewidth}}
\toprule
字母 & 含义（小写希腊：函数与各种表） \\
\midrule
$\omega_n^t$ & $e^{2\pi t\ii/n}$，单位根；辅助根写作 $\omega_Q$、$\omega_H$、$\omega_L$ \\
$\delta_G(x)$ & 编号表，由 $G^{\delta_G(x)}\equiv x$ 与 $0\le\delta_G(x)<Q$ 唯一确定 \\
$\chi_{r,j}$ & 权重表 \\
$\varphi_r$ & 折叠表 \\
$\xi_{r,j}$ & 交叉表 \\
$\iota$ & 指示函数 \\
$\lambda_{r,h}$ & 幂表 \\
$\theta_s$ & 恢复多项式 \\
$\rho_{r,b}$ & 第 $b$ 个细块的指示表 \\
$\mu_Q,\ \mu_H,\ \mu_P$ & 分圆多项式 \\
$\alpha_t,\ \beta_t$ & 收尾时实部、虚部该查哪一行实迹的两个指标 \\
$\varepsilon_t$ & 虚部符号，取值 $\pm1$ \\
$\zeta,\ \eta$ & 算子定义式中任意表的占位符 \\
$\sigma,\ \tau$ & 两个域自同构，$\sigma:\omega_P\mapsto\omega_P^{G}$，$\tau=\sigma^{D_{r-1}}$ \\
$\pi$ & 圆周率（通用常数） \\
\bottomrule
\end{tabular}
\end{center}

\begin{center}
\begin{tabular}{c|>{\raggedright\arraybackslash}p{0.72\linewidth}}
\toprule
字母 & 含义（大写希腊：算子、泛函与矩阵；西里尔：补充） \\
\midrule
$\Psi_\eta$ & 折叠算子，$\Psi_\eta(\zeta)(x)=\sum_{x'}\zeta(x')\eta([x-x']_P)$ \\
$\Xi_{r-1}$ & 求值泛函，$\Xi_{r-1}(\zeta)=\zeta(0)+\sum_s\zeta(G^{\,s})S_{r-1,s}$ \\
$\Theta_r$ & 恢复矩阵，其元素记 $\Theta_{r,s,h}=\lambda_{r,h}(G^s)-\lambda_{r,h}(0)$ \\
$\textcyrillic{Б}_{r,s_1,s_2,s_3}$ & 块乘法的计数系数 \\
\bottomrule
\end{tabular}
\end{center}

\subsubsection{根式去嵌套部分}

\begin{center}
\begin{tabular}{c|>{\raggedright\arraybackslash}p{0.72\linewidth}}
\toprule
字母 & 含义（小写拉丁：变量与下标） \\
\midrule
$a$ & 单生成元的底数，$\theta^m=a\in\mathbb{Q}$ \\
$b,\ c$ & 两层封闭公式中的参数（$\sqrt b$ 与配平系数 $c$） \\
$d$ & 扩张次数，$d\mid n$ \\
$e$ & 指数，$e_r\in\{0,\dots,n-1\}$ \\
$f$ & 剩余域次数 $f(\mathfrak{p}/p)$ \\
$j$ & 生成元下标（$\theta_j$） \\
$k$ & 指数与卷积下标：$\mathfrak{q}^k$、$c_k=\sum_{r+t=k}a_ra_t$（不是域） \\
$m$ & 生成元的个数 \\
$n$ & 外层根式的次数 \\
$p$ & 有理素数 \\
$q$ & 范数条件 $|N(\alpha)|=q^n$ 中的有理数 \\
$r$ & 生成元下标与求和下标 \\
$s$ & 两层公式中的二次层量、求和下标 \\
$t$ & 表达深度 \\
$u$ & 实嵌入的坐标 \\
$v$ & 赋值 $v_p$、$v_{\mathfrak{p}}$ \\
$w$ & 分式 $\beta^n=\gamma^n/w^n$ 中的分母 \\
$x,\ y$ & 辅助变量（如 $x=2^{1/5}$） \\
\bottomrule
\end{tabular}
\end{center}

\begin{center}
\begin{tabular}{c|>{\raggedright\arraybackslash}p{0.72\linewidth}}
\toprule
字母 & 含义（大写拉丁与花体拉丁） \\
\midrule
$A,\ B$ & 两层公式中 $\alpha=A+B\theta$ 的两个系数 \\
$N$ & 域范数 $N_{\mathscr{K}/\mathbb{Q}}$ \\
$P$ & 配平法中两两乘积 $P_{jk}$ \\
$S$ & 单项式的多重下标，$\beta=\sum_S b_S\theta^S$ \\
$T$ & 工具箱条目编号 $T1$–$T7$ \\
$\mathscr{D}$ & Kummer 对应中的子群（代数结构，用花体） \\
$\mathscr{F}$ & 底域（被扩张的域） \\
$\mathscr{K}$ & 内层根式生成的数域（与分圆部分一致） \\
$\mathscr{L}$ & 中间域 \\
$\mathscr{O}_{\mathscr{K}}$ & 整数环，$\mathscr{O}_{\mathscr{K}}^*$ 为其单位群 \\
$\mathscr{U}_n$ & $n$ 次单位根群，$\mathscr{U}_n=\langle\omega_n\rangle$ \\
\bottomrule
\end{tabular}
\end{center}

\begin{center}
\begin{tabular}{c|>{\raggedright\arraybackslash}p{0.72\linewidth}}
\toprule
字母 & 含义（希腊字母、哥特体与镂空） \\
\midrule
$\alpha$ & 被开方数（外层根式的底） \\
$\beta$ & 去嵌套的结果，$\beta^n=\alpha$ \\
$\gamma$ & Kummer 定理中的辅助元素 \\
$\delta$ & Kummer 配对中的子群元素 \\
$\varepsilon$ & 单位（$N(\varepsilon)=\pm1$ 的代数整数） \\
$\theta$ & 单生成元，$\theta^m=a\in\mathbb{Q}$ \\
$\kappa$ & 范数中间量 $N_{\mathscr{K}(\beta)/\mathscr{K}}(\beta)$ \\
$\mu$ & 分圆多项式（与分圆部分一致） \\
$\pi$ & 圆周率（通用常数） \\
$\sigma$ & 域嵌入（影子） \\
$\varphi$ & Euler 函数 $\varphi(n)$ \\
$\omega$ & 单位根；$\omega_n$ 专指本原 $n$ 次单位根（与分圆部分一致） \\
$\mathfrak{p},\ \mathfrak{q},\ \mathfrak{Q}$ & 素理想（哥特体，见分工总则） \\
$\mathfrak{b},\ \mathfrak{B}$ & 分式理想与由元素生成的素理想 \\
$\mathbb{C},\ \mathbb{F},\ \mathbb{Q},\ \mathbb{R},\ \mathbb{Z}$ & 复数域、有限域、有理数域、实数域、整数环（镂空） \\
\bottomrule
\end{tabular}
\end{center}

\subsubsection{两部分之间的字母复用}

分圆部分与根式去嵌套部分处理的是两类不同的对象。凡两部分含义相同的字母
（单位根 $\omega$、数域 $\mathscr{K}$、分圆多项式 $\mu$、虚数单位 $\ii$、镂空数域）
已经统一；只有下面一组字母在两个部分中承担不同含义，正文中各自只在本部分内部使用，不跨部分混用。\par

\begin{center}
\begin{tabular}{c|c|>{\raggedright\arraybackslash}p{0.58\linewidth}}
\toprule
字母 & 所属部分 & 含义 \\
\midrule
$\alpha,\ \beta$ & 分圆部分 & 收尾查表指标（恒带下标 $t$） \\
$\alpha,\ \beta$ & 根式去嵌套部分 & 被开方数与去嵌套结果 \\
\bottomrule
\end{tabular}
\end{center}

\subsection{与三套原始记号的对照}
本文为统一叙述改动了若干记号，换算关系如下。
\begin{center}
\begin{tabular}{c|c|c|c}
\toprule
本文 & 全阶法原稿 & 半阶法原稿 & 周期下降法手算稿 \\
\midrule
$P$ & $p$ & $p$ & $P$ \\
$G$ & $g$ & $g$ & $G$ \\
$\omega_Q$ & $\eta$ & --- & --- \\
$Z^{(Q)}_j$ & $u_j$（亦记$u_p^j$） & --- & --- \\
$U^{(Q)}$ & $R$ & --- & --- \\
$E^{(Q)}_{P,G}$ & $F_p$ & --- & --- \\
$\omega_H$ & --- & $\rho$ & --- \\
$Z^{(H)}_k$ & --- & $v_j$ & --- \\
$U^{(H)}$ & --- & $U$ & --- \\
$E^{(H)}_{P,G}$ & --- & $G_{p,g}$ & --- \\
$S_{r,s}$ & --- & --- & $S_{r,s}$ \\
$R_r$ & --- & --- & $R_r$ \\
\bottomrule
\end{tabular}
\end{center}
虚数单位在本文一律写作$\ii$，原稿中写作$i$的地方都按下表读；
复数一律写成「实部加虚部乘$\ii$」的形式，
$a+b\,\ii$中的$a$与$b$在本文的算法里都是实数。
\subsection{术语速查}
\begin{center}
\begin{tabular}{c|>{\raggedright\arraybackslash}p{0.44\linewidth}|>{\raggedright\arraybackslash}p{0.26\linewidth}}
\toprule
术语 & 一句话 & 详见 \\
\midrule
数域 & 对加、减、乘、除封闭的数集 & 预备：数域与嵌入 \\
影子（共轭） & 保持运算、固定有理数的映射 & 预备：嵌入（影子）与范数 \\
范数 & 全部影子的乘积 & 预备：嵌入（影子）与范数 \\
折叠 & $\theta^{qm+r}=(\theta^m)^q\theta^r=a^q\theta^r$ & 预备：数域与嵌入 \\
$v_p$、$v_{\mathfrak{q}}$ & 有理素数、素理想上的赋值 & 预备：赋值、素理想与无平方因子部分 \\
$\operatorname{sqfree}$ & 削去平方因子 & 同上 \\
$e,\ f$ & 分歧指数与剩余域次数 & 同上 \\
单位 & $N(\varepsilon)=\pm1$ 的代数整数 & 预备：单位与单位根 \\
D1、D2 & 内去嵌套、深度约化 & 问题与两种「可去」的定义 \\
单项式引理 & 答案必是底域元乘生成元的单项式 & 两个主定理：Landau 单项式引理 \\
$\omega$-陷阱 & 在 $\mathscr{K}(\omega_n)$ 中是幂，不等于在 $\mathscr{K}$ 中是幂 & 两个主定理：在 $\mathscr{K}(\omega_n)$ 中检验的危险与挽救 \\
\bottomrule
\end{tabular}
\end{center}
\subsection{参考文献}
{\sloppy
\begin{enumerate}

 \item J. S. Milne. \emph{Fields and Galois Theory}, version 5.10, 2022.
       第 3 章与第 5 章，特别是循环扩张的 Proposition 5.27.\\
       \texttt{https://www.jmilne.org/math/CourseNotes/FT.pdf}
 \item Paul Garrett. \emph{Kummer, Eisenstein, computing Gauss sums as Lagrange resolvents},
       July 28, 2010. 第 1 节的循环预解式与后文 Gauss--Jacobi 乘法关系.\\
       \texttt{https://www-users.cse.umn.edu/\~{}garrett/m/v/kummer\_eis.pdf}

\item S.~Landau, \emph{Simplification of nested radicals}, SIAM J.~Comput.\ 21(1)
(1992) 85--110. [doi:\allowbreak10.1137/\allowbreak0221009]
\item S.~Landau, \emph{How to tangle with a nested radical},
The Mathematical Intelligencer 16(2) (1994). [doi:\allowbreak10.1007/\allowbreak BF03024284]
\item S.~Landau, G.~L.~Miller, \emph{Solvability by radicals is in polynomial time},
J.~Comput.\ System Sci.\ 30 (1985) 182--207.
\item A.~Borodin, R.~Fagin, J.~E.~Hopcroft, M.~T.~Tompa,
\emph{Decreasing the nesting depth of expressions involving square roots},
J.~Symbolic Comput.\ 1 (1985) 169--188.
\item R.~Zippel, \emph{Simplification of expressions involving radicals},
Proc.\ SYMSAC 1985, ACM.
\item H.~Cohen, \emph{A Course in Computational Algebraic Number Theory},
GTM 138, Springer, 1993（类群、单位群、主理想算法）。
\item L.~C.~Washington, \emph{Introduction to Cyclotomic Fields}, 2nd ed., GTM 83,
Springer, 1997（Kummer 理论）。
\end{enumerate}
}

\end{document}